% DER_XI-XIIa: Grothendieck's text with the commentary set in, then the overview, inside a group that applies the chunk's own definitions
\dgrchapter{XI}{Chapter XI}{Hot-fibrations, foncteurs propres et foncteurs lisses etc. (dans Cat)}
\dgrpart{0}{XI-XIIa}{Manuscript pages 127--136}{XI-XIIa\enspace\textperiodcentered\enspace Chapter XI\enspace\textperiodcentered\enspace manuscript pages 127--136}{Chapter XI \textperiodcentered\ \textup{XI-XIIa}}
\dgrchunkbegin
% DER_XI-XIIa: the definitions of the chunk's preamble, applied inside its group
\let\printedpage\relax\newcommand{\printedpage}[1]{\marginpar{\footnotesize p.~#1}}
\let\msp\relax\newcommand{\msp}[1]{\marginpar{\footnotesize [#1]}}
\let\Hot\relax\newcommand{\Hot}{\mathrm{Hot}}
\let\Cat\relax\newcommand{\Cat}{\mathrm{Cat}}
\let\Ob\relax\newcommand{\Ob}{\operatorname{Ob}}
\let\Ch\relax\newcommand{\Ch}{\underline{\mathrm{Ch}}}
\let\Dl\relax\newcommand{\Dl}{\boldsymbol{\Delta}}
\let\fl\relax\newcommand{\fl}{\longrightarrow}
\let\hot\relax\newcommand{\hot}{\mathrm{hot}}
\let\uHom\relax\newcommand{\uHom}{\underline{\mathrm{Hom}}}
\let\id\relax\newcommand{\id}{\mathrm{id}}
\let\pr\relax\newcommand{\pr}{\mathrm{pr}}
\let\Fl\relax\newcommand{\Fl}{\mathrm{Fl}}
\let\card\relax\newcommand{\card}{\operatorname{card}}
\let\uFib\relax\newcommand{\uFib}{\underline{\mathrm{Fib}}}
\let\HOT\relax\newcommand{\HOT}{\mathrm{HOT}}
\let\HOTlc\relax\newcommand{\HOTlc}{\mathrm{HOT}^{\mathrm{lc}}_W}
\let\Hom\relax\newcommand{\Hom}{\mathrm{Hom}}
\let\Wd\relax\newcommand{\Wd}{W^{\mathrm{d}}_S}
\let\Wf\relax\newcommand{\Wf}{W^{\mathrm{f}}_S}
\let\dashedbox\relax\newcommand{\dashedbox}[1]{\tikz[baseline=(n.base)]\node[draw,dashed,inner sep=2pt](n){#1};}
\let\Fib\relax\newcommand{\Fib}{\mathrm{Fib}}
\let\Parf\relax\newcommand{\Parf}{\mathrm{Parf}}

% ---- from Der-ChXI-890_p001-005.tex
% ===================== page 1 =====================
\printedpage{--}
% head: (none) -- title page, no page number printed

\begin{center}
\vspace*{5cm}
{\LARGE\textsf{LES DÉRIVATEURS}}

\vspace{1.5em}
{\large\textbf{ALEXANDRE GROTHENDIECK}}

\vspace{3em}
{\large\textbf{Chapitre XI}}

\vspace{0.5em}
{\large $\Hot$\textbf{-fibrations, foncteurs propres et foncteurs lisses etc.}\\
\textbf{(dans} $\Cat$\textbf{)}}
\end{center}

\clearpage

% ===================== page 2 =====================
\printedpage{--}
% head: (none) -- no page number printed

\vspace*{6cm}
\noindent Ce texte a été déchiffré et transcrit en \LaTeXe\ par M. Künzer. Il a été édité par M. Künzer, J. Malgoire, et G. Maltsiniotis. La transcription est aussi fidèle que possible au manuscrit. Pour les quelques corrections évidentes, ou rares commentaires des éditeurs, ainsi que pour la numérotation originale des pages du manuscrit, les caractères de machine à écrire \texttt{[typewriter]} entre crochets sont utilisés. Un point d'interrogation entre crochets signifie que l'on n'est pas sûr du mot qui précède.

\noindent Cette édition est provisoire. Les remarques, commentaires et corrections sont bienvenus. Envoyer un message à :

\noindent\textbf{[e-mail address removed]}

\hfill G. Maltsiniotis\hspace*{2.5em}

\clearpage

% ===================== page 3 =====================
\printedpage{1}
% head: A. Grothendieck : Dérivateurs, Ch. XI, Édition des Universités de Montpellier II et Paris VII
% foot: Version 28 mai 2002

\noindent\texttt{[Il n'y a pas de pages 1-126]}

\msp{page 127}

\section*{1\quad $\Hot$-fibrations dans $\Cat$}

\noindent\textbf{Théorème 1.} \textit{Soit $f : X \longrightarrow Y$ dans} $\Cat$\textit{. Conditions équivalentes :}

\dgrblock{dgrb450}{The universal meaning of a Hot-fibration}
\begin{dcomment}{exposition}{p127:Conditions équivalentes :}{The universal meaning of a Hot-fibration}
Grothendieck wants the ordinary pullback of categories to calculate a homotopy pullback whenever one of its sides belongs to a recognisable class. This is what makes a fibre a homotopically meaningful object, rather than an accident of a presentation. Write \(BC=|NC|\), where \(N\) is the nerve, and let \(W_\infty\) consist of the functors inducing weak homotopy equivalences of classifying spaces. Condition a) says that
\[
 f^*:\Cat/Y\longrightarrow\Cat/X,\qquad
 (T\longrightarrow Y)\longmapsto(X\times_Y T\longrightarrow X),
\]
preserves \(W_\infty\). The quantifier ranges over maps between \emph{arbitrary} objects of \(\Cat/Y\), and not only over weak equivalences with target \(Y\). It therefore includes every further change of base.

In the terminology of \citet[\S2, Proposition 2.7]{Rezk1998Sharp}, these are precisely the sharp maps in the Thomason model category. The identification uses the definition of sharpness and the right properness of that model category; it does not identify this class with the smaller class of its model fibrations. Every strict pullback square along such a map is homotopy cartesian. This interpretation also explains why composition and arbitrary base change will be available without choosing liftings.

There is a small but decisive quantifier in a') and a''). The point \(e\longrightarrow Y'\) is arbitrary; it need not select the final, respectively initial, object whose existence is assumed. Over the interval \(0\longrightarrow1\), both endpoint inclusions must therefore become weak equivalences after pullback. Asking only for the distinguished final endpoint, for every category having a final object, leads instead to the later notion of preproperness. Asking only over intervals is weaker again. These distinctions become the subject of Chapter XII, and they should not be erased by replacing every occurrence of “fibration” with a single modern definition.
\end{dcomment}


\begin{itemize}[leftmargin=3em]
\item[\textit{a)}] \textit{Pour tout $Y'' \xrightarrow{p} Y'$ sur $Y$, d'où $X'' = X \times_Y Y'' \xrightarrow{q} X' = X \times_Y Y'$, si $p$ est une} $\Hot$\textit{-équivalence, $q$ l'est aussi.}

\item[\textit{a$'$)}] \textit{Comme a), avec $Y'' = e$, $Y'$ ayant un objet final.}

\item[\textit{a$''$)}] \textit{Comme a), avec $Y'' = e$, $Y'$ ayant un objet initial} (\footnote{garder a$''$) et b$''$), conditions duales de a$'$), b$'$) \texttt{[a$''$) et b$''$) étaient barrés]}.})\textit{.}

\item[\textit{b$'$)}] \textit{Pour tout $y \in \Ob Y$, $X_y \longrightarrow X/y$ est une} $\Hot$\textit{-équivalence, pour tout $y \longrightarrow y'$ dans $Y$, $X/y \longrightarrow X/y'$ est une} $\Hot$\textit{-équivalence; enfin ces conditions sont encore vérifiées pour $f' : X' \longrightarrow Y'$, pour tout changement de base $Y' \longrightarrow Y$ dans} $\Cat$\textit{. (En fait, il suffit de le demander pour la première des deux conditions, concernant $X_y \longrightarrow X/y$.)}

\item[\textit{b$''$)}] \textit{Dual de b$'$), avec $X_y \longrightarrow y\backslash X$ et $y'\backslash X \longrightarrow y\backslash X$.}

\item[\textit{c)}] \textit{Pour tout faisceau d'ensembles localement constant $F$ sur $X$, $f_*(F)$ est localement constant et se calcule `fibre par fibre', de même pour $R^1 f_*(F)$ si $F$ est un faisceau en groupes, et pour les $R^i f_*(F)$ si $F$ est un faisceau abélien. De plus, on veut que ces conditions restent satisfaites après tout changement de base.} (\footnote{\textit{c$'$) (version homologique). 1$^o$) en termes d'images directes de groupoïdes localement constants, et 2$^o$) en termes des $R^i f_!(L)$} \texttt{[plutôt $L_i f_!(L)$]}\textit{, pour $L$ faisceau abélien localement constant.}})
\end{itemize}

\noindent\textsc{Démonstration}. On a les implications $\pm$ tautologiques
\[
\mathrm{a} \implies \mathrm{a}' \implies \mathrm{b}' \implies \mathrm{c},
\]
et il reste à prouver que c $\implies$ a. Comme a est autoduale, étant équivalente à a$'$, b$'$, elle est aussi équivalente aux conditions duales de celles-ci, a$''$ et b$''$ (cf. cor. 2, page 75). On est ramené à prouver que si c) est satisfait, et si on a
\[
\begin{tikzcd}[row sep=large, column sep=large]
X \arrow[d, "f"'] & X' \arrow[l, "q"'] \arrow[d, "f'"] \\
Y & Y' \arrow[l, "p"]
\end{tikzcd}
\]
\printedpage{2}
% head: A. Grothendieck : Dérivateurs, Ch. XI, Édition des Universités de Montpellier II et Paris VII
% foot: Version 28 mai 2002
cartésien dans $\Cat$, si $p$ est une $\Hot$-équivalence, $q$ aussi. On utilise le critère d'Artin-Mazur.

\dgrblock{dgrb451}{Comma categories connect the geometric and coefficient tests}
\begin{dcomment}{exposition}{p127:On utilise le critère d'Artin-Mazur.}{Comma categories connect the geometric and coefficient tests}
The comma categories in b') serve two purposes at once: they express a pointwise right Kan extension, and they supply canonical transition maps. We use \(X/y=(f\downarrow y)\), with objects \((x,f(x)\longrightarrow y)\). Its strict fibre \(X_y\) is the full subcategory on identity arrows. A morphism \(u:y\longrightarrow z\) gives a genuine functor \(u_*:X/y\longrightarrow X/z\), so the coherence needed for transport is already present before localization.

Here is the short argument hidden in a') \(\Rightarrow\) b'). Apply a') over \(Y/y\), which has final object \((y,1_y)\). The inclusion of the strict fibre over that object is \(X_y\longrightarrow X/y\). Next work over \(Y/z\) and select its object \((y,u)\). The resulting weak equivalence \(X_y\longrightarrow X/z\) is the composite of \(X_y\longrightarrow X/y\) and \(u_*\). Two-out-of-three proves that \(u_*\in W_\infty\). The same argument works after every base change.

The parenthesis in b') retains the transition condition on the original base; it says that only the fibre-to-comma condition needs to be imposed universally. Indeed, after \(v:Y'\longrightarrow Y\), the comparison
\[
 X'/y'\longrightarrow X/v(y')
\]
is a weak equivalence: both its source and target receive a weak equivalence from the same strict fibre \(X_{v(y')}\). Comparing this square for an arrow of \(Y'\) with the corresponding arrow of \(Y\) then proves the required transition condition on \(Y'\).

For a locally constant coefficient \(F\), the value of the derived direct image at \(y\) is computed on \(X/y\). Thus the preceding weak equivalences identify it with the cohomology of \(X_y\), and identify its transition maps with isomorphisms. In degrees zero and one this means sections and torsors; for abelian coefficients it means every derived degree. The use of local coefficients is essential: untwisted integral homology alone cannot recognise the fundamental group or its action on higher information. The homological form of the categorical Whitehead criterion is stated explicitly in \citet[Proposition 2.4]{Thomason1980}; the cohomological formulation used here is the criterion to which Grothendieck gives the name Artin--Mazur.
\end{dcomment}


% ===================== page 4 (printed 2) =====================

\noindent 1$^o$) Si $F$ est un faisceau d'ensembles localement constant sur $X$, il faut prouver que $H^0(X,F) \xrightarrow{\sim} H^0(X', q^*F)$.

\msp{page 128}

\noindent Mais
\[
H^0(X,F) \simeq H^0(Y,\ \underbrace{f_*(F)}_{\substack{\text{localement constant}\\ \text{par hypothèse c)}}}\ )
\qquad
\underbrace{\xrightarrow{\qquad\sim\qquad}}_{\substack{\text{car } p\\ \text{Hot-équivalence}}}
\qquad
H^0(Y', p^* f_*(F)),
\]
mais
\[
p^* f_*(F) \xrightarrow{\ \sim\ } f'_*(q^*(F)),
\]
car $f_*(F)$ et $f'_*(q^*F)$ se calculent `fibre par fibre' par hypothèse. Donc
\[
H^0(X,F) \simeq H^0(Y', f'_*(q^*(F))) \simeq H^0(X', q^*(F)),
\]

\hfill q.e.d.

\noindent 2$^o$) Si $F$ est un faiscau en groupes localement constant sur $X$, il faut prouver % [sic] faiscau -> faisceau, misspelling
\[
H^1(X,F) \xrightarrow{\ \sim\ } H^1(X', q^*F).
\]
Le fait que c'est injectif résulte déjà de 1$^o$), qui implique que le foncteur
\[
\begin{array}{rcl}
F\text{-torseurs sur } X & \longrightarrow & q^*(F)\text{-torseurs sur } X' \\
P & \longmapsto & f^*P % [sic] f^*P -> q^*P, the functor is the inverse image by q from X to X'
\end{array}
\]
est pleinement fidèle, par l'argument bien connu utilisant les faisceaux $\mathcal{H}$om. Il faut voir que ce foncteur est essentiellement surjectif.

\msp{page 129}

\noindent On a un homomorphisme de suites exactes d'ensembles ponctués
\[
\begin{tikzcd}[column sep=large, row sep=large,
  every matrix/.append style={name=m},
  execute at end picture={
    \coordinate (lb) at ([xshift=0.4em]m-2-3.east);
    \coordinate (lt) at ([xshift=1.8em]m-2-3.east);
    \draw (lt |- m-1-3.north) -- (lb |- m-1-3.north) -- (lb |- m-2-3.south) -- (lt |- m-2-3.south);
    \coordinate (rb) at ([xshift=0.6em]m-2-4.east);
    \coordinate (rt) at ([xshift=-0.6em]m-2-4.east);
    \draw (rt |- m-1-3.north) -- (rb |- m-1-3.north) -- (rb |- m-2-3.south) -- (rt |- m-2-3.south);
  }]
H^1(Y, f_*F) \arrow[r, "u"] \arrow[d, "\alpha"] & H^1(X,F) \arrow[r, "v"] \arrow[d, "\beta"] & \overbrace{H^0(Y, R^1 f_*(F))}^{\ni\, \eta} \arrow[r, "w"] \arrow[d, "\gamma"] & H^2(Y, f_*F) \arrow[d, "\delta"] \\
H^1(Y', f'_*F) \arrow[r, "u'"] & \underbrace{H^1(X', F')}_{\ni\, \xi'} \arrow[r, "v'"] & \underbrace{H^0(Y', R^1 f'_*(F'))}_{\ni\, \eta'} \arrow[r, "w'"] & H^2(Y', f'_*F') % [sic] H^1(Y', f'_*F) -> H^1(Y', f'_*F'), the other entries of the lower row are in F'
\end{tikzcd}
\]
(\footnote{où on a posé $F' = q^*(F)$.}), où les faisceaux $f_*F$, $R^1 f_*(F)$ sont localement constants par hypothèse et se calculent fibre par fibre, les faisceaux $f'_*(F')$ et $R^1 f'_*(F')$ aussi, qui implique que ces derniers sont
\printedpage{3}
% head: A. Grothendieck : Dérivateurs, Ch. XI, Édition des Universités de Montpellier II et Paris VII
% foot: Version 28 mai 2002
canoniquement isomorphes (par les flèches de changement de base) aux images inverses de leurs homologues $f_*(F)$, $R^1 f_*(F)$. Comme $p$ est une $\Hot$-équivalence, cela implique que les flèches $\alpha$, $\gamma$ sont \texttt{[des]} isomorphismes. Donc si $\xi' \in H^1(X', F')$, son image $\eta'$ dans $H^0(Y', R^1 f'_*(F'))$ provient d'un $\eta \in H^0(Y, R^1 f_*(F))$ bien déterminé. Il faut voir que cet $\eta$ provient d'un $\xi \in H^1(X,F)$. Ceci (si je me rappelle bien) est possible si et seulement si une certaine obstruction $w(\eta)$ dans $H^2(Y, f_*(F))$ est dans le sous-ensemble neutre du $H^2$. \textbf{NB} Le champ sur $Y^\wedge$ des $F$-torseurs sur les $X^\wedge/U$ qui remontent $\eta$ est une $\mathcal{F} = f_*F$-gerbe $\mathcal{G}_\eta$, et il faut voir qu'elle est neutre. Mais il y a mieux à faire :

\dgrblock{dgrb452}{Torsor descent and the band of the obstruction}
\begin{dcomment}{exposition}{p129:Mais il y a mieux à faire}{Torsor descent and the band of the obstruction}
The difficult part of the coefficient argument is essential surjectivity for torsors. Full faithfulness is already secure: if \(P,Q\) are \(F\)-torsors, the sheaf \(\mathcal I=\underline{\operatorname{Isom}}_F(P,Q)\) is locally constant, and
\[
 \operatorname{Hom}_F(P,Q)=H^0(X,\mathcal I).
\]
An equivariant map between two torsors is invertible, and the degree-zero calculation identifies these Hom sets after pullback.

The displayed obstruction sequence should, however, be read as a provisional sketch. For a general nonabelian \(F\), the gerbe of lifts of a section \(\eta\) of \(R^1f_*F\) need not have the fixed band of \(f_*F\). Its band is obtained from automorphisms of local representatives of \(\eta\), hence involves inner forms of \(F\). For a direct check, take \(f:BC_2\to e\), constant \(F=S_3\), and the torsor represented by the inclusion of a transposition subgroup \(C_2\hookrightarrow S_3\). Its automorphism group is \(C_{S_3}(C_2)=C_2\), whereas \(f_*F=S_3\). The component of this torsor therefore cannot have the fixed band \(S_3\). Nor does nonabelian \(H^2\) behave as a group with a distinguished zero: neutrality is a condition on a gerbe. The notation \(G\) in the next paragraph is not defined. These are reasons to replace that portion of the argument, not reasons to reject the theorem. The classical framework is the theory of bands and gerbes in \citet[Chapter IV]{Giraud1971}; we give a direct repair that avoids choosing an obstruction class.

For each \(y\), let \(\mathcal T(y)\) be the groupoid of torsors under the restriction of \(F\) to \(X/y\). Restriction along \(X/y\longrightarrow X/z\) makes \(\mathcal T\) a groupoid-valued contravariant diagram on \(Y\). We claim that every transition functor is an equivalence. Essential surjectivity follows from the local constancy of \(R^1f_*F\): its value at \(y\) is precisely the set of isomorphism classes of \(\mathcal T(y)\). To prove full faithfulness for a transition \(u:y\longrightarrow z\), take \(P,Q\in\mathcal T(z)\). Their isomorphism sheaf is a locally constant set-valued coefficient on \(X/z\). Apply the universally quantified degree-zero hypothesis to the base change
\[
 f_z:X/z\longrightarrow Y/z.
\]
The values at \((z,1_z)\) and \((y,u)\) are respectively
\(\operatorname{Hom}(P,Q)\) and
\(\operatorname{Hom}(u^*P,u^*Q)\).
Local constancy makes the restriction map bijective. This is the place where coefficients on a base-changed total category, rather than only coefficients originally chosen on \(X\), are needed.

The same hypotheses identify \(\mathcal T(y)\) with the groupoid of \(F|_{X_y}\)-torsors. On isomorphism classes this is the fibre formula for \(R^1(f_y)_*F\), applied over \(Y/y\). On Hom sets it is the fibre formula for \((f_y)_*\mathcal I\), for the isomorphism sheaf just used. Thus the identification is an equivalence of groupoids, not merely a bijection on \(H^1\).

Now let \(p:Y'\longrightarrow Y\) be a weak equivalence, and form the cartesian square of the source. The natural comparison from \(p^*\mathcal T\) to the torsor diagram \(\mathcal T'\) on \(Y'\) is an objectwise equivalence: both values at \(y'\) restrict equivalently to the torsors on the same strict fibre \(X_{p(y')}\). Finally,
\[
 \operatorname{Tors}_F(X)
 \simeq \operatorname*{holim}_{y\in Y^{\mathrm{op}}}\mathcal T(y).
\]
This is ordinary descent for torsors, expressed on the representable objects \(Y/y\). Explicitly, a compatible family of torsors on these slices determines its torsor over \(x\in X\) by evaluation at \((x,1_{f(x)})\); the comparison maps along \(f(a)\), for \(a:x\longrightarrow x'\), give its restriction maps. The descent cocycle is exactly their composition law. Pulling a torsor on \(X\) to all slices gives the inverse construction.

Homotopy limits of locally constant groupoid diagrams are invariant under a weak equivalence of indexing categories. One can express this as invariance of sections under equivalence of fundamental \(2\)-groupoids, or obtain it by applying the homotopy-invariance theorem for locally constant model-category diagrams to groupoids and their derived right adjoint of constant diagrams; see \citet[Theorem 1.3]{Cisinski2009LocallyConstant}. Consequently
\[
 \operatorname{Tors}_F(X)\longrightarrow
 \operatorname{Tors}_{q^*F}(X')
\]
is an equivalence. This proves the missing essential surjectivity and the required bijection on \(H^1\). It also supplies the local-constancy argument for the lifting gerbe that the source compresses into a parenthesis.
\end{dcomment}


% ===================== page 5 (printed 3) =====================

\msp{page 130}

\noindent On note que cette gerbe $\mathcal{G}$ est localement constante (puisqu'elle correspond à un lien $f_*(G)$ localement constant) et que $\mathcal{G}'$ est son image inverse. Mais comme $p$ est une $\Hot$-équivalence, il s'ensuit que le foncteur
\[
\underline{\Gamma}(Y, \mathcal{G}) \longrightarrow \underline{\Gamma}(Y', \mathcal{G}')
\]
est une équivalence de catégories, donc $\xi'$ provient (à isomorphisme près) d'un unique $\xi$ de $\underline{\Gamma}(Y, \mathcal{G})$, q.e.d.

\noindent 3$^o$) Soit $F$ faisceau abélien localement constant sur $X$, il faut voir \texttt{[que]}
\[
H^i(X,F) \xrightarrow{\ \sim\ } H^i(X', F') \qquad (F' = q^*F), \tag*{$(*)$}
\]
pour tout $i \in \mathbf{N}$. Mais les deux termes sont l'aboutissement des suites spectrales de Leray, de termes initiaux
\[
\begin{aligned}
E_2^{pq} &\simeq H^p(Y, R^q f_*(F)) \\
E^{\prime pq}_2 &\simeq H^p(Y', R^q f'_*(F')),
\end{aligned}
\]
et $(*)$ est associé à un homomorphisme de suites spectrales, qui pour les termes initiaux est
\[
H^p(Y, R^q f_*(F)) \xrightarrow{\sim} H^p(Y', R^q f'_*(F)), \tag*{$(**)$} % [sic] R^q f'_*(F) -> R^q f'_*(F'), cf. the E'_2 term just above
\]
associé à
\[
p^*(R^q f_*(F)) \longrightarrow R^q f'_*(F'). \tag*{$(*{*}*)$}
\]
Mais par hypothèse les $R^q f_*(F)$, $R^q f'_*(F')$ sont localement constants et se calculent fibre par fibre,

\msp{page 131}

\noindent d'où il résulte que les homomorphismes de changement de base $(*{*}*)$ sont \texttt{[des]} isomorphismes, et comme il s'agit de faisceaux localement constants et que $p$ est une $\Hot$-équivalence, il s'ensuit que $(**)$ est \texttt{[un]} isomorphisme, donc aussi $(*)$ sur les aboutissements, q.e.d.

\dgrblock{dgrb453}{The Leray comparison has finite diagonals}
\begin{dcomment}{exposition}{p131:donc aussi $(*)$ sur les aboutissements, q.e.d.}{The Leray comparison has finite diagonals}
Grothendieck can now return to abelian coefficients, where the formal argument is both shorter and complete. For each \(q\), the canonical comparison
\[
 p^*R^qf_*F\longrightarrow R^qf'_*q^*F
\]
is an isomorphism because its values are the cohomology groups of identical strict fibres. Both sides are locally constant. Since \(p\) is a weak equivalence, it therefore induces an isomorphism on every \(E_2^{p,q}\) term of the Leray spectral sequences.

There is no convergence obstruction here. These are first-quadrant cohomological spectral sequences, and in total degree \(n\) only the \(n+1\) pairs \(p+q=n\) occur. Comparison gives isomorphic associated graded groups of finite filtrations of \(H^n(X,F)\) and \(H^n(X',q^*F)\); induction through those filtrations gives the asserted isomorphism. No boundedness of the cohomological dimension of \(X\) or \(Y\) is required. The missing prime on \(F\) in the printed \(E_2\)-comparison is an edition slip, retained in the transcription and recorded separately. Together with the torsor argument, this completes the coefficient criterion.
\end{dcomment}


\noindent\textbf{Définition 1.} Sous les conditions équivalentes du théorème 1, on dit que $f$ est une $\Hot$-\textit{fibration.}

% ---- from Der-ChXI-890_p006-008.tex
\printedpage{4}
% head: A. Grothendieck : Dérivateurs, Ch. XI, Édition des Universités de Montpellier II et Paris VII
% foot: Version 28 mai 2002

\textbf{Proposition 1.} \textit{Cette notion est stable par changement de base (trivial), par composition (formel sur a)). Pour que $f$ soit une \textup{Hot}-fibration, il faut et il suffit que $f$ \textup{\texttt{[plutôt}} $f^{o}$\textup{\texttt{]}} le soit.}

\dgrblock{dgrb454}{From strict fibres to homotopy fibres}
\begin{dcomment}{later}{p131:Cette notion est stable par}{From strict fibres to homotopy fibres}
The stability assertions now have a geometric interpretation. For composable maps \(X\longrightarrow Y\longrightarrow Z\), pulling a weak equivalence back along the composite is the succession of the two pullbacks. Opposite categories preserve strict pullbacks, and reversing a nerve simplex identifies the classifying spaces of \(C\) and \(C^{\mathrm{op}}\). Thus both composition and self-duality follow directly from a), without another coefficient argument.

The strict fibre also calculates the homotopy fibre. By b'), the categories \(X/y\) vary by weak equivalences along all arrows of \(Y\), and \(X_y\longrightarrow X/y\) is a weak equivalence. The dual form of Quillen's Theorem B therefore identifies
\[
 BX_y \simeq \operatorname{hofib}_y(BX\longrightarrow BY).
\]
The exact citation is \citet[\S1, Theorem B and its corollary]{Quillen1973}; its hypotheses concern the comma-category transition functors. Merely knowing the homotopy types of the strict fibres, without the comparison to comma categories, would not suffice.

This perspective explains the force of the theorem while leaving room for Grothendieck's chosen constructions. \citet[Proposition 4.4]{Cisinski2009LocallyConstant} later proves that weak equivalences of nerves together with the maps that are both smooth and proper form a category of fibrant objects, with functorial factorization. These smooth-and-proper maps are a particularly useful subclass of the present Hot-fibrations. The statement does not say that all Hot-fibrations are smooth and proper. Their advantages are the explicit categories of lifts and the simultaneous control of left and right homotopy Kan extensions.
\end{dcomment}


\textbf{Proposition 2.} \textit{Soit $f : X \longrightarrow S$ une \textup{Hot}-fibration. Alors $s \longmapsto \mathrm{hot}(X_s)$ est un système local sur $S$, à valeurs dans \textup{Hot}. Donc les $X_s$ pour $s$ dans une même composante connexe de $S$ sont \textup{Hot}-isomorphes entre eux (mais pas canoniquement, si cette composante n'est pas \textup{1}-connexe).}

\dgrblock{dgrb455}{A group extension and its monodromy}
\begin{dcomment}{examples}{p131:système local sur $S$, à valeurs dans}{A group extension and its monodromy}
The local system of fibre homotopy types already contains nontrivial information in a finite group example. Regard a group as its one-object category. If
\[
 1\longrightarrow K\longrightarrow G
 \xrightarrow{\varphi}Q\longrightarrow1
\]
is exact, the functor from the one-object groupoid of \(G\) to that of \(Q\) is an isofibration. Its nerve is a Kan fibration. Pullbacks of weak equivalences along Kan fibrations are weak equivalences, and the nerve preserves pullbacks, so this functor is a Hot-fibration. Its strict fibre is the one-object groupoid \(K\).

The comma category over the unique object of \(Q\) has objects \(q\in Q\). A morphism from \(q\) to \(q'\) is an element \(g\in G\) satisfying \(q=q'\varphi(g)\), with a harmless reversal if the other convention for composition is chosen. Surjectivity of \(\varphi\) makes this a connected groupoid; the automorphism group of any object is \(K\). Thus the inclusion of the strict fibre is an equivalence of categories. If \(\varphi\) is not surjective, its connected components are the corresponding cosets of \(\varphi(G)\), and that inclusion is not a weak equivalence. This computes precisely why an arbitrary homomorphism of groups need not be a Hot-fibration.

For the nontrivial monodromy, take \(G=\mathfrak S_3\), \(K=\mathfrak A_3=C_3\), and \(Q=C_2\), with \(\varphi\) the sign. A lift of the nonidentity element of \(C_2\) is a transposition. Conjugation by it sends a three-cycle to its inverse. A different lift changes this automorphism by an inner automorphism of \(K\); inner automorphisms induce naturally isomorphic functors on the one-object groupoid and hence the same map in \(\Hot\). The action on the fibre is consequently well defined there. It is nontrivial, since inversion acts nontrivially on \(\pi_1(BC_3)=C_3\).

With constant coefficient \(\mathbb F_3\) on \(B\mathfrak S_3\), the degree-one fibre cohomology is
\[
 H^1(C_3;\mathbb F_3)
 =\operatorname{Hom}(C_3,(\mathbb F_3,+))
 \cong\mathbb F_3,
\]
and the base generator acts by \(-1\). Its invariant subgroup is zero. The low-degree Leray sequence gives \(H^1(\mathfrak S_3;\mathbb F_3)=0\): the other term \(H^1(C_2;\mathbb F_3)\) is zero, and positive-degree cohomology of \(C_2\) with these coefficients vanishes by the averaging operator \(\frac12(1+\tau)\). Independently, the same answer follows from \(\mathfrak S_3^{\mathrm{ab}}=C_2\). The fibre cohomology is locally constant on the base, but it is not a constant coefficient with trivial monodromy. This is the distinction that the printed qualifier about noncanonical identifications is meant to preserve; a nonsimply connected base can also carry trivial monodromy, as a product shows.
\end{dcomment}


% (a vertical change bar is printed in the left margin beside Théorème 2)
\textbf{Théorème 2.} {\itshape Soit
\begin{tikzcd}[ampersand replacement=\&, baseline=(\tikzcdmatrixname-1-1.base), column sep=small, row sep=small]
X' \arrow[rr, "f"] \arrow[dr, "p'"'] \& \& X \arrow[dl, "p"] \\
\& S \&
\end{tikzcd}
un diagramme commutatif dans \textup{Cat}, avec $p$, $p'$ des \textup{Hot}-fibrations, et soit $S_0 \subseteq S$ une partie de $S$ telle que $S_0 \longrightarrow \pi_0(S)$ soit surjectif. Pour que $f$ soit une \textup{Hot}-équivalence, il faut et il suffit que les $f_s : X_s \longrightarrow Y_s$ $(s \in S_0)$ le soient.} % [sic] f_s : X_s -> Y_s -> f_s : X'_s -> X_s, the diagram has X' and X (cf. f_s : X'_s -> X_s p. 6)

a) La suffisance se prouve comme tantôt, avec le critère de Artin-Mazur.

\msp{page 132}

Notons que l'hypothèse sur $S_0$ assure que si $F$, $F'$ sont deux faisceaux d'ensembles localement constants sur $S$, et $u : F \longrightarrow F'$ un homomorphisme tel que $u_s : F_s \xrightarrow{\ \sim\ } F'_s$ pour $s \in S$, alors $u$ est \texttt{[un]} isomorphisme. % [sic] s \in S -> s \in S_0, otherwise the hypothesis on S_0 plays no role

1$^{o}$) Soit $F$ un faisceau d'ensembles localement constant sur $X$, $F'$ son image inverse sur $X'$, il faut prouver
\[
H^0(X,F) \ \xrightarrow{\ \sim\ }\ H^0(X',F').
\]
Or
\[
H^0(X,F) \simeq H^0(Y, \underbrace{f_*(F)}_{\substack{\text{localement}\\ \text{constant sur } S}}) \ \longrightarrow\ H^0(Y, f'_*(F')) \simeq H^0(X',F'),
\] % [sic] H^0(Y, f_*(F)), H^0(Y, f'_*(F')) -> H^0(S, p_*(F)), H^0(S, p'_*(F')), the brace says "localement constant sur S" and (*) uses p_*, p'_*
il faut prouver que
\[
p_*(F) \ \xrightarrow{\ \sim\ } p'_*(F') \tag*{$(*)$}
\]
est \texttt{[un]} isomorphisme. Or par hypothèse sur $p$, $p'$, les deux faisceaux sont localement constants et se calculent fibre par fibre, d'où résulte que $(*)$ est un isomorphisme pour les fibres en $s \in S_0$, donc un isomorphisme partout. OK.

2$^{o}$) Il en résulte déjà, si $F$ est un faisceau en \textit{groupes,} que le foncteur
\[
(F\text{-torseurs sur } X) \ \longrightarrow\ (F\text{-torseurs sur } X')
\]
est pleinement fidèle, donc $H^1(X,F) \longrightarrow H^1(X',F')$ \texttt{[est]} injectif; il reste à prouver que ce foncteur est essentiellement surjectif. À cause de la pleine fidélité, c'est \textit{local sur $S$} (\textbf{NB} On voit directement que l'hypothèse fait sur les $s \in S_0$ implique la même hypothèse pour tout $s \in S$, en utilisant le fait que $p$, $p'$ sont des fibrations) (\footnote{Les hypothèses faites sont bel et bien locales sur $S$ (du moins si $S_0 = S$).}). Donc on est ramené au cas où $S = S/s_0$ a un objet final $s_0$.

\dgrblock{dgrb456}{Testing one fibre in each component}
\begin{dcomment}{exposition}{p132:est pleinement fidèle, donc}{Testing one fibre in each component}
Theorem 2 is the fibrewise Whitehead theorem for these categorical fibrations. Its fibre map is \(f_s:X'_s\longrightarrow X_s\); the displayed \(X_s\longrightarrow Y_s\) is an edition slip. Transport is natural with respect to \(f\), so if \(f_s\) is a weak equivalence at one point of a component then it is a weak equivalence at every point of that component. The choice of \(S_0\) therefore loses no information.

For sufficiency, the source compares direct images of local coefficients. Its degree-zero comparison should be written
\[
 H^0(X,F)=H^0(S,p_*F)\longrightarrow
 H^0(S,p'_*F')=H^0(X',F').
\]
Both direct images are locally constant, and their values on \(S_0\) are isomorphic by the fibre hypothesis. A morphism between locally constant diagrams is an isomorphism everywhere if it is so at one object of each component: propagate the property along an unoriented chain of arrows, whose transition maps are invertible. This proves degree zero and, by the isomorphism-sheaf argument, full faithfulness on torsors.

The local reduction for essential surjectivity is legitimate once it is applied to these groupoids of torsors, rather than only their sets of isomorphism classes. On each representable \(S/s\), both total categories still satisfy the hypotheses. The torsor diagrams constructed above are locally constant there, and \(S/s\) has a final object. Their groupoids of global sections are thus equivalent to their values at that object. Fibrewise equivalence supplies local preimages of any torsor on \(X'\). Full faithfulness uniquely lifts the overlap isomorphisms between these preimages; their cocycle follows from its validity after pullback and faithfulness. Effective descent for torsors gives the global preimage. This justifies both the phrase “local on \(S\)” and the degree-one edge identifications on manuscript page 133. The abelian comparison follows from Leray as before.

For necessity there is a shorter modern proof, complementary to the printed path-category argument. The preceding block identifies both strict fibres with homotopy fibres. An equivalence \(BX'\longrightarrow BX\) over \(BS\) induces equivalences on all homotopy pullbacks along \(\{s\}\longrightarrow BS\), so every \(Bf_s\) is a weak equivalence. Conversely, maps of spaces over a fixed base are detected on their homotopy fibres: on each base component the long exact homotopy sequences give the higher homotopy isomorphisms, and the action of \(\pi_1(BS,s)\) on the components of the fibre identifies the components of the total space with its orbit set. A fibrewise equivalence preserves that action and those orbit sets. This also treats disconnected or empty fibres, which a statement only about based groups would overlook.
\end{dcomment}


\printedpage{5}
% head: A. Grothendieck : Dérivateurs, Ch. XI, Édition des Universités de Montpellier II et Paris VII
% foot: Version 28 mai 2002

\msp{page 133}

On a alors
\[
\begin{aligned}
H^1(X,F) &\ \simeq\ H^0(S, R^1p_*(F)) \\
H^1(X',F') &\ \simeq\ H^0(S, R^1p'_*(F')),
\end{aligned}
\]
et on voit encore que
\[
R^1p_*(F) \ \longrightarrow\ R^1p'_*(F')
\]
est \texttt{[un]} isomorphisme (car il se calcule fibre par fibre, et on utilise l'hypothèse sur les $f_s$), donc induit un isomorphisme
\[
H^1(X,F) \ \longrightarrow\ H^1(X',F').
\]

3$^{o}$) Si $F$ est même abélien, on doit prouver
\[
H^i(X,F) \ \xrightarrow{\ \sim\ }\ H^i(X',F')
\]
pour tout $i$, et cela résulte encore de la suite spectrale de Leray pour $(p,F)$ et pour $(p',F')$.

b) Inversement, supposons que $f$ soit une Hot-équivalence, et prouvons que les $f_s : X'_s \longrightarrow X_s$ sont des Hot-équivalences. Je vais utiliser le

% (a vertical change bar is printed in the left margin beside the Lemme)
\textbf{Lemme.} \textit{Soit $S \in \Ob\Cat$, $S$ \textup{0}-connexe. Alors il existe un $\overline{S} \in \Ob\Cat$, $\overline{S}$ \textup{Hot}-asphérique, et $\overline{S} \longrightarrow S$ une \textup{Hot}-fibration.}

Voyons comment ceci implique le théorème 2 (partie réciproque). On peut supposer $S$ 0-connexe, et on fait le changement de base par $\overline{S} \longrightarrow S$

\msp{page 134}

comme dans le lemme :
\[
\begin{tikzcd}[column sep=small, row sep=small]
X' \arrow[rr, "f"] \arrow[dr, "p'"'] & & X \arrow[dl, "p"] & & & & \\
& S & & & \overline{X'} \arrow[rr, "\overline{f}"] \arrow[dr, "\overline{p'}"'] & & \overline{X} \arrow[dl, "\overline{p}"] \\
& & & & & \overline{S} \arrow[ullll, "\varphi"] &
\end{tikzcd}
\quad .
\]

Alors $\overline{p'}$ et $\overline{p}$ sont des Hot-fibrations, et comme $\varphi$ est une Hot-fibration, le changement de base par $\varphi$ transforme la Hot-équivalence $f$ en une Hot-équivalence $\overline{f}$. Considérons alors $\overline{s} \in \Ob\overline{S}$, et le diagramme
\[
\begin{tikzcd}
\overline{X'} \arrow[r, "\overline{f}"] & \overline{X} \\
\overline{X'}_{\overline{s}} \arrow[u, hook, "i'"] \arrow[r, "\overline{f}_{\overline{s}}"] & \overline{X}_{\overline{s}} \arrow[u, hook, "i"']
\end{tikzcd}
\qquad .
\]

\printedpage{6}
% head: A. Grothendieck : Dérivateurs, Ch. XI, Édition des Universités de Montpellier II et Paris VII
% foot: Version 28 mai 2002

Comme $\{\overline{s}\} \longrightarrow \overline{S}$ est une Hot-équivalence, la flèche $i$ qui s'en déduit par le changement de base $\overline{p} : \overline{X} \longrightarrow \overline{S}$ l'est aussi, puisque $\overline{p}$ est une Hot-fibration, et de même pour $i'$. Ainsi dans le carré commutatif précédent, $\overline{f}$, $i$ et $i'$ sont des Hot-équivalences, donc aussi $\overline{f}_{\overline{s}}$. Mais si $s = \varphi(\overline{s})$, $\overline{f}_{\overline{s}}$ s'identifie à $f_s : X'_s \longrightarrow X_s$. D'ailleurs, il est immédiat que $\varphi$, étant une Hot-fibration

\msp{page 135}

à source non-vide, et à but connexe, est surjective. (\textbf{NB} toutes les fibres d'une même composante connexe de $S$ sont Hot-équivalentes, donc si l'une est non-vide, toutes le sont.) Cela achève la démonstration, modulo le lemme.

\dgrblock{dgrb457}{A contractible replacement without a derivator}
\begin{dcomment}{exposition}{p135:Cela achève la démonstration, modulo le lemme.}{A contractible replacement without a derivator}
The lemma is not asserting the existence of a contractible covering space: a general connected homotopy type has no such covering. It asks for a fibration whose fibre retains the entire loop homotopy type. The source's based path category is designed to do exactly this. The dependency it leaves open at the end of the chapter concerns the proof that this specific categorical path construction is both proper and smooth.

For the existential lemma alone, Thomason's theory already gives an independent proof. Choose \(s_0\in S\) and factor its inclusion in the Thomason model category as
\[
 e\xrightarrow{\ i\ }E\xrightarrow{\ p\ }S,
 \qquad i\text{ a trivial cofibration},\quad
 p\text{ a model fibration}.
\]
Then \(BE\) is weakly contractible. Since \(\operatorname{Ex}^2N\) preserves pullbacks and sends \(p\) and all its base changes to Kan fibrations, right properness of simplicial sets proves that \(p\) is a Hot-fibration in the sense of a). The existence of the factorization is \citet[Theorem 4.9]{Thomason1980}; its weak equivalences and fibrations are Definitions 2.1--2.2, and properness is Corollary 5.5.

The map \(p\) is surjective on objects when \(S\) is nonempty and connected. The object \(i(*)\) supplies one nonempty fibre. Pull back over an arrow of \(S\); both endpoint fibres have the homotopy type of that pullback, so one is nonempty if and only if the other is. A zigzag in \(S\) reaches every object of its connected component. This is the elementary justification of the surjectivity used on manuscript page 135. It gives the printed proof of the converse in Theorem 2 without invoking any derivator: pull back along \(E\longrightarrow S\), use that the inclusions of its fibres are weak equivalences, and apply two-out-of-three.

This replacement does not diminish the purpose of Chapter XII. The abstract model factorization proves existence, whereas Grothendieck needs constructions with explicit categories of lifts, duality, and control of both direct images. The later smooth-and-proper factorization quoted above addresses that stronger requirement.
\end{dcomment}


Il reste à prouver celui-ci. Comme $S = \emptyset$ \texttt{[plutôt} $S \neq \emptyset$\texttt{]}, je peux choisir une ‘origine’
\[
s_0 \in S,
\]
et regarder la catégorie
\[
\Ch_{\infty}(S; s_0, -)
\]
des chemins infinis dans $S$, d'origine $s_0$, fibre de $\Ch_{\infty}(S)$ par le foncteur-source
\[
\sigma : \Ch_{\infty}(S) \ \longrightarrow\ S.
\]
On sait que les fibres de $\sigma$ sont Hot-asphériques (VII), et la catégorie construite ici n'est autre que (avec les notations de VII, p.~152) que la catégorie $\overline{S'}$ qui entre dans la factorisation canonique de
\[
f : \{s_0\} = S' \ \longrightarrow\ S
\]
en
\[
\{s_0\} = S' \ \underbrace{\longrightarrow}_{\text{Hot-équivalence}}\ \overline{S'} \ \xrightarrow{\ \overline{f}\ }\ S,
\]
où ici $\overline{f}$ s'interprète simplement comme le foncteur but (i.e.\ le foncteur induit sur

\msp{page 136}

$\Ch_{\infty}(S; s_0, -)$ par le foncteur but
\[
\beta : \Ch_{\infty}(S) \ \longrightarrow\ S \ \ ).
\]
On a prouvé dans loc.\ cit.\ que $\overline{f}$ est une Hot-fibration (‘propre’ et ‘lisse’), mais dans un contexte qui semblait faire appel tacitement à la théorie du dérivateur HOT. Donc pour mettre les résultats ‘topologiques’ de loc.\ cit.\ sur une base solide, il faut développer les notions pertinentes de foncteur propre, foncteur lisse, sans référence à une théorie HOT. C'est ce que j'avais d'ailleurs fait déjà en 1983, du temps de Pursuing Stacks, mais sans rédaction en forme à ce sujet. Il me faut donc à présent revenir sur ces notions, de façon un peu systématique. D'ailleurs, j'aurais dû mal à construire un seul foncteur Hot-fibrant (sauf les exemples triviaux, tels les isomorphismes, ou les produits $Y \times Z \xrightarrow{\mathrm{pr}_1} Y$), si je ne disposais d'une théorie des foncteurs lisses et des foncteurs propres (\footnote{J'exagère un peu quand même -- après tout il doit être assez évident que les foncteurs de Thomason-Kan sont des Hot-fibrations (mais il y en a encore un peu trop peu !).}). % [sic] dû mal -> du mal, "avoir du mal à" (partitive article, no accent)

\dgrblock{dgrb458}{Why the next chapter starts again}
\begin{dcomment}{context}{p136:Il me faut donc à présent revenir sur ces notions}{Why the next chapter starts again}
Grothendieck's return to proper and smooth functors is an effort to remove a dependency, not a rejection of the homotopy theory just obtained. A path construction had been used to produce fibrations while some of its asserted properties still appealed to the derivator that those fibrations were meant to construct. Categories of lifts offer a direct combinatorial foundation for those properties.

The historical background is visible in \emph{Pursuing Stacks}, \S69: the coefficient theory is to remember categories of coherent diagrams together with both adjoints of restriction, rather than merely a single homotopy category \citep{Grothendieck1983}. The letter to Thomason of 2 April 1991 makes a related distinction explicit: the weak equivalences carry the intrinsic homotopy theory, while classes of fibrations and cofibrations make its operations calculable \citep[\S1]{Grothendieck1991Letter}. We refer to the commentary on Chapter~I, manuscript p.~126 (“The later letter and the meaning of a fibration”), for the broader historical discussion.

The next pages isolate two different issues. One is to describe the maps along which categorical pullback behaves well. The other is to decide how much of the usual homotopy theory is forced by the proposed axioms on \(W\). The revisions of terminology in Chapter XII are necessary because an axiom suitable for the full homotopy type can be too strong for a localization that deliberately forgets higher homotopy groups.
\end{dcomment}
\dgrchapter{XII}{Chapter XII}{\texorpdfstring{Caractérisation de $W_\infty$. Foncteurs $W$-propres, $W$-lisses etc. Sommes amalgamées et carrés $W$-cocartésiens dans Cat}{Caractérisation de W∞. Foncteurs W-propres, W-lisses etc. Sommes amalgamées et carrés W-cocartésiens dans Cat}}
\dgrpart{0}{XI-XIIa}{Sections 1--4, manuscript pages 1--58}{XI-XIIa\enspace\textperiodcentered\enspace Chapter XII, sections 1--4\enspace\textperiodcentered\enspace manuscript pages 1--58}{Chapter XII \textperiodcentered\ \textup{XI-XIIa}}
% ---- from Der-ChXIIa-012_p001-005.tex
\printedpage{--}
% head: (none; no running head and no page number printed on this page)

\begin{center}
{\Large \textsf{LES DÉRIVATEURS}}

\bigskip

\textbf{ALEXANDRE GROTHENDIECK}

\bigskip\bigskip

\textbf{Chapitre XII}

\medskip

\textbf{Caractérisation de $W_\infty$. Foncteurs $W$-propres,\\
$W$-lisses etc. Sommes amalgamées et carrés\\
$W$-cocartésiens dans $\Cat$.}
\end{center}

\clearpage

\printedpage{--}
% head: (none; no running head and no page number printed on this page)

Ce texte a été déchiffré et transcrit en \LaTeXe{} par M. Künzer. Il a été édité par M. Künzer, J. Malgoire, et G. Maltsiniotis. La transcription est aussi fidèle que possible au manuscrit. Pour les quelques corrections évidentes, ou rares commentaires des éditeurs, ainsi que pour la numérotation originale des pages du manuscrit, les caractères de machine à écrire \texttt{[typewriter]} entre crochets sont utilisés. Un point d'interrogation entre crochets signifie que l'on n'est pas sûr du mot qui précède.

Cette édition est provisoire. Les remarques, commentaires et corrections sont bienvenus. Envoyer un message à :

\textbf{[e-mail address removed]}

\hfill G. Maltsiniotis\hspace*{2em}

\clearpage

\printedpage{1}
% head: A. Grothendieck : Dérivateurs, Ch. XII, Édition des Universités de Montpellier II et Paris VII
% foot: Version 7 décembre 2002

\msp{page 1}

\hfill 19.11.90

\section*{I\quad Caractérisation de $W_\infty$.}

Je voudrais caractériser $W_\infty$ en tant que partie de $\mathrm{Fl}(\Cat)$, par ses propriétés comme localiseur, et sans faire appel à la connaissance de notions utilisant déjà $W_\infty$, comme la propreté etc. Je voudrais tout au moins des propriétés sur un localiseur $W$, familières pour $W_\infty$, qui soient si fortes qu'elles impliquent $W \subseteq W_\infty$. Bien sûr, dans les démonstrations je me permets d'utiliser tout ce que je sais sur $W_\infty$.

Bien sûr, il faudra au moins une forme de saturation, peut-être très faible, je présume le suivant suffira (\footnote{(non, j'ai dû le renforcer en 2 [?]).}).

\begin{itemize}[leftmargin=*]
\item[\textbf{W(1)}] (Axiome de saturation.) a) Les identités sont dans $W$, b) si $f$, $g$ \texttt{[sont]} composables et deux parmi $f$, $g$, $gf$ sont dans $W$, le troisième aussi, c) si $fg \in W$, $gf \in W$, alors $f, g \in W$.
\end{itemize}

Une propriété particulièrement cruciale, qui déjà distingue $W_\infty$ de l'ensemble des homotopismes (notion trop stricte à beaucoup d'égards), est l'axiome que j'ai appelé L4 bis :

\dgrblock{dgrb459}{A contractible nerve without a categorical contraction}
\begin{dcomment}{examples}{p1:Une propriété particulièrement cruciale, qui déjà distingue}{A contractible nerve without a categorical contraction}
Grothendieck's insistence that homotopisms are too restrictive has a concrete explanation. Let \(P\) be the ordered set with objects \(0,1,2,\ldots\) and nonidentity relations
\[
  0<1>2<3>4<5>\cdots .
\]
There are no composable pairs of nonidentity arrows. Its nerve therefore realizes to the half-line subdivided at the integers; in particular \(P\to e\) belongs to \(W_\infty\).

Nevertheless \(P\to e\) is not a homotopism in the sense used on the next manuscript page. Suppose that the identity of \(P\) could be joined to the constant functor at \(c\) by a zigzag of \(m\) natural transformations. Evaluating the zigzag at an integer \(n\) would give a path of at most \(m\) edges from \(n\) to \(c\) in the underlying undirected graph of \(P\). Every arrow changes the integer by at most one. Thus \(|n-c|\leq m\) for every \(n\), which is impossible. A homotopy inverse to \(P\to e\) would select precisely such a constant functor, so no homotopy inverse exists.

This example separates two finite conditions that are easily conflated. Each individual vertex can be connected to \(0\) by a finite path. A natural-transformation contraction would require one finite zigzag working simultaneously at every vertex. The localization axioms allow the classifying space to forget that uniform bound.
\end{dcomment}


\begin{itemize}[leftmargin=*]
\item[\textbf{W(2)}] (Axiome de localisation.) Soit $f : X \fl Y$ une flèche dans $\Cat/S$, telle que $\forall\ s \in \Ob S$, $f_s : X/s \fl Y/s$ soit dans $W$ (on dira que $f$ est dans $W$ `localement sur $S$\,'). Alors $f \in W$.
\end{itemize}

\dgrblock{dgrb460}{The two directions of the characterization}
\begin{dcomment}{exposition}{p1:soit dans $W$ (on dira que $f$ est dans $W$}{The two directions of the characterization}
Grothendieck wants the equivalences of ordinary homotopy theory to emerge from operations on categories. There are two logically different assertions to establish. The first says that the proposed axioms cannot identify too many maps: \(W\subseteq W_\infty\). The second says that a sufficiently coherent theory cannot identify too few: \(W_\infty\subseteq W\). This section attacks the first assertion by repeatedly taking fibres. The later minimal-localizer theorem supplies the second. Keeping these directions separate makes the successive changes of axioms intelligible.

Write \(BA=|NA|\) for the classifying space of a small category, and let \(e\) be the terminal category. Here \(W_\infty\) consists of the functors \(u\) for which \(Bu\) is a weak homotopy equivalence. The expression \(W\)-aspherical means only \(A\to e\in W\). It does not mean that \(A\) has vanishing higher homotopy groups until a comparison with \(W_\infty\) has been established. In particular, it is not the usual topological convention under which a connected \(K(G,1)\) is called aspherical.

The local-to-global axiom W(2) tests a map \(u:X\to Y\) over \(S\) on the comma categories \(X/s\to Y/s\). These are not the strict fibres \(X_s\to Y_s\). An object of \(X/s\) is a pair \((x,p(x)\to s)\), so the test retains all ways of approaching \(s\). W(2 bis) is the absolute theorem-A form: contractibility, in the sense of \(W\), of every \(X/y\) implies \(u\in W\). The distinction matters: the source explicitly uses the relative form for Proposition~3, whereas the loop argument and the final nondegeneracy repair need only the absolute form.

The commentary on Chapter~VI, manuscript pp.~44--45, explained why a fibre axiom begins to select ordinary homotopy theory. The present pages give the mechanism: any nonzero homotopy group of a \(W\)-aspherical category can be lowered, by looping, to a forbidden nontrivial set of components.
\end{dcomment}


\texttt{[Les axiomes W(1) etc. sont appelés aussi} \textcircled{1} \texttt{etc. dans l'original.]}

Dans le cas $Y = S$, cela donne

\textbf{Corollaire :} \textit{Si $f : X \fl Y$ est $W$-asphérique, i.e. $y\backslash X \fl y\backslash Y$} \texttt{[plutôt $X/y \fl Y/y$]} \textit{dans $W$ pour tout $y \in \Ob Y$, alors $f \in W$} (\footnote{On peut considérer la validité du corollaire comme une variante affaiblie W(2 bis) de W(2) ci-dessus. Dans la proposition 1, on n'utilise que W(2 bis), forme faible de W(2).}).

\textbf{Proposition 1.} \textit{Si $W$ satisfait} W(1) \textit{et} W(2)\textit{, les propriétés suivantes sur $W$ sont équivalentes :}

\begin{itemize}[leftmargin=*]
\item[\textit{a)}] \textit{Pour tout $X$ dans $\Cat$, $X \times \Dl^1 \xrightarrow{\pr_1^X} X$ est dans $W$ (i.e. $\Dl^1 \fl e$ est `universellement dans $W$\,').}

\item[\textit{b)}] \textit{Si $X$ dans $\Cat$ a un objet final, alors $X$ est $W$-asphérique, i.e. $X \fl e$ est dans $W$.}

\item[\textit{c)}] \textit{Tout homotopisme est dans $W$.}
\end{itemize}

\clearpage

\printedpage{2}
% head: A. Grothendieck : Dérivateurs, Ch. XII, Édition des Universités de Montpellier II et Paris VII
% foot: Version 7 décembre 2002

\textsc{Démonstration.} b) $\Longrightarrow$ a) (\footnote{n'utilise W(1) que sous forme faible, sans la condition de saturation c). NB Seul b) $\Longrightarrow$ c) utilise W(2 bis), les implications a) $\Longrightarrow$ c) $\Longrightarrow$ b) sont valables sous la seule hypothèse W(1).}), car pour voir que $\pr_1 \in W$, en appliquant le corollaire ci-dessus, il faut prouver que $\forall\ x \in X$
\[ (X \times \Dl^1)/x \fl X/x \tag*{$(*)$} \]
est dans $W$. Or
\[ (X \times \Dl^1)/x \simeq X/x \times \Dl^1, \]
\[
\begin{tikzcd}
\Dl^1 \arrow[d] & X \times \Dl^1 \arrow[l] \arrow[d] & X/x \times \Dl^1 \arrow[l] \arrow[d] \\
e & X \arrow[l] & X/x \arrow[l]
\end{tikzcd}
\]
\texttt{[la flèche verticale à droite est ajoutée]} et les deux catégories $X/x$, $X/x \times \Dl^1$ sont $W$-asphériques car elles ont un objet final, donc $(*)$ est dans $W$ par saturation faible W(1).

a) $\Longrightarrow$ c). Je rappelle que
\[ f : X \fl Y \]
\msp{page 2}
est un homotopisme s'il existe
\[ g : Y \fl X \]
telle que
\[ gf \sim \id_X, \quad fg \sim \id_Y, \tag*{$(*)$} \]
où la relation $(\sim)$ désigne l'homotopie dans $\uHom(X,X)$ et $\uHom(Y,Y)$ respectivement (définie par chemins allant d'un élément à l'autre). Prouvons le

\textbf{Lemme.} \textit{Supposons que $W$ satisfait} W(1) \textit{et le} a) \textit{de la proposition, i.e. $X \times \Dl^1 \fl X$ dans $W$ pour tout $X$ dans} $\Cat$. \textit{Alors si $X$, $Y$} \texttt{[sont]} \textit{dans} $\Cat$ \textit{et $f, g$ dans} $\uHom(X,Y)$ \textit{sont homotopes, on a}
\[ \hot_W(f) = \hot_W(g). \]
\textit{De plus, $f \in W$ si et seulement si $g \in W$.}

\dgrblock{dgrb461}{Saturation at the point where it is used}
\begin{dcomment}{exposition}{p2:si et seulement si $g \in W$.}{Saturation at the point where it is used}
The proof is carefully arranged so that equality after localization is not mistaken for membership in \(W\). For a natural transformation \(f\Rightarrow g\), let \(F:X\times[1]\to Y\) be its corresponding functor and let \(p:X\times[1]\to X\) be the projection. If \(p\in W\), its sections \(i_0,i_1\) lie in \(W\) by two-out-of-three. Since \(pi_0=pi_1=\id_X\), their images in \(\Cat[W^{-1}]\) are both \(p^{-1}\). Thus \(f=Fi_0\) and \(g=Fi_1\) become equal there. Independently, two-out-of-three gives
\[
 f\in W\quad\Longleftrightarrow\quad F\in W
 \quad\Longleftrightarrow\quad g\in W.
\]
Iteration handles a finite zigzag of transformations.

Now suppose \(f\) has a categorical homotopy inverse \(g\). The preceding argument puts both \(gf\) and \(fg\) in \(W\). The separate clause W(1c) then puts \(f\) and \(g\) in \(W\). This last step would not follow merely from the existence of their inverse images in the localization unless strong saturation had already been proved.

Three usages of saturation should therefore remain distinct. Two-out-of-three is W(1b); the two-cycle condition is W(1c); strong saturation says that every map inverted by the localization already belongs to \(W\). The modern weak-saturation axiom used for fundamental localizers asks, in addition to identities and two-out-of-three, only that \(ri=\id\) and \(ir\in W\) imply \(r\in W\) \citep[Definition~1.1.1]{Cisinski2004Minimal}. W(1c) implies this condition. Consequently the later results apply to the source's axioms without silently inserting strong saturation into this elementary proof.
\end{dcomment}


Il suffit de le prouver quand il existe un homomorphisme $f \fl g$. Cela peut s'interpréter par l'existence d'une flèche
\[ X \times \Dl^1 \xrightarrow{F} Y \]
dont les restrictions à $X \times \{0\}$ et $X \times \{1\}$ respectivement soient $f$, $g$. Pour prouver que $\hot_W(f) = \hot_W(g)$, il suffit de prouver que
\[ i_0, i_1 : X \fl X \times \Dl^1 \qquad x \longmapsto (x,0),\ x \longmapsto (x,1) \]
\printedpage{3}
% head: A. Grothendieck : Dérivateurs, Ch. XII, Édition des Universités de Montpellier II et Paris VII
% foot: Version 7 décembre 2002
sont égales dans $\Hot_W$, puisque
\[ f = F \circ i_0, \quad g = F \circ i_1. \]
Or cela résulte du fait que ce sont des sections de $X \times \Dl^1$ sur $X$, et que dans $\Hot_W$ la projection $X \times \Dl^1 \xrightarrow{p} X$ est \texttt{[un]} isomorphisme (donc $i_0 = i_1 = p^{-1}$ dans $\Hot_W$). De plus, le diagramme
\[
\begin{tikzcd}
X \arrow[r, "i_0"] \arrow[dr, "f"'] & X \times \Dl^1 \arrow[d, "F"] \\
& Y,
\end{tikzcd}
\]
où $i_0 \in W$, montre que $f \in W$ si et seulement si $F \in W$, et de même $g \in W$ si et seulement si $F \in W$, donc $f \in W \Longleftrightarrow g \in W$.

Revenons à $(*)$, on voit que dans $\Hot_W$ $g \circ f = \id_X$, $f \circ g = \id_Y$, donc $f$ et $g$ sont des $\Hot_W$-isomorphismes inverses l'un de l'autre. D'autre part, $gf \in W$, $fg \in W$, puisque $\id_X, \id_Y \in W$. Donc on gagne, si on a formulé W(1) de façon assez forte. On a prouvé donc a) $\Longrightarrow$ c).

c) $\Longrightarrow$ b) est trivial, puisque si $X$ a un objet final, $X \fl e$ est un homotopisme.

\msp{page 3}

On doit bien sûr par la suite supposer que $W$ satisfait aux conditions équivalentes de la proposition. On prendra p. ex. la forme

\begin{itemize}[leftmargin=*]
\item[\textbf{W(3)}] (Axiome de l'élément final.) Si $X$ a un objet final, alors $X$ est $W$-asphérique.
\end{itemize}

\dgrblock{dgrb462}{Terminal objects, initial objects, and duality}
\begin{dcomment}{exposition}{p3:alors $X$ est $W$-asphérique}{Terminal objects, initial objects, and duality}
The terminal-object axiom gives a usable supply of trivial categories. If \(t\) is terminal in \(A\), the natural transformation \(\id_A\Rightarrow \operatorname{const}_t\) makes \(A\to e\) a homotopism; an initial object gives the transformation in the reverse direction. Proposition~1 therefore derives initial-object asphericity from the same axioms. No second axiom is being assumed.

For a functor \(u:A\to B\), the category \(B/b\) has the terminal object \((b,\id_b)\). Hence two-out-of-three identifies the two formulations
\[
 A/b\longrightarrow B/b\in W
 \quad\Longleftrightarrow\quad A/b\longrightarrow e\in W .
\]
This is why the corrected slice orientation in the source's first corollary is essential: it determines whether the subsequent notion is asphericity or coasphericity.

Proposition~2 next turns this one-sided language into a self-dual theory. The bridge is the twisted-arrow category \(\Phi(A)\), with comparison functors to \(A\) and \(A^{\mathrm{op}}\). Their comma categories have explicit contractions, so absolute theorem A puts both comparisons in \(W\). Naturality then compares \(u\), \(\Phi(u)\), and \(u^{\mathrm{op}}\) by two-out-of-three. The construction and its variance were treated in the commentary on Chapter~I, manuscript pp.~73--75, and are not repeated here. What is new in the present argument is its economy of hypotheses: the comparisons use W(2 bis), not the full relative W(2). Applying the same construction to \(u^{\mathrm{op}}\) gives the reverse implication, so \(u\in W\) if and only if \(u^{\mathrm{op}}\in W\).
\end{dcomment}


\textbf{Corollaire 1.} \textit{Pour que $f : X \fl Y$ soit $W$-asphérique, il faut et il suffit que les $X/y$ soient $W$-asphériques} ($y \in \Ob Y$).

Car pour que
\begin{tikzcd}[baseline=(\tikzcdmatrixname-1-1.base), column sep=small, row sep=small]
X/y \arrow[rr] \arrow[dr] & & Y/y \arrow[dl] \\
& e &
\end{tikzcd}
soit dans $W$, comme $Y/y$ est $W$-asphérique (objet final), il faut et il suffit que $X/y$ le soit aussi.

\textbf{Corollaire 2.} \textit{Si $X$ dans} $\Cat$ \textit{a un objet initial, il est $W$-asphérique.} (En effet, $X \fl e$ est un homotopisme, et on applique prop. 1.) (\footnote{C'est donc l'énoncé dual de l'axiome W(3) de l'objet final.})

Je mettrai pour mémoire

\textbf{Corollaire 3.} \textit{Tout homotopisme dans} $\Cat$ \textit{est dans $W$. Si $f, g : X \rightrightarrows Y$ dans} $\Cat$ \textit{sont homotopes, alors elles sont égales dans} $\Hot_W$\textit{, et $f$ est dans $W$ si et seulement si $g$ l'est.}

\textbf{Proposition 2.} Soit $f : X \fl Y$ dans $\Cat$. Si $f \in W$, $f^\circ \in W^\circ$ \texttt{[plutôt $f^\circ \in W$].}

% ---- from Der-ChXIIa-012_p006-010.tex
\printedpage{4}
% head: A. Grothendieck : Dérivateurs, Ch. XII, Édition des Universités de Montpellier II et Paris VII
On utilise la démonstration de I, p. 75, à l'aide de la catégorie nommée $\Phi(X)$ ($\mathrm{Ob}\,\Phi(X) = \mathrm{Fl}(X)$, mais avec une notion de flèche dans $\Phi(X)$ un peu particulière). On trouve que $f \in W \iff \Phi(f) \in W$, et de même $f^\circ \in W \iff \Phi(f) \in W$.

NB. Dans la démonstration de la proposition 2, comme dans celle de la proposition 1 (pour l'implication b) $\Longrightarrow$ a)), on n'a utilisé la propriété de localisation W(2) que sous la forme faible du corollaire, i.e. sous la forme $f$ $W$-asphérique $\Longrightarrow f \in W$.

\textbf{Corollaire} (de la prop. 2) : \textit{Soit $W \subseteq \mathrm{Fl}(\mathrm{Cat})$ satisfaisant l'axiome} W(1) (saturation). \textit{Alors la conjonction des axiomes} W(2), W(3) \textit{équivaut à celle des axiomes duaux} W($2'$), W($3'$) (\footnote[5]{Itou si on remplace dans cet énoncé W(2) par W(2 bis).}).

Il suffit de prouver que 1,2,3 $\Longrightarrow$ $2'$,$3'$.

\msp{page 4}

Pour $2'$ \texttt{[plutôt} $3'$\texttt{]}, c'est le corollaire 2 ci-dessus. Pour $3'$ \texttt{[plutôt} $2'$\texttt{]}, cela résulte aussitôt de la proposition 2. (NB Cet énoncé resterait valable, si on aurait pris W(2) sous la forme faible seulement, cf. le NB page précédente.)

Je dois encore mettre un axiome de la forme L5, L6 (cf. VI, p. 43-44). Mais l'énoncé de ces axiomes utilise la notion de Hot-fibrés, ou de foncteurs parfaits, i.e. fait appel à la connaissance de $W_\infty$, ce que j'aimerais éviter (si je peux). Je dois donc introduire la notion de $W$-fibration, sur le modèle de ce que j'ai fait dans le contexte des dérivateurs, dans I (cf. p. 111-126). On introduit la notion de \textbf{D}-fibration à la page 117 (prop. 5), par \texttt{[des]} conditions équivalentes, dont certaines ont un sens en termes de la donnée de $W$ sans plus. Soit donc
\[ f : X \longrightarrow Y \]
dans Cat et considérons les conditions

\dgrblock{dgrb463}{What the four fibration tests quantify over}
\begin{dcomment}{exposition}{p4:Je dois encore mettre un axiome de la forme}{What the four fibration tests quantify over}
Grothendieck is looking for a notion of family whose strict pullbacks compute homotopy-invariant changes of base. His strongest condition says that
\[
  f^*:\Cat/Y\longrightarrow\Cat/X
\]
preserves the underlying \(W\)-maps. The intermediate condition tests this only between \(W\)-aspherical bases. The left-hand test asks that, for every \(v:T\to Y\) whose domain \(T\) has a terminal object and every object \(t\in T\), the inclusion of the strict fibre
\[
   X_{v(t)}\longrightarrow X\times_Y T
\]
belong to \(W\). The chosen \(t\) need not be the terminal object. The right-hand test uses bases with an initial object. This universal quantification over \(T\to Y\) distinguishes the definition from an isolated comparison \(X_y\to X/y\).

The arrows in the source's diagram follow at once because a terminal or initial category is \(W\)-aspherical. Their converses express additional mathematics. For \(W_\infty\), Chapter~XI supplies them. For a general \(W\), the step from a one-sided test to invariance under every \(W\)-map is precisely the proposed axiom W(4). One should not import it from the formalism of a model category: an arbitrary categorical functor is not a fibration of a model structure merely because it satisfies one of these tests.

The eventual terminology is settled only at XII, manuscript pp.~51--52. The early adjective \emph{ultraforte} names preservation of all \(W\)-maps, and it is useful to retain that descriptive meaning while reading these pages. The later \(W\)-fibration is defined using categorical homotopisms in the base. In a truncated theory the distinction survives: the path family over an \((n+1)\)-sphere passes the ordinary homotopy tests but does not preserve \(W_n\)-equivalences. The complete computation is given below at manuscript pp.~16--17.
\end{dcomment}


\begin{enumerate}
\item[a)] Si $Y'' \longrightarrow Y'$ dans $\mathrm{Cat}/Y$ est dans $W$, alors $X'' \longrightarrow X'$ aussi (le changement de base par $X \longrightarrow Y$ transforme $W$ dans $W$).
\item[a$'$)] Comme a), mais avec $Y''$, $Y'$ $W$-asphériques.
\item[a$_1$)] Comme a), dans le cas particulier où $Y'$ a un objet final et $Y''$ est la catégorie ponctuelle.
\item[a$_2$)] Dual de a$_1$), avec objet initial de $Y'$.
\end{enumerate}

On a donc
\[
\begin{tikzcd}[row sep=large, column sep=1.4em]
 & & \mathrm{a} \arrow[d, Rightarrow] & \text{$W$-fibration ultraforte} & \\
 & & \mathrm{a}' \arrow[dl, Rightarrow] \arrow[dr, Rightarrow] & \text{$W$-fibration forte} & \\
\text{$W$-fibration gauche} & \mathrm{a}_1 & & \mathrm{a}_2 & \text{$W$-fibration droite}
\end{tikzcd}
\]
% foot: Version 7 décembre 2002
\printedpage{5}
% head: A. Grothendieck : Dérivateurs, Ch. XII, Édition des Universités de Montpellier II et Paris VII
et dans le cas $W = W_\infty$ toutes ces conditions sont équivalentes. J'ai envie \texttt{[de]} prendre la plus forte de ces conditions, mais je dois me donner des moyens pour trouver des $W$-fibrations en un sens aussi fort. Par exemple j'aurais envie de savoir que si $X$
\msp{page 5}
est dans Cat, alors $\Ch_\infty(X) \longrightarrow X \times X$ ou $\Ch(X) \longrightarrow X \times X$ est une $W$-fibration. Dans les deux cas, et (disons) pour $W_\infty$, je le prouve en prouvant que $p$ est à la fois lisse, et propre. Il faudrait donc revoir les notions de lissité et de propreté dans le contexte de la donnée d'un $W$.

Si je me rappelle bien, la propriété a$'$) \texttt{[plutôt a}$_1$\texttt{)]} était quelque chose d'assez aisément vérifiable, il suffisait de supposer que les
\[ X_y \longrightarrow X/y \qquad \text{et} \qquad X/y \longrightarrow X/y' \]
(pour $y \in \mathrm{Ob}\,Y$ resp. $u : y \longrightarrow y'$ dans $\mathrm{Fl}\,Y$) sont dans $W$, et que ça reste vrai par tout changement de base - ce qui revient à la condition a$_1$) sur $f$. On pourrait donc formuler une condition sur $W$ ainsi :

\begin{enumerate}[leftmargin=*]
\item[\textbf{W(4)}] (Axiome des $W$-fibrations (\footnote[6]{Voir variante affaiblie suffisante, dans le NB p. 10. Voir encore [?] de variantes pages 13, 14 ci-dessous.}).) Soit $f : X \longrightarrow Y$ dans Cat satisfaisant la condition de $W$-fibration a$_1$) ci-dessus (p. 4), alors elle satisfait même à a), i.e. le changement de base par $f$ transforme flèches dans $W$ en flèches dans $W$ (\footnote[7]{i.e. une $W$-fibration à gauche est une fibration ultraforte.}).
\end{enumerate}

Grâce à la proposition 2, il s'ensuit que la condition duale a$_2$) sur $f$ implique la condition de $W$-fibration forte, donc l'axiome W(4), en apparence asymétrique, signifie aussi que les quatre conditions de fibration sur $f$ sont équivalentes, il est donc en fait symétrique.

\textbf{NB} Je ne vois pas, sans axiomes supplémentaires sur $W$, comment passer de a$_1$ à a$'$ seulement. La difficulté, c'est de montrer, moyennant a$_1$ seulement, que si $Y$ est $W$-asphérique, l'inclusion $X_y \longrightarrow X$ (où $y \in \mathrm{Ob}\,Y$) est dans $W$. Dans le contexte des dérivateurs \textbf{D}, cela se voyait par voie \textbf{D}-cohomologique.

Cet axiome très fort suffit à impliquer que si $X$, $Y$ sont $W$-fibrés sur $S$, et si $f : X \longrightarrow Y$ est un $S$-morphisme qui est dans $W$, alors $f_s : X_s \longrightarrow Y_s$ est dans $W$ ($\forall\ s \in S$). Montrons déjà au moins l'inverse.

\msp{page 6}

\textbf{Proposition 3.} \textit{Soient $X$, $Y$ dans} Cat$/S$\textit{, $f : X \longrightarrow Y$ une $S$-flèche. Supposons ($X$, $Y$ $W$-fibrés) que $\forall\ s \in S$, $f_s : X_s \longrightarrow Y_s$ soit dans $W$, alors $f \in W$} (\footnote[8]{\textit{Sans} axiome W(4). C'est la première fois qu'on a besoin de W(2) sous sa forme forte, à l'exclusion de W(2 bis).}).

Il s'impose d'utiliser l'axiome de localisation, donc de regarder les $f/s$ induits par $f$
% foot: Version 7 décembre 2002
\printedpage{6}
% head: A. Grothendieck : Dérivateurs, Ch. XII, Édition des Universités de Montpellier II et Paris VII
\[
\begin{tikzcd}[row sep=large]
X/s \arrow[r, "f/s"'] & Y/s \\
X_s \arrow[u, "i_s^X"] \arrow[r, "f_s"] & Y_s \arrow[u, "i_s^Y"'] \quad .
\end{tikzcd}
\]

\dgrblock{dgrb464}{From strict fibres to the total category}
\begin{dcomment}{exposition}{p6:Il s'impose d'utiliser l'axiome de localisation}{From strict fibres to the total category}
Proposition~3 is a descent statement, and the comma categories supply the descent data. Suppose \(p:X\to S\) and \(q:Y\to S\) satisfy the left fibration test and \(u:X\to Y\) lies over \(S\). Pulling \(p\) back to \(S/s\), with terminal object \((s,\id_s)\), shows that the canonical inclusion \(X_s\to X/s\) is in \(W\); similarly for \(Y\). The square in the text therefore gives
\[
 u_s\in W\quad\Longleftrightarrow\quad u/s\in W .
\]
If this holds at every \(s\), the relative axiom W(2) gives \(u\in W\). The dual proof applies when both maps satisfy the right test. The maps must satisfy a common variance; a mixed pair is not covered by this argument.

There is also a direct reason why one representative per connected component suffices. For an arrow \(a:s\to t\), the left fibration condition makes \(X/s\to X/t\) and \(Y/s\to Y/t\) weak equivalences. Comparing the two horizontal maps in the resulting square shows that \(u_s\in W\) exactly when \(u_t\in W\). Propagate this statement along a finite zigzag of arrows joining the chosen representative to any other object.

None of these steps uses W(4). That distinction is the point of the proposition: fibrewise equivalences imply a total equivalence by local-to-global descent, while a total equivalence need not be detectable on fibres without a stronger hypothesis on the theory. The path-space example below will show the failure sharply.
\end{dcomment}


La condition a$_1$ suffit à impliquer que $i_s^X$ et $i_s^Y$ sont dans $W$, donc $f/s$ est dans $W$ si et seulement si $f_s$ l'est. On n'a pas eu à utiliser l'axiome W(4), et il suffisait de supposer que $X$, $Y$ soient fibrés ‘$W$-cohomologiques’ (ou $W$-fibrés à gauche) sur $S$, ou dualement, des $W$-fibrés à droite.

De plus :

\textbf{Corollaire} (\footnote[9]{Sans axiome W(4).}). \textit{Soit $S_0 \subseteq \mathrm{Ob}\,S$ une partie de $S$ telle que $S_0 \longrightarrow \pi_0(S)$ surjectif. Alors dans la proposition 3 il suffit de supposer $f_s \in W$ pour $s \in S_0$.}

Il suffit de voir que si $u : s \longrightarrow s'$ est une flèche dans $S$, alors $f_s \in W \iff f_{s'} \in W$. On utilise le fait que $f_s \in W \iff f/s \in W$, et de même pour $f_{s'}$. Donc à voir que $f/s \in W \iff f/s' \in W$. Or on a
\[
\begin{tikzcd}[row sep=large]
X/s \arrow[r, "f/s"] \arrow[d, "X/u"'] & Y/s \arrow[d, "Y/u"] \\
X/s' \arrow[r, "f/s'"] & Y/s' \quad ,
\end{tikzcd}
\]
où $X/u$ et $Y/u$ sont dans $W$, d'où le résultat.

Mais revenons à la réciproque : si $f \in W$, alors a-t-on $f_s \in W_S$? \texttt{[plutôt} $f_s \in W$\texttt{?]} Il y a un argument dans ce sens dans VI, p. 23, que je vais reprendre. L'idée est celle-ci. On considère la composante connexe $Y_0$ de $Y$ \texttt{[plutôt} $S_0$ \texttt{de} $S$\texttt{]} contenant $s$, et on construit un
\msp{page 7}
$W$-fibré (au \textit{sens fort} a))
\[ S' \longrightarrow S_0 \subseteq S \]
avec $S'$ \textit{$W$-asphérique.} Comme $X \longrightarrow Y$ sur $S$ est dans $W$, il en sera donc de même de $X' \longrightarrow Y'$ sur $S'$. Comme les fibres de $X'$, $Y'$ en $s'$ s'identifient à celles de $X$, $Y$ en l'image de $s'$ dans $S$, on est ramené au cas où $S$ est $W$-asphérique. Mais alors on a
\[
\begin{tikzcd}[row sep=large]
X \arrow[r, "f"] & Y \\
X_s \arrow[u, "i_s^X"] \arrow[r, "f_s"] & Y_s \arrow[u, "i_s^Y"'] \quad ,
\end{tikzcd}
\]
% foot: Version 7 décembre 2002
\printedpage{7}
% head: A. Grothendieck : Dérivateurs, Ch. XII, Édition des Universités de Montpellier II et Paris VII
où $i_s^X$, $i_s^Y$ sont dans $W$, car déduites de $\{s\} \longrightarrow S$ (qui est dans $W$, $S$ étant $W$-asphérique) par les changements de base $W$-fibrants $X \longrightarrow S$, $Y \longrightarrow S$. Donc $f \in W \implies f_s \in W$.

NB On n'a utilisé la propriété de fibration pour $X$, $Y$ qu'au sens a$'$) (\footnote[10]{Et même, dans a$'$) on pourrait exiger que $Y'' = e$, $Y'$ $W_\omega$-asphérique, où $W_\omega$ \texttt{[est]} défini page 14, $W_\omega \subseteq W \cap W_\infty \ldots$}).

Il faut donc prouver, si on peut, le

\textbf{Lemme} (?) (\footnote[11]{Est prouvé moyennant l'axiome W(4) (+ des petites vérifications ultérieures).}). \textit{Soit $S$ un objet} $0$\textit{-connexe de} Cat\textit{. Alors il existe un objet $W$-asphérique $S'$ de} Cat\textit{, et une $W$-fibration (au sens} fort \textit{de a), page 4) $S' \longrightarrow S$.}

\dgrblock{dgrb465}{A path cover makes the converse local}
\begin{dcomment}{exposition}{p7:une $W$-fibration (au sens}{A path cover makes the converse local}
The converse fibre test needs a cover on which the base becomes aspherical without changing the homotopy type of the fibres. Fix \(s_0\) in a connected component of \(S\), and use the endpoint map
\[
 \beta:\Ch(S;s_0,-)\longrightarrow S .
\]
Its source is intended to model based paths with free endpoint. Every object in the chosen component is the endpoint of a finite zigzag starting at \(s_0\), so this cover reaches each object of that component.

Assume the two properties under construction: \(E=\Ch(S;s_0,-)\) is \(W\)-aspherical and \(\beta\) is ultraforte. If \(u:X\to Y\) over \(S\) is in \(W\), then the pullback \(u_E:X_E\to Y_E\) is in \(W\). For any \(e_0\in E\), the point map \(e\to E\) is in \(W\), by two-out-of-three with \(E\to e\). The strong-aspherical-base condition for \(X_E\) and \(Y_E\) now places both fibre inclusions in \(W\). Their comparison square gives \(u_{\beta(e_0)}\in W\).

This formulation avoids an unnoticed component restriction: one pulls back \(u\) along the actual composite \(E\to S\). The inclusion \(S_0\to S\) of the selected component is a pullback of the point inclusion \(e\to\pi_0S\). A functor to a discrete category is perfect, so this inclusion is perfect as well; W(4) makes it preserve \(W\) under base change. The composite of this inclusion with the perfect path cover has the required property. No prior assertion that restriction of an arbitrary \(W\)-map to one component remains in \(W\) is needed here.

The dependency that remains is real but explicitly acknowledged: asphericity of \(E\) must be proved before the converse fibre theorem is used. Deriving it from the desired conclusion about fibres of \(\Ch(S)\to S\) would be circular. The next paragraph changes the base from \(S\) to the category of path types precisely to avoid that circle.
\end{dcomment}


L'idée naturelle, c'est de prendre, pour $s_0 \in \mathrm{Ob}\,S$ fixé,
\[ S' = \Ch(S; s_0, -) \qquad \text{ou} \qquad S' = \Ch_\infty(S; s_0, -), \]
et la flèche
\[ \beta : c \longmapsto \text{but } c : S' \longrightarrow S. \]
On sait que $S'$ est $W_\infty$-aphérique et que $\beta$ est $W_\infty$-fibrante (et même parfaite). Il faut voir que $S'$ est aussi $W$-asphérique, et $\beta$ $W$-fibrante. Il faut donc revoir la démonstration dans le cas de $W_\infty$, et prouver qu'elle se transforme à $W$, % [sic] aphérique -> asphérique, lettre s omise
\msp{page 8}
moyennant les axiomes qu'il faut sur $W$. J'ai envie de regarder d'abord le choix $S' = \Ch(S; s_0, -)$, sans passage à une limite inductive qui m'obligerait à mettre un axiome de limites sur $W$. Les propriétés homotopiques délicates et essentielles de $\Ch(S)$ sont établies dans VII, p. 157-169, savoir que dans le diagramme
\[ S \xrightarrow{i_S} \Ch(S) \xrightarrow{p_S} S \times S, \]
$i_S \in W_\infty$ et $p_S$ est une $W_\infty$-fibration (en fait, \textit{propre} et \textit{lisse} pour $W_\infty$) (\footnote[12]{Dans l'énoncé des ‘propriétés essentielles’ de $\Ch(S)$, on peut travailler avec $W_\omega$, la plus petite partie de Fl(Cat) satisfaisant les axiomes W(1), W(2 bis), W(3) (cf. page 14), donc $W_\omega \subseteq W \cap W_\infty$.}). Le fait que $i_S \in W$ \texttt{[plutôt} $i_S \in W_\infty$\texttt{]} implique que les foncteurs source et but
\[ \begin{tikzcd} \Ch(S) \arrow[r, shift left, "\sigma"] \arrow[r, shift right, "\beta"'] & S \end{tikzcd} \qquad \sigma = \mathrm{pr}_1 \circ p_S,\ \beta = \mathrm{pr}_2 \circ p_S \]
sont dans $W_\infty$ (puisque $\sigma i_S = \beta i_S = \mathrm{id}_S$). Comme ce sont des Hot-fibrations, il en résulte que les fibres $\Ch(S; s, -)$ resp. $\Ch(S; -, s)$ sont Hot-asphériques. Mais même en admettant que $\sigma$, $\beta$ sont dans $W$ et sont des $W$-fibrations, je ne peux en conclure tout de suite que les fibres sont $W$-asphériques - il y aurait cercle vicieux! Il me faut donc une démonstration directe.

Considérons le foncteur
\[
\begin{tikzcd}[row sep=huge]
\Ch(S; s, -) \arrow[d, "p"'] \\
\underline{Ty} \quad ,
\end{tikzcd}
\]
% foot: Version 7 décembre 2002
\printedpage{8}
% head: A. Grothendieck : Dérivateurs, Ch. XII, Édition des Universités de Montpellier II et Paris VII
c'est un foncteur Cat-fibrant, donc $W$-lisse (je devrai développer la notion de $W$-lissité (\footnote[13]{$f : X \longrightarrow Y$ est $W$-lisse si les $X_y \longrightarrow y\backslash X$ sont $W$-asphériques.}) avec soin, j'admettrai pour le moment les propriétés formelles familières, et les justifierai par la suite). À cause de la $W$-lissité, pour prouver $p \in W$, il suffit de
\msp{page 9}
prouver que les fibres de $p$ sont $W$-asphériques (\footnote[14]{On utilise ici W(2 bis), pas $W(2)$.}). Or les fibres sont les $Z_\tau = \Ch_\tau(S; s, -)$, qui sont des catégories \textit{contractiles} (i.e. $Z_\tau \longrightarrow e$ est un homotopisme), donc $W$-asphériques par W(3) (cf. cor. 3, page 3).

\dgrblock{dgrb466}{Why finite path types avoid a limiting axiom}
\begin{dcomment}{exposition}{p9:prouver que les fibres de $p$ sont}{Why finite path types avoid a limiting axiom}
The projection to path types separates combinatorics from homotopy. For a fixed finite oriented interval \(\tau\), the category \(\Ch_\tau(S;s,-)\) consists of its diagrams in \(S\) whose initial endpoint is \(s\). Replace the last vertex by its predecessor and the last edge by the identity. This defines a functor on the diagram category; the last arrow of each diagram, with identities at the other vertices, defines a natural transformation to or from the identity functor, according to the orientation of the edge. Its image is the subcategory of diagrams with final edge an identity, canonically identified with the diagram category for the shortened interval. Repeat there and include back into the original category. Induction on the number of edges gives a finite zigzag from the identity to the constant-path functor at \(s\). This supplies the fixed-type contraction which Chapter~VII states and later leaves unwritten \citep[manuscript pp.~152, 160, 170]{GrothendieckDerivateursVII}.

The total category over \(\underline{Ty}\) is a Grothendieck fibration. Therefore it is \(W\)-smooth: every required category of lifts with prescribed target has a terminal cartesian lift. Aspherical strict fibres then give a \(W\)-equivalence to \(\underline{Ty}\), by the dual of absolute theorem A. The unit path type is terminal in \(\underline{Ty}\), so W(3) makes this base \(W\)-aspherical as well \citep[manuscript pp.~160--161]{GrothendieckDerivateursVII}. A claim that the fibres alone make the total category aspherical would omit this base hypothesis. Thus the path cover is built from categorical contractions and theorem A, before invoking any stability of filtered colimits.

The same economy explains the lifting calculation for
\(\Ch(S)\to S\times S\). For endpoint arrows \(u:a'\to a\), \(v:b'\to b\) and a path \(c:a\leadsto b\), form the concatenation \(\widetilde c=\operatorname{ch}(u)c\operatorname{ch}(v)^{\mathrm{op}}\). The category of lifts is compared with \(\Ch(S;a',b')/\widetilde c\), which has a terminal object. Chapter~VII constructs both comparison functors and links their composites to the identities by natural transformations \citep[manuscript pp.~161--166]{GrothendieckDerivateursVII}. These are actual categorical homotopies, so W(1)--W(3) put the comparisons in \(W\). Applying the construction to the opposite category supplies the other variance. This establishes properness and smoothness independently of W(4). Only afterwards does W(4) upgrade the family to preservation of arbitrary \(W\)-maps.

The lifting-category comparison is a proof by reference to Chapter~VII; the audit checked the cited construction and retains that dependency. The mapping-path factorization then follows from base-change and composition of perfect maps, with its first map a section of a universally aspherical source projection \citep[manuscript pp.~170--171]{GrothendieckDerivateursVII}. The distinction between finite zigzags and the infinite completion \(\Ch_\infty\), already developed in the commentary on Chapter~VII, manuscript pp.~157--170, explains why the present proof need not assume W(6).
\end{dcomment}


D'autre part, si on admet que $\Ch(S) \xrightarrow{p_S} S{\times}S$ est $W$-lisse (et $W$-propre), donc $\Ch(S) \xrightarrow{\sigma_S} S$ aussi, étant à fibres $W$-asphériques comme on vient de voir, il s'ensuit que $\sigma_S \in W$, d'où $S \xrightarrow{i_S} \Ch(S)$ aussi dans $W$ (puisque $\sigma_S \circ i_S = \mathrm{id}_S$). Donc il reste à prouver seulement que
\[ \Ch(S) \xrightarrow{p_S} S \times S \tag*{$(*)$} \]
est $W$-propre et $W$-lisse. Il en résultera que
\[ \Ch(S; s_0, -) \xrightarrow{\beta} S \]
est lui aussi $W$-propre et $W$-lisse, car déduit de $p_S$ par le changement de base
\[ S \longrightarrow S \times S \qquad s \longmapsto (s_0, s). \]

Il faut ensuite voir que $W$-lisse et $W$-propre implique la propriété de $W$-fibration \textit{sous sa forme la plus forte} a) (alors que je prévois qu'a priori on ne l'aura que sous les formes faibles a$_1$, a$_2$ - pas même sous la forme a$'$). C'est ici donc que l'axiome W(4) des $W$-fibrations sera essentiel (\footnote[15]{sous la forme affaiblie W(4a$''$) (p. 13).})!

Voici l'argument pour la lissité de $(*)$, dans loc. cit. p. 161 ff. Il faut prouver que pour $c \in \mathrm{Ob}\,\Ch(S)$, et une flèche $u \in \mathrm{Fl}(S \times S)$ \texttt{[plutôt} $(u,v) \in$ \texttt{Fl(}$S \times S$\texttt{)]}, une certaine catégorie de ‘relèvements de $u$ de source $c$ ’ \texttt{[plutôt ‘relèvements de} $(u,v)$ \texttt{de but} $c$ \texttt{’]}, $\mathcal{C} = \mathcal{C}(c,u,v)$, est $W$-asphérique (\footnote[16]{En fait, on prouve que $\mathcal{C}(c,u,v)$ est \textit{contractile}, donc elle est $W_\omega$-asphérique.}). Pour ceci, on la compare à une catégorie $\mathcal{C}'$ à objet final, donc contractile et a fortiori $W$-asphérique, à l'aide de deux foncteurs
\[ \begin{tikzcd} \mathcal{C} \arrow[r, shift left, "\varphi"] & \mathcal{C}', \arrow[l, shift left, "\psi"] \end{tikzcd} \]
et des homotopies $\varphi\psi \sim \mathrm{id}_{\mathcal{C}'}$ (sans mérite, $\mathcal{C}'$ étant contractile)
\msp{page 10}
et $\psi\varphi \sim \mathrm{Id}_{\mathcal{C}}$ (on trouve même une flèche $\psi\varphi \longrightarrow \mathrm{id}_{\mathcal{C}}$). Ainsi $\varphi$ est un homotopisme, donc une $W$-équivalence (cor. 3, p. 3), donc $\mathcal{C}$ est $W$-asphérique, $\mathcal{C}'$ l'étant. La propreté de $p_S$ équivaut à la lissité de $p_{(S^\circ)}$, elle est donc prouvée du même coup.
% foot: Version 7 décembre 2002

% ---- from Der-ChXIIa-012_p011-015.tex
\printedpage{9}
% head: A. Grothendieck : Dérivateurs, Ch. XII, Édition des Universités de Montpellier II et Paris VII

Ainsi, le lemme de la page 7 est prouvé, moyennant l'axiome W(4) qui assure qu'une flèche qui est à la fois $W$-lisse et $W$-propre, satisfait à la condition très forte a) des $W$-fibrations.

\textbf{NB} En fait, il suffirait, au lieu de l'axiome W(4) sous sa forme forte, de supposer qu'une $f : X \longrightarrow Y$ qui est à la fois `homotopiquement lisse' et `homotopiquement propre', i.e. les fibres des $X/x \longrightarrow Y/y$ et des $x\backslash X \longrightarrow y\backslash Y$ sont contractiles (donc $W$-asphériques pour tout $W$ satisfaisant W(1,2,3)), soit $W$-fibration forte. Énonçons quand même ce que nous avons prouvé.

\textbf{Proposition 4.} \textit{Si $W$ satisfait} W(1,2,3) (\footnote[17]{Pour W(2), il le suffit sous la forme W(2 bis).}) (nous le supposons toujours par la suite) \textit{plus} W(4) (sous une forme un peu affaiblie suffisante, cf. remarque précédente), \textit{on a ceci: Soit $f : X \longrightarrow Y$ une flèche dans} Cat$/S$\textit{, avec $X$, $Y$ des $W$-fibrés sur $S$. Si $f \in W$, alors $f_s : X_s \longrightarrow Y_s$ est dans $W$ pour tout $s \in S$.}

Je me demande si j'en sais assez pour pouvoir prouver
\[ W \subseteq W_\infty . \]
Il faut quand même éliminer des cas canularesques, comme $W = \mathrm{Fl}(\mathrm{Cat})$, donc mettre l'axiome

\begin{itemize}[leftmargin=4em, labelwidth=3em, labelsep=0.6em]
\item[W(5)] (Axiome de connexité (\footnote[18]{Voir à ce sujet le corollaire page 22.}).) Si $f : X \longrightarrow Y$ est dans $W$, $\pi_0(f) : \pi_0(X) \longrightarrow \pi_0(Y)$ \texttt{[est]} bijectif.
\end{itemize}

\msp{page 11}Je factorise alors $f$ en (\footnote[19]{À nouveau, dans ces arguments on a intérêt à travailler avec $W_\omega$.})
\[ X \xrightarrow{i} \bar{X} \xrightarrow{\bar{f}} Y, \]
avec $\bar{f}$ une $W_\infty$-fibration et une $W$-fibration (\footnote[20]{en fait, $\bar{f}$ est à la fois $W_\infty$-parfait (i.e. $W_\infty$-lisse et propre) et $W$-parfait (en fait, $W_\omega$-parfait).}), $i \in W \cap W_\infty$ - cela marche finallement, % [sic] finallement -> finalement, spelling
par la construction de VII, p. 170-171 (de \textsc{Cartan-Serre}), à partir des propiétés de % [sic] propiétés -> propriétés, missing letter
$Y \longrightarrow \Ch(Y) \longrightarrow Y \times Y$ sans plus. Comme $\bar{f} i = f \in W$ et $i \in W$, on en déduit $\bar{f} \in W$. Comme $i \in W_\infty$, on en conclut $f \in W_\infty \iff \bar{f} \in W_\infty$. Donc pour prouver que $f$ est dans $W_\infty$, on est ramené au cas où $f$ est une $W$-fibration et une $W_\infty$-fibration. Mais alors la proposition 4 nous montre que les fibres de $f$ sont $W$-asphériques (\footnote[21]{on n'a utilisé ici que cette propriété sur les fibres, comme consequence de W(4). (On utilise W(4) sous sa forme la plus faible W(4b$''$), cf. pages 13, 14).}), et il suffit de voir qu'elles sont $W_\infty$-asphériques (proposition 3 (\footnote[22]{on n'utilise la proposition 3 que dans le cas $Y = S$, pour lequel W(2 bis) suffit.}) appliquée à $W_\infty$). Ainsi on est ramené au cas où $Y = e$, i.e. à prouver que si $X$ dans Cat est $W$-asphérique, il est $W_\infty$-asphérique.
% [sic] (footnote 21) consequence -> conséquence, missing accent

Soit $x_0 \in \mathrm{Ob}\, X$, considérons comme tantôt
\[ E = \Ch(X, x_0, -) \xrightarrow{p} X, \]
% foot: Version 7 décembre 2002
\printedpage{10}
% head: A. Grothendieck : Dérivateurs, Ch. XII, Édition des Universités de Montpellier II et Paris VII
on sait que $E$ est $W \cap W_\infty$-asphérique, et que c'est une fibration pour $W$ et pour $W_\infty$. Comme $X$ est $W$-asphérique (comme $E$), $p \in W$, donc à nouveau par la proposition 4 (\footnote[23]{en fait, la condition W(4b$''$) sur $W$.}), les fibres de $p$ sont $W$-asphériques, en particulier (par W(5)) 0-connexes. On a donc prouvé le

\msp{page 12}\textbf{Lemme.} (\footnote[24]{Utilise W(1, 2 bis, 3, 4b$''$) sur $W$.}) \textit{Soit $(X, x_0)$ dans} Cat$_\bullet$\textit{. Si $X$ est $W$-asphérique, alors $\Omega(X, x_0) \overset{\text{déf}}{=} \Ch(X; x_0, x_0)$ l'est aussi} (moyennant les axiomes W(1,2,3,4) sur $W$)\textit{. Donc (par itération) les $\Omega^n(X, x_0)$ sont $W$-asphériques} ($\forall\, n \in \mathbf{N}$).

\textbf{Corollaire.} \textit{Si de plus $W$ satisfait l'axiome} W(5) \textit{de connexité, alors $X$ est $W_\infty$-asphérique.}

En effet, `on sait' que cela signifie que tous les $\Omega^n(X, x_0)$ sont 0-connexes. Or l'axiome de connexité W(5) assure que $W$-asphérique implique 0-connexe.

\dgrblock{dgrb467}{The loop comparison that closes Theorem 1}
\begin{dcomment}{exposition}{p12:que cela signifie que tous les}{The loop comparison that closes Theorem 1}
Grothendieck has isolated the decisive argument: a \(W\)-aspherical category remains \(W\)-aspherical after forming its based loop category. To obtain the theorem, these loops must detect the ordinary homotopy groups. This is the comparison left open in the running text, and it has a short proof once the ordinary fibration properties already used are available.

Let \(E=\Ch(X;x,-)\) and \(F=\Ch(X;x,x)\). The endpoint map \(E\to X\) is a \(W_\infty\)-fibration. Its strict fibre over \(x\) maps by a weak equivalence to the corresponding comma category, and the comma categories vary by weak equivalences along arrows of \(X\). Quillen's theorem B therefore identifies \(BF\) with the homotopy fibre of \(BE\to BX\) \citep[Theorem~B and its corollary, pp.~97--98]{Quillen1973}. Since \(BE\) is weakly contractible, this gives a pointed equivalence
\[
 B\Ch(X;x,x)\simeq\Omega_x BX.
\]
At the next step choose the constant loop as basepoint. Repeating the same argument gives
\[
 B\Omega^r(X,x)\simeq\Omega^r_x BX,\qquad
 \pi_0\Omega^r(X,x)\cong\pi_r(BX,x)
 \quad(r\geq1),
\]
with the right side regarded as its underlying pointed set where only connectivity is used.

Now W(4b\(^{\prime\prime}\)) applies to the perfect path map \(E\to X\): both total space and base are \(W\)-aspherical, so their map is in \(W\), and its fibre is \(W\)-aspherical. Induction gives this for every \(\Omega^r(X,x)\). W(5) makes each of them nonempty and connected, hence all positive homotopy groups of \(BX\) vanish; W(5) applied to \(X\to e\) gives connectivity in degree zero. Thus \(X\to e\in W_\infty\).

Finally, for \(u:A\to B\in W\), the canonical perfect replacement \(A\to\bar A\to B\) has first map in both \(W\) and \(W_\infty\). The second map has \(W\)-aspherical fibres by W(4b\(^{\prime\prime}\)); the preceding argument makes those fibres weakly contractible. Its ordinary homotopy-fibre sequence makes it a weak equivalence, and two-out-of-three gives \(u\in W_\infty\). The conclusion is proved; the source presents it with the loop comparison still deferred.
\end{dcomment}


Mais le `on sait' est un tout petit peu léger (\footnote[25]{C'est prouvé fort simplement dans le corollaire du `theorème 1', page 190.}). Pour le justifier, il faudrait, ou bien % [sic] (footnote 25) theorème -> théorème, missing accent
que je fasse le lien avec la théorie 1/2-simpliciale classique, et prouver que les $\Omega^n$ que je construis sont `les bons', i.e. que leur nerf redonne bien (à isomorphisme près dans Hot) les $\Omega^n$ classiques (??). En particulier, que le $\pi_0(\Omega^n(X, x_0))$ est bien le $\pi_n(X, x_0)$ `de tout le monde'. Ou bien que je prouve directement que la 0-connexité des $\Omega^n(X, x_0)$ équivaut à la $W_\infty$-asphéricité de $X$, définie cohomologiquement à la \textsc{Artin-Mazur}. C'est donc un travail qu'il me faudra faire absolument. Mais ce n'est plus qu'un simple travail de routine, et je peux énoncer dès maintenant le

\textbf{Théorème 1.} (\footnote[26]{Il suffit même de W(4) sous une forme affaiblie, avec des foncteurs `fibrants' dans un sens `absolu'. Cf. le NB p. 10 - et surtout, cf. pages 13, 14.},\footnote[27]{NB Au lieu de W(5), il suffit de supposer $W \neq \mathrm{Fl}(\mathrm{Cat})$, cf. th. 2, p. 20.},\footnote[28]{Les hypothèses utiles sur $W$ sont finalement W(1), W(2 bis), W(3), W(4b$''$), et $W \neq \mathrm{Fl}(\mathrm{Cat})$.}) \textit{Soit $W \subseteq$} Fl(Cat) \textit{satisfaisant les conditions} W(1,2,3,4,5) (axiome de saturation, de localisation, de l'objet final, de fibration, de connexité)\textit{. Alors}
\[ W \subseteq W_\infty . \]

\textbf{Corollaire.} \textit{On peut caractériser $W_\infty$ comme la plus grande partie de} Fl(Cat) \textit{satisfaisant aux}
\msp{page 13}\textit{conditions} W(1) \textit{à} W(5) \textit{précédentes.}

Bien sûr, on aimerait un axiome 6 plus intéressant, qui joint aux précédents assure que $W = W_\infty$ (\footnote[29]{Cf. axiome W(6) page suivante.}). Ou est-ce que ces propriétés assurent déjà que $W = W_\infty$? Pour le voir, il faudrait peut-être que j'en sache plus sur la structure des objets $W_\infty$-asphériques? Si c'était faux, je songe à un axiome de `descente homotopique', qui impliquerait qu'un objet $W_\infty$-asphérique est aussi $W$-asphérique.

\dgrblock{dgrb468}{The conjectural minimum is now known}
\begin{dcomment}{status}{p13:Bien sûr, on aimerait un axiome 6 plus intéressant}{The conjectural minimum is now known}
The question on this page has a positive answer without the proposed additional axiom W(6). Cisinski proved that \(W_\infty\) is the smallest fundamental localizer \citep[Theorem~2.2.11]{Cisinski2004Minimal}. The refinement to the absolute theorem-A axioms identifies it also as the smallest \emph{weak} fundamental localizer \citep[Theorem~6.1.18]{Cisinski2006}. The source's W(1c) implies the weak-saturation condition in that theorem. Hence W(1), W(2 bis), and W(3) already give \(W_\infty\subseteq W\).

Combining this inclusion with Theorem~1 gives \(W=W_\infty\) under W(4b\(^{\prime\prime}\)) and W(5). The nondegeneracy argument at manuscript pp.~20--22 replaces W(5) by \(W\ne\Fl(\Cat)\). Accordingly, the conjecture at manuscript p.~15 is no longer open, and adding closure under retracts is unnecessary.

There is a local notational change to observe. On manuscript p.~13, \(W_\omega\) is introduced as the intersection of classes satisfying W(1)--W(5). On the next page it denotes the minimum satisfying only the basic localizer axioms, preferably W(2 bis). These are different definitions during the exploration. The later theorem makes both classes equal to \(W_\infty\), but that equality must not be used to conceal their different roles in the original argument.

The proof strategy here also reappears in Cisinski's characterization by a theorem-B property: a fundamental localizer preserving components and making every locally constant equivalence aspherical must be \(W_\infty\) \citep[Theorem~2.3.6]{Cisinski2004Minimal}. Its loop-space argument is the modern counterpart of the preceding manuscript proof.
\end{dcomment}


% foot: Version 7 décembre 2002
\printedpage{11}
% head: A. Grothendieck : Dérivateurs, Ch. XII, Édition des Universités de Montpellier II et Paris VII

Je note dès maintenant que si des $W_i$ ($i \in I$) satisfont aux axiomes W(1) à W(5), il en est de même de $\bigcap_i W_i$ (\footnote[30]{C'est vrai pour tout ensemble d'axiomes dans la nature d'axiomes de stabilité.}). Donc il y a un plus petit ensemble de flèches $W_\omega$ satisfaisant à ces conditions. S'il n'est égal à $W_\infty$, i.e. si les axiomes W(1) à W(5) ne caractérisent déjà $W_\infty$, ce serait un localiseur vraiment très intéressant.

Je me rends compte que je n'ai utilisé qu'une forme assez faible de l'axiome W(4) - essentiellement, la validité de la proposition 4. Je vais résumer différentes variantes de l'axiome de fibration.

\textbf{W(4)}

\begin{itemize}[leftmargin=4em, labelwidth=2.5em, labelsep=0.6em, align=right]
\item[a$_1$)] Si $f : X \longrightarrow Y$ dans Cat est une $W$-fibration gauche, c'est une $W$-fibration ultraforte.

\item[a$_2$)] Si $f : X \longrightarrow Y$ est une $W$-fibration droite, c'est une $W$-fibration ultraforte.

\item[a$'$)] Si $f : X \longrightarrow Y$ est $W$-parfait, i.e. à la fois $W$-propre et $W$-lisse (\footnote[31]{ce qui implique que c'est une $W$-fibration gauche \textit{et} droite.}), alors $f$ est une $W$-fibration ultraforte.

\item[a$''$)] Si $f : X \longrightarrow Y$ est à la fois homotopiquement propre et homotopiquement lisse (cf. NB page 10), $f$ est une $W$-fibration ultraforte. (Il vaut mieux dire : Si $f$ est $W_\omega$-parfait, c'est une $W$-fibration ultraforte.)

\item[b)] Soit $f : X \longrightarrow Y$ une flèche dans Cat$/S$, avec $X$ et $Y$ des $W$-fibrés sur $S$. Si $f \in W$, alors $\forall\, s \in S$, $f_s : X_s \longrightarrow Y_s$ est \texttt{[dans]} $W$.

\item[b$'$)] Soit $f : X \longrightarrow Y$ une flèche de Cat$/S$ qui soit $W$-parfaite et $W_\infty$-parfaite, alors $f \in W$ implique que les fibres de $f$ sont $W$-asphériques. (Il vaut mieux dire : Si $f$ est $W_\omega$-parfait, il est une $W$-fibration ultraforte.) \texttt{[Plutôt : Soit $f : X \longrightarrow Y$ une flèche dans \textrm{Cat}$/S$ telle que $X \longrightarrow S$ et $Y \longrightarrow S$ soient $W_\omega$-parfaites. Si $f \in W$, alors $\forall\, s \in S$, $f_s : X_s \longrightarrow Y_s$ est dans $W$.]}

\item[b$''$)] \msp{page 14}Soit $W_\omega \subseteq W \cap W_\infty$ la plus petite partie de Fl(Cat) satisfaisant les conditions W(1,2,3) (\footnote[32]{Il vaut mieux prendre W(2 bis) au lieu de W(2), pour trouver $W_\omega$ la plus petite possible!}) (axiomes de saturation, de localisation, de l'objet final). Soit $f : X \longrightarrow Y$ une flèche dans Cat qui soit $W_\omega$-parfaite. Si $f \in W$, alors les fibres de $f$ sont $W_\omega$-asphériques \texttt{[plutôt $W$-asphériques]}.
\end{itemize}

Entre ces propriétés de $W$, on a les implications tautologiques (sauf a$''$) $\Longrightarrow$ b), qui est essentiellement la proposition 4)

% foot: Version 7 décembre 2002
\printedpage{12}
% head: A. Grothendieck : Dérivateurs, Ch. XII, Édition des Universités de Montpellier II et Paris VII

\[
\begin{tikzcd}[row sep=huge, column sep=0pt]
\mathrm{a} \arrow[d, Rightarrow] & {}= (\mathrm{a}_1 \Longleftrightarrow \mathrm{a}_2) & \hspace{2.5cm} \\
\mathrm{a}' \arrow[d, Rightarrow] & & \smash{\left.\rule[-1.8cm]{0pt}{3.6cm}\right\}\ \text{\begin{tabular}{l}En fait,\\ toutes ces\\ conditions sont\\ équivalentes\end{tabular}}} \\
\mathrm{a}'' \arrow[d, Rightarrow] & & \\
\mathrm{b} \arrow[d, Rightarrow] & & \\
\mathrm{b}' \arrow[d, Rightarrow] & & \\
\mathrm{b}'' & &
\end{tikzcd}
\]

(\footnote[33]{Je reviens sur les diverses variantes de W(4) plus bas, p. 53-56.}).

Eh bien, dans le lemme page 12, il suffisait de supposer que $W$ satisfait à la version la plus faible b$''$) de l'axiome des fibrations (sept versions!). En effet, les fibrations $\bar{X} \longrightarrow Y$ et $E \longrightarrow X$ utilisées dans la démonstration du lemme sont $W_\omega$-parfaites, et il suffisait de savoir que si elles sont dans $W$, alors les fibres sont $W$-asphériques.

\dgrblock{dgrb469}{What the weakest fibre axiom actually needs}
\begin{dcomment}{exposition}{p14:Eh bien, dans le lemme page 12}{What the weakest fibre axiom actually needs}
The list of seven variants is an attempt to reduce the amount of fibration theory needed for the characterization. Its least demanding member, W(4b\(^{\prime\prime}\)), asks only this: a \(W\)-equivalence which is \(W_\omega\)-perfect has \(W\)-aspherical strict fibres. The conclusion is relative to \(W\), as the editors correct; it does not assert that the fibres are already \(W_\omega\)-aspherical.

For the proof of Theorem~1 this is enough because the maps to which the axiom is applied have a special origin. The replacement \(\bar A\to B\) and the based-path projection \(\Ch(X;x,-)\to X\) are \(W_\omega\)-perfect before any fibration axiom for \(W\) is assumed. Two-out-of-three puts these maps in \(W\), and W(4b\(^{\prime\prime}\)) then supplies exactly the asphericity required. Arbitrary \(W\)-fibrations need not be considered.

The distinction between preservation and detection also explains the middle variants. W(4a\(^{\prime\prime}\)) requires base change along an absolute perfect map to preserve \(W\). W(4b\(^{\prime}\)) concerns a map \emph{between} two absolute perfect families over the same base. W(4b\(^{\prime\prime}\)) concerns one perfect map and compares it to the identity family on its target. The editorial replacement of the entire statement of b\(^{\prime}\) restores this intended progression.

The brace printed on this page spans the three a-conditions; it should not be read as a proof of every converse in the seven-term chain. Section~IV returns to those converses and to the saturation needed for the localization arguments. In the modern setting, minimality and the nondegeneracy criterion below identify the two possibilities as \(W_\infty\) and the class of all maps, so these different rigidity axioms ultimately select the same possibilities. That later conclusion is stronger than the argument supplied by the diagram itself.
\end{dcomment}


Mais je ne pense pas que même la version la plus forte W(4a) de l'axiome des fibrations, jointe aux trois axiomes précédents et à W(5) (connexité), implique $W = W_\infty$. Il faudrait tout au moins y joindre l'axiome

\begin{itemize}[leftmargin=4em, labelwidth=3em, labelsep=0.6em]
\item[\textbf{W(6)}] (Axiome des limites filtrantes.) Soit $X = \varinjlim\limits_{i} X_i$ une limite inductive filtrante dans Cat. Si les $X_i$ sont $W$-asphériques, $X$ aussi.
\end{itemize}

Ceci devrait être équivalent à la propriété, en apparence plus forte, en termes de flèches de $W$ :

\textbf{Proposition 5.} \textit{Supposons qu'on ait} W(6) (en plus de W(1) à W(4) (\footnote[34]{W(4) sous la forme la plus faible W(4b$''$) .})),
\msp{page 15}\textit{et soit}
\[ (f_i)_{i\in I} : (X_i)_{i\in I} \longrightarrow (Y_i)_{i\in I} \]
% foot: Version 7 décembre 2002
\printedpage{13}
% head: A. Grothendieck : Dérivateurs, Ch. XII, Édition des Universités de Montpellier II et Paris VII
\textit{un homomorphisme de systèmes inductifs dans} Cat\textit{, indexé par un ensemble ordonné filtrant, soient $X = \varinjlim X_i$, $Y = \varinjlim Y_i$, $f = \varinjlim f_i$, $f : X \longrightarrow Y$. Si $f_i$ est dans $W$ pour tout $i$, alors $f \in W$.}

On factorise chaque $f_i$ en
\[ X_i \xrightarrow{\alpha_i} \bar{X}_i \xrightarrow{\bar{f}_i} Y_i, \qquad \alpha_i \in W_\omega,\ \bar{f}_i \in \mathrm{Parf}\, W_\omega \]
(\footnote[35]{$W_\omega$ introduit dans la condition W(4b$''$), page 14.}) suivant le système breveté. La construction étant fonctorielle, les $\bar{X}_i$ forment eux aussi un système inductif. De plus, comme le foncteur $S \longrightarrow \Ch(S)$ qui donne naissance à la factorisation canonique des morphismes, commute aux $\varinjlim$ filtrantes (immédiat), il s'ensuit que, si on pose
\[ \bar{X} = \varinjlim_{i} \bar{X}_i \]
et qu'on regarde les flèches
\[ X \xrightarrow{\alpha} \bar{X} \xrightarrow{\bar{f}} Y, \qquad \alpha = \varinjlim \alpha_i,\ \bar{f} = \varinjlim \bar{f}_i, \]
avec $\bar{f}\alpha = \bar{f}$, cette factorisation de $f$ est bien sa factorisation canonique. Donc $\alpha \in W$, % [sic] $\bar{f}\alpha = \bar{f}$ -> $\bar{f}\alpha = f$, the composite should be f
$\bar{f} \in \mathrm{Parf}\, W_\omega$. Ceci vu, en vertu de prop. 3, $f \in W$, i.e. $\bar{f} \in W$, est vrai si les fibres de $\bar{f}$ sont $W$-asphériques. Mais les fibres de $\bar{f}$ s'interprètent comme limites inductives filtrantes de fibres des $\bar{f}_i$, lesquelles sont $W$-asphériques par la condition W(4b$''$) des fibrations. Donc l'axiome W(6) donne ce qu'il faut.

\dgrblock{dgrb470}{Filtered fibres require a moving basepoint}
\begin{dcomment}{exposition}{p15:La construction étant fonctorielle}{Filtered fibres require a moving basepoint}
Proposition~5 turns an axiom about aspherical objects into an axiom about maps by making the fibres of those maps explicit. Two finite-presentability assertions are implicit in the sketch, and both can be checked directly.

First, a path of a fixed finite type is a functor from a finite category, and a morphism between two such paths is specified by finitely many arrows and equations. If \(S=\varinjlim_i S_i\) is a filtered colimit, any path or morphism of paths in \(S\) is represented at one stage, and an equality between representatives holds at some later common stage. Thus \(\Ch(-)\), including its source and target maps, commutes with filtered colimits. Since filtered colimits of sets commute with finite limits, the finite pullbacks defining the mapping-path replacement commute with them as well.

Second, a fibre over \(y\in Y=\varinjlim_iY_i\) is not the colimit of fibres over an unspecified fixed \(y\). Choose a representative \(y_i\in Y_i\), and let \(y_j\) be its image for \(j\geq i\). There is a canonical isomorphism
\[
 (\bar X)_y\ \cong\
 \varinjlim_{j\geq i}(\bar X_j)_{y_j}.
\]
For objects, a representative \(x_j\) mapping to \(y\) has image equal to \(y_k\) at a sufficiently late stage \(k\), so it becomes an object of the indicated fibre. For arrows, the condition that the image be the identity of \(y\) is likewise witnessed at a later stage. The same finite-stage argument proves injectivity and compatibility with composition. Passing to \(j\geq i\) is cofinal and remains filtered.

Each map \(\bar f_j\) belongs to \(W\), since \(f_j\) and its replacement inclusion do. W(4b\(^{\prime\prime}\)) makes its fibres \(W\)-aspherical, W(6) makes their displayed colimit \(W\)-aspherical, and theorem A for a proper or smooth map gives \(\bar f\in W\). Finally \(\alpha\in W_\omega\subseteq W\), so \(f=\bar f\alpha\in W\). The printed equation \(\bar f\alpha=\bar f\) is a slip retained in the source.

For fundamental localizers in the modern sense, filtered-colimit stability is already a theorem, with no accessibility assumption \citep[Corollary~4.2.22]{Cisinski2006}. The calculation here remains useful because it shows exactly how a fibre argument reduces that stability to the object form.
\end{dcomment}


\textbf{Conjecture} (hardie? sotte?). \textit{Si $W \subseteq$} Fl(Cat) \textit{satisfait les axiomes} W(1) \textit{à} W(6) \textit{(avec} W(4) \textit{sous sa forme la plus faible} W(4b$''$) (\footnote[36]{et peut-être W(2) sous la forme plus faible W(2 bis).})\textit{), alors $W = W_\infty$.}

\msp{page 16}Mais il est bien possible qu'il faille renforcer un peu W(1) (saturation), en exigeant que tout rétracte d'un $f \in W$ soit dans $W$.

\dgrblock{dgrb471}{The proposed retract axiom is unnecessary}
\begin{dcomment}{status}{p16:Mais il est bien possible qu'il faille renforcer}{The proposed retract axiom is unnecessary}
The possible need for an extra retract axiom is resolved by the later theory. Every fundamental localizer with the relative theorem-A axiom W(2) is strongly saturated and consequently closed under retracts \citep[Proposition~4.2.4]{Cisinski2006}. In the weaker setting using W(2 bis), the specific equality conjectured here already follows from the weak minimal-localizer theorem and the fibre argument, so no independent retract hypothesis is required for that conclusion either. This does not assert strong saturation for every weak fundamental localizer before the fibre axiom is imposed.

This does not retroactively make every appeal to an isomorphism in \(\Hot_W\) valid under two-out-of-three alone. The elementary proofs above deliberately establish membership in \(W\) directly. The later strong-saturation theorem explains when such arguments may be shortened; it should be cited at the point where that stronger conclusion is used.
\end{dcomment}


\textbf{Remarque.} Considérons, pour $n \in \mathbf{N}$, l'ensemble $W_n \subseteq$ Fl(Cat) des `$n$-équivalences d'homotopie' dans Cat :
\[ f \in W_n \iff \pi_0(f) \text{ bijectif, et } \pi_i(f, x) : \pi_i(X, x) \longrightarrow \pi_i(Y, f(x)) \text{ bijectif pour } 0 < i \leq n. \]

Je croyais que ça satisfaisait à tout ce qu'on peut demander à un localiseur, sans avoir pris la peine de l'écrire avec soin. Ma ça ne peut satisfaire à W(1) jusqu'à W(4b$''$), puisqu'on % [sic] Ma -> Mais, word truncated
n'a pas $W_n \subseteq W_\infty$!

\dgrblock{dgrb472}{Truncated equivalences really are localizers}
\begin{dcomment}{exposition}{p16:n'a pas $W_n \subseteq W_\infty$!}{Truncated equivalences really are localizers}
The class \(W_n\) is not a defective homotopy theory. It is the ordinary theory observed through the first \(n\) homotopy groups. More precisely,
\[
  u\in W_n\quad\Longleftrightarrow\quad
  \tau_{\leq n}Bu\text{ is an equivalence},
\]
where \(\tau_{\leq n}\) is Postnikov truncation. This functor is left adjoint to the inclusion of \(n\)-types \citep[Proposition~5.5.6.18]{Lurie2009HTT}. Its universal property gives a direct modern proof of the assertions the source recalls without repeating.

Identities, two-out-of-three, and the two-cycle condition hold because they hold for equivalences of \(n\)-types. Strong saturation follows too: the functor \(\tau_{\leq n}B\), regarded as taking values in the homotopy category of \(n\)-types, factors through \(\Cat[W_n^{-1}]\). Hence any map inverted there induces an equivalence of \(n\)-types and therefore was in \(W_n\) already. Categories with terminal objects have contractible nerves, giving W(3).

For the relative axiom W(2), let \(X\to S\) be any functor. The simplicial bar construction on the covariant diagram \(s\mapsto N(X/s)\) has bisimplices consisting of
\[
 x_0\to\cdots\to x_q,\qquad
 p(x_q)\to s_0\to\cdots\to s_r.
\]
For a fixed \(q\)-simplex of \(NX\), its \(r\)-direction is the nerve of \(p(x_q)\backslash S\), which has an initial object. Realizing this direction therefore gives a weak equivalence of the diagonal bar construction with \(NX\). In other words,
\[
 BX\simeq\operatorname*{hocolim}_{s\in S}B(X/s).
\]
This is a homotopy colimit of spaces, retaining the arrows of \(S\), rather than a colimit of their homotopy categories. If \(X/s\to Y/s\) is in \(W_n\) for every \(s\), apply \(\tau_{\leq n}\) to the displayed colimits. As a left adjoint, truncation preserves them as colimits in the category of \(n\)-types. The resulting objectwise equivalence induces an equivalence of colimits, proving \(X\to Y\in W_n\).

Finally, the nerve commutes with filtered colimits because each finite chain of arrows and each relation is witnessed at a finite stage. A Kan replacement such as \(\operatorname{Ex}^{\infty}\) also commutes with filtered colimits: each simplex of \(\operatorname{Ex}\) is represented by a map from a finite subdivision, and its sequential iteration commutes with colimits. Homotopy groups of Kan complexes are computed by finite simplicial spheres and finite homotopies. They consequently commute with filtered colimits, with basepoints represented at sufficiently late stages. This proves W(6) for \(W_n\). What fails is the fibre axiom, as the next computation shows.
\end{dcomment}


W(1) est trivial sur la définition (et on a sans larmes même la \textit{saturation forte} de $W_i$, % [sic] $W_i$ -> $W_n$ (also in the following lines), the remark concerns W_n
i.e. si $f$ dans Cat est \texttt{[un]} isomorphisme dans $\mathrm{Hot}_{W_i} = (\mathrm{Cat})W_i^{-1}$, alors $f \in W_i$). Je crois me rappeler avoir vérifié W(2), par la même méthode que pour $W_\infty$, en utilisant un critère cohomologique à la \textsc{Artin-Mazur}, pour exprimer $f \in W_i$, et en utilisant les suites spectrales de Leray pour $X \longrightarrow S$ et $Y \longrightarrow S$. La condition W(3) est trivialement satisfaite puisque $W_n \supseteq W_\infty$ et qu'elle l'est pour $W_\infty$. Donc ça canule là où je m'y
% foot: Version 7 décembre 2002

% ---- from Der-ChXIIa-012_p016-020.tex
\printedpage{14}
% head: A. Grothendieck : Dérivateurs, Ch. XII, Édition des Universités de Montpellier II et Paris VII
attendais le moins : l'axiome des fibrations. Prenons la forme faible : Soit $f : X \longrightarrow Y$ une $W_\omega$-fibration et $f \in W_i$, les fibres $X_s$ sont elles $W_i$-asphériques (\footnote[37]{$\omega > \infty$, $W_\omega \subset W_\infty \subset W_n$})? On peut supposer $X$ et $Y$ 0-connexes, soit $x \in X_s$, la suite exacte d'homotopie donne
\[
\pi_{i+1}(X) \xrightarrow[\text{si } i \leq n-1]{\text{iso}} \pi_{i+1}(Y) \longrightarrow \pi_i(X_s) \longrightarrow \pi_i(X) \xrightarrow[\text{si } i \leq n]{\text{iso}} \pi_i(Y),
\]
donc
\[
\pi_i(X_s) = 1 \qquad \text{si } 0 \leq i \leq n-1,
\]
mais pour $i = n$ on trouve\msp{page 17}
\[
\pi_n(X_s) = \mathrm{Coker}\left(\pi_{n+1}(X) \longrightarrow \pi_{n+1}(Y)\right),
\]
il n'y a aucune raison que ce soit nul, comme ça devrait si $X_s \longrightarrow e$ était dans $W_n$.

\dgrblock{dgrb473}{The missing group occurs one degree higher}
\begin{dcomment}{examples}{p17:il n'y a aucune raison que ce soit nul}{The missing group occurs one degree higher}
Grothendieck's exact sequence identifies the precise loss of information. For a connected homotopy fibre \(F\) of an \(n\)-equivalence \(E\to B\), with \(n\geq1\), exactness gives
\[
 \pi_iF=0\quad(0\leq i<n),\qquad
 \pi_nF\cong
 \operatorname{coker}\bigl(\pi_{n+1}E\to\pi_{n+1}B\bigr).
\]
At \(n=1\), the cokernel is still meaningful: the connecting homomorphism from the abelian group \(\pi_2B\) surjects onto \(\pi_1F\), so the latter group is abelian in this situation. The long exact sequence and its low-dimensional interpretation are recalled in \citet[Theorem~4.41]{Hatcher2002}.

For an example entirely within \(\Cat\), take \(S\) to be the poset of nonempty proper faces of an \((n+2)\)-simplex. Then \(BS\) is the barycentric subdivision of its boundary, hence has the homotopy type of \(S^{n+1}\). Choose a vertex \(s\), and use
\[
 E=\Ch(S;s,-)\xrightarrow{\beta}S .
\]
The total space \(BE\) is contractible, and \(\beta\) is \(W_\omega\)-perfect by the construction recalled on manuscript pp.~8--11. Since the homotopy groups of the sphere vanish through degree \(n\), \(\beta\in W_n\). Its fibre \(F=\Ch(S;s,s)\), however, satisfies
\[
 BF\simeq\Omega S^{n+1},\qquad \pi_n BF\cong\mathbb Z.
\]
Thus \(F\to e\notin W_n\), contradicting W(4b\(^{\prime\prime}\)). The same square shows that \(\beta\) does not preserve all \(W_n\)-equivalences under base change: \(e\to S\) is in \(W_n\), whereas its pullback \(F\to E\) is not.

For \(n=0\), the sphere example remains valid, with \(\pi_0\Omega S^1\cong\mathbb Z\), but the displayed \emph{group cokernel} in the source requires correction. In general \(\pi_0F\) is the pointed coset set of the image of \(\pi_1E\) in \(\pi_1B\), not its quotient by the normal closure. Take \(G=S_3\), \(H=\langle(12)\rangle\), and the action groupoid \(E=(G/H)\mathbin{/\mkern-6mu/}G\). The projection \(E\to BG\) is both a Grothendieck fibration and an opfibration; its strict fibre is the discrete three-element set \(G/H\). Both \(E\simeq BH\) and \(BG\) are connected, so the projection is in \(W_0\). Yet the group cokernel of \(H\to G\) is trivial, since the conjugates of \((12)\) generate \(S_3\). The fibre has three components. This is a genuine degree-zero qualification, not a defect in the argument for positive \(n\).
\end{dcomment}


Il est vrai que parfois on renforce justement la condition pour $W_n$, disons qu'on définit $W'_n \subseteq W_n$ par
\[
f \in W'_n \iff \left\{
\begin{array}{ll}
\text{a)} & \pi_0(f) \text{ bijectif,} \\[1ex]
\text{b)} & \pi_i(f,x) : \pi_i(X,x) \longrightarrow \pi_i(Y,f(x)) \text{ bijectif pour } 0 < i \leq n \text{ et } x \in X, \\[1ex]
\text{c)} & \pi_{n+1}(f,x) : \pi_{n+1}(X,x) \longrightarrow \pi_{n+1}(Y,f(x)) \ \textit{surjectif,}
\end{array}
\right.
\]
c'est à dire aussi
\[
f \in W'_n \iff \left\{
\begin{array}{l}
\text{les `fibres homotopiques' } Z \text{ de } X \longrightarrow Y \text{ (i.e. les fibres de } \bar{X} \xrightarrow{\bar{f}} Y\text{)} \\
\text{sont } n\text{-connexes, i.e. 0-connexe et } \pi_i(Z_y,z) = 0 \text{ pour } 1 \leq i \leq n.
\end{array}
\right.
\]
Alors $X$ dans Cat sera $W'_n$-asphérique si et seulement s'il est $W_n$-asphérique, i.e. $n$-connexe. Donc par la forme 2 de la définition de $W'_n$, l'axiome W(4b$''$) est satisfait. Il faut croire que ça pèche ailleurs! Peut-être déjà la saturation.

\dgrblock{dgrb474}{Connectivity repairs fibres but loses cancellation}
\begin{dcomment}{examples}{p17:Alors $X$ dans Cat sera}{Connectivity repairs fibres but loses cancellation}
The strengthened class \(W'_n\) measures a connectedness condition on a map: its homotopy fibres are \(n\)-connected. The long exact sequence makes this equivalent to bijectivity on \(\pi_0,\ldots,\pi_n\) and surjectivity on \(\pi_{n+1}\). For the terminal map \(X\to e\), the last surjectivity condition is automatic, so \(W_n\) and \(W'_n\) have exactly the same aspherical objects.

The failure of two-out-of-three can be seen with the same finite category \(S\) modeling \(S^{n+1}\). Let \(a:e\to S\) select a vertex and let \(b:S\to e\) be the terminal functor. Then \(ba=\id_e\in W'_n\). Also \(b\in W'_n\): it is an isomorphism on homotopy groups through degree \(n\), and the map \(\mathbb Z=\pi_{n+1}S\to0\) is surjective. But \(a\notin W'_n\), because \(0\to\mathbb Z\) is not surjective. This works also for \(n=0\), using the circle.

Thus the problematic cancellation is exactly
\[
 b,\ ba\in W'_n\quad\nRightarrow\quad a\in W'_n.
\]
The source's informal argument about surjective homomorphisms is correct; the example supplies actual functors and verifies every homotopy group involved. A class of connected maps can be stable under many geometric operations without being the class inverted by a localization. The modern theory treats connectedness as one side of a factorization system; it is not required to satisfy the cancellation law for equivalences.
\end{dcomment}


\hspace*{3em}Si $gf$ et $g$ sont dans $W'_i$, il n'est pas forcément de même de $f$,

car si $\bar{g}\bar{f}$ (et donc $\bar{g}$) est surjectif, il n'en est par forcément de même pour $\bar{f}$! Donc $W'_n$ ne mérite pas le nom de `localiseur'. % [sic] par forcément -> pas forcément, misprint

Par contre, $W_n$ me paraît excellent, si on ne lui demande pas l'impossible - de satisfaire à un `axiome de fibration' qui semble bien (à peu de choses près) caractériser le seul $W_\infty$ parmi les `bons' localiseurs. En tous cas l'axiome de connexité W(5) est satisfait, et l'axiome des limites filtrantes W(6) aussi. Par ailleurs W(1) (et même la saturation forte), W(2), W(3) sont satisfaits par tout $W_{\mathbf{D}}$ associé à un dérivateur - ce qui confirme à nouveau leur importance. Les axiomes W(4) (l'axiome `draconien' des fibrations, $\pm$ spécial à $W_\infty$), W(5) (connéxité) et % [sic] connéxité -> connexité, misspelling (written connexité elsewhere)
\printedpage{15}
% head: A. Grothendieck : Dérivateurs, Ch. XII, Édition des Universités de Montpellier II et Paris VII
\msp{page 18}W(6) (limites) m'apparaissent comme de nature plus spéciale, je ne songerais à imposer aucun d'eux pour l'usage terminologique de `localiseur fondamental'.

\newpage
\printedpage{16}
% head: A. Grothendieck : Dérivateurs, Ch. XII, Édition des Universités de Montpellier II et Paris VII
\section*{II\quad Perplexités sur les axiomes et la terminologie.}

Dans I et la plupart du temps, en travaillant avec les dérivateurs, j'ai supposé \fbox{$\ W_{\mathbf{D}} \supseteq W_\infty\ $}, hypothèse qui m'a parue essentielle (à tort (\footnote[38]{à \textit{tort} il s'avère maintenant.}) ou à raison). On voit donc que si $W_{\mathbf{D}}$ satisfait à l'axiome de fibrations W(4), même sous sa forme la plus faible W(4b$''$), il est \textit{égal} à $W_\infty$. Cela montre bien que cet axiome doit être vu comme pratiquement une \textit{caractérisation} de $W_\infty$ (parmi les localiseurs fondamentaux qui contiennent $W_\infty$ -- et ces localisateurs me paraissent les plus importants, et peut-être les seuls vraiment utiles).

La théorie des \textbf{D}-fibrations dans I, p. 111-126, a été faite dans l'esprit qu'on vient de dire, en supposant $W_{\mathbf{D}} \supseteq W_\infty$. La notion de \textbf{D}-fibrations à laquelle je m'étais finalement arreté était par suite différente des notions envisagées ici dans un esprit différent (page 4). Elle revient à tester la condition a) (trop forte presque toujours, sauf pour $W = W_\infty$), caractérisant les $W$-fibrations ultrafortes, par des $Y'' \longrightarrow Y'$, où $Y''$, $Y'$ sont tous deux $W_\infty$-asphériques (ce qui est plus exigeant que $W$-asphérique). C'est seulement quand $Y''$ et $Y'$ sont $W_\infty$-asphériques (au lieu de $W$-asphériques), qu'on demande que $X'' \longrightarrow X'$ soit dans $W$. On trouve donc\msp{page 19} une notion moins exigeante que celle de $W$-fibration forte. Si on dénote ici cette condition par b), on trouve (sous réserve que $W_\infty \subseteq W$), avec les notations de la page 4 (\footnote[39]{Cf. la bonne notion de $W$-fibration dans un contexte où on ne suppose pas $W_\infty \subseteq W$, dans th. 3 et déf. 3, p. 51, 52.}) % [sic] arreté -> arrêté, missing circumflex
% In the print, a curved arrow starts at "b" in the diagram below and points up to "b)" in the line "par b), on trouve ..." above; it is drawn here to the empty node at the top left.
\[
\begin{tikzcd}[column sep=small, row sep=large]
{} & & & & \\
& & \mathrm{a} \arrow[d, Rightarrow] & & W\text{-fibration ultraforte} \\
& & \mathrm{a}' \arrow[d, Rightarrow] & & W\text{-fibration forte} \\
& & \mathrm{b} \arrow[dl, Rightarrow] \arrow[dr, Rightarrow] \arrow[uuull, bend left=30] & & W\text{-fibration} \\
W\text{-fibration gauche} & \mathrm{a}_1 & & \mathrm{a}_2 & W\text{-fibration droite.}
\end{tikzcd}
\]
a$_1$), a$_2$) étaient aussi appelées \textbf{D}-fibrations cohomologique resp. \textbf{D}-fibrations homologiques (\footnote[40]{Non, c'est une notion encore plus forte, à coups de coefficients localement constants, cf. I, p. 121.}). Quand \textbf{D} satisfait Der 3a$'$ (I, p. 111) sur les \textbf{D}-coefficients localement constants sur les $X$ $W_\infty$-asphériques, alors les notions b), a$_1$), a$_2$) sont équivalentes. (Je présume que c'est le cas pour tous les \textbf{D} et tous les $W$ vraiment intéressants.) Par contre, j'ai l'impression que la notion de $W$-fibration forte, voire ultraforte, n'est guère utilisable que dans le cas $W = W_\infty$, où elle coïncide avec les précédentes. C'est dire qu'on peut probablement oublier la terminologie, et se rappeler par contre que dans le cas $W = W_\infty$, les % [sic] cohomologique -> cohomologiques, agreement with "D-fibrations"
\printedpage{17}
% head: A. Grothendieck : Dérivateurs, Ch. XII, Édition des Universités de Montpellier II et Paris VII
$W$-fibrations ont des propriétés d'une force extrême, s'exprimant par les deux propriétés dégagées dans W(4a) et W(4b).

\dgrblock{dgrb475}{Changing the test bases changes the theory}
\begin{dcomment}{exposition}{p19:Par contre, j'ai l'impression que la notion}{Changing the test bases changes the theory}
The revised hierarchy separates the homotopy type of the indexing category from the coefficients used to observe it. A \(W_\infty\)-aspherical category is an absolutely contractible test shape. A \(W\)-aspherical category may instead be invisible only after truncation or after a chosen cohomological localization. Demanding preservation over the latter shapes is stronger when \(W_\infty\subseteq W\).

For \(W_n\), the sphere category used above is \(W_n\)-aspherical but not \(W_\infty\)-aspherical. It is therefore an admissible base in the early strong test \(a'\), and the path family fails that test; it is not an admissible base in the new condition \(b\). This is the concrete reason that the revised notion can accommodate truncated homotopy theories. The source's footnote~40 also corrects a different confusion: a condition tested using every locally constant coefficient is stronger than a condition on the underlying localizer alone. The two distinctions should not be merged.

The opening inference of section~II needs the nondegeneracy hypothesis that Theorem~2 will make explicit. From \(W_{\mathbf D}\supseteq W_\infty\) and W(4), one cannot conclude \(W_{\mathbf D}=W_\infty\) without excluding \(W_{\mathbf D}=\Fl(\Cat)\). The terminal derivator, whose value on every indexing category is the terminal category, gives \(W_{\mathbf D}=\Fl(\Cat)\); all the fibration implications then hold trivially. With \(W_{\mathbf D}\ne\Fl(\Cat)\), the corrected inference follows from Theorems~1 and~2. This is the degenerate exception already recognized at manuscript p.~10, rather than a new obstruction.
\end{dcomment}


\msp{page 20}Dans la théorie de $\mathrm{HOT}_W$, influencé par la réflexion générale sur les dérivateurs et l'axiome Der 3b, j'ai travaillé sous l'hypothèse $W \supseteq W_\infty$, mais en utilisant un axiome de fibration L5 $=$ W(4b) qui, il s'avère, implique finalement $W = W_\infty$. J'ai l'impression que je n'ai jusqu'à présent (propriétés de $f_!$ pour HOT) pas utilisé $W \supseteq W_\infty$, donc ça peut être utile, comme `ça ne coûte pas plus cher' (et oblige d'être plus attentif à ce qu'on utilise), de laisser tomber $W \supseteq W_\infty$ (impliquant $W = W_\infty$), et de mettre par contre sur $W$ des axiomes aussi forts (notamment de fibration) que les propriétés de $W_\infty$, `sans se priver en rien'. Je supposerai donc W(1), W(2), W(3), W(4a) (la forme la plus forte de l'axiome de fibration, cf. pages 13, 14) - c'est tout ce qui a été utile jusqu'à présent. Si on y joint W(5) (axiome de connexité) ça implique $W \subseteq W_\infty$. Mais pour l'instant je ne crois pas en avoir un besoin.

\dgrblock{dgrb476}{A fibre criterion that avoids component restriction}
\begin{dcomment}{development}{p20:Mais pour l'instant je ne crois pas en avoir un besoin.}{A fibre criterion that avoids component restriction}
Grothendieck wants to deduce preservation of components from the nontriviality of the theory. The displayed proof begins by restricting a \(W\)-map to each component of its target; footnote~42 correctly observes that this invokes a stronger fibre-detection statement than W(4b\(^{\prime\prime}\)). The theorem can be proved with the weaker axiom. The following formulation isolates the repair. We did not find this exact formulation in the sources checked for this audit; it is an elementary reorganization of the fibre argument, with no claim of priority.

\textbf{Proposition (two-step nondegeneracy criterion).}\ours\
Let \(V\) be a class of functors containing the categorical homotopisms and satisfying two-out-of-three and absolute theorem A: if every \(A/b\to e\) lies in \(V\), then \(A\to B\) lies in \(V\). Suppose also that for every \(u:A\to B\) in \(V\) and every \(b\in B\), its ordinary homotopy fibre has the weak homotopy type of a category \(F_b\) with \(F_b\to e\in V\). Then either every map of \(V\) induces a bijection on components, or \(V=\Fl(\Cat)\).

\emph{Proof.}
We first record two consequences of the stated axioms. If \(\emptyset\to e\in V\), then every \(\emptyset\to C\) belongs to \(V\) by absolute theorem A, since all its comma categories are empty. Two-out-of-three applied to
\(\emptyset\to C\to e\) gives \(C\to e\in V\). Applying it once more to \(C\to D\to e\) gives every functor \(C\to D\in V\).

If \(D\to e\in V\), then \(D\times D\to e\in V\). Indeed, for \(d\in D\), the comma category of the second projection is
\((D\times D)/d\cong D\times(D/d)\).
The category \(D/d\) has a terminal object, so the projection \(D\times(D/d)\to D\) is a categorical homotopism: its section at that terminal object has composite naturally joined to the identity. Hence \((D\times D)/d\to e\in V\). Absolute theorem A gives \(D\times D\to D\in V\), and composition gives the assertion. When \(D\) is empty the assertion is immediate.

Now take \(u\in V\) for which \(\pi_0u\) is not bijective. If a component of \(B\) is missed, choose \(b\) in it. The homotopy fibre of \(Bu\) at \(b\) is empty: a point in that fibre would supply a path from an image point to \(b\). Consequently its categorical model \(F_b\) is empty as well, because a category with an object has nonempty realization. The preceding paragraph gives \(V=\Fl(\Cat)\).

Suppose instead that two distinct components of \(A\) map into the same component of \(B\). Choose one object in each and paths from their images to some \(b\). These give two points of the homotopy fibre which cannot lie in the same component, since projection to \(BA\) would connect the chosen objects. Its categorical model \(D=F_b\) is therefore disconnected and \(V\)-aspherical. The product argument makes \(D\times D\) \(V\)-aspherical, so its diagonal \(d:D\to D\times D\) belongs to \(V\). Choose objects \(d_0,d_1\) in different components of \(D\). The homotopy fibre of \(Bd\) at \((d_0,d_1)\) is empty: it would otherwise provide a point connected to both \(d_0\) and \(d_1\). Apply the fibre hypothesis to this diagonal. It again gives \(\emptyset\to e\in V\), and hence \(V=\Fl(\Cat)\). This proves the proposition.

For the source's class \(W\), the hypothesis about homotopy fibres follows from its canonical factorization \(A\to\bar A\to B\). The first arrow belongs to \(W_\omega\subseteq W\cap W_\infty\); the second is \(W_\omega\)-perfect and lies in \(W\), so W(4b\(^{\prime\prime}\)) makes its strict fibres \(W\)-aspherical. Ordinary theorem B identifies those fibres with the homotopy fibres of the original map. The proposition therefore proves Theorem~2 using W(1), W(2 bis), W(3), and W(4b\(^{\prime\prime}\)), including the weakened hypotheses asserted in footnote~28 of Theorem~1.
\end{dcomment}


Mais méditant sur la question, je crois avoir trouvé le

\textbf{Théorème 2} (!) (\footnote[41]{Cf. pour un résultat plus fort XV 9.12, p. 133 bis. \textbf{NB} Le résultat est plus général, mais pas vraiment plus fort, car il ne donne pas $W = \Fl(\Cat)$.}) : \textit{Soit} $W \subseteq \Fl(\Cat)$ \textit{satisfaisant} W(1,2,3,4b$''$)\textit{. Si} $W$ \textit{ne satisfait pas à} W(5)\textit{, alors} $W = \Fl(\Cat)$.

\textsc{Démonstration}. Soit $f : X \longrightarrow Y$ dans Cat, tel que $\pi_0(f)$ ne soit pas bijectif. Soit $I$ la catégorie discrète ayant pour ensemble d'objets l'ensemble $\pi_0(Y)$, on a donc un diagramme commutatif
\[
\begin{tikzcd}[column sep=small]
X \arrow[rr, "f"] \arrow[dr, "p"'] & & Y \arrow[dl, "q"] \\
& I, &
\end{tikzcd}
\]
d'autre part un foncteur d'une catégorie dans une catégorie discrète est toujours $W_\omega$-parfait, donc $p$, $q$ sont des $W$-fibrations. Comme\msp{page 21} $f \in W$, il s'ensuit $f_i \in W$ pour tout $i \in I$ (\footnote[42]{On utilise W(4) sous forme plus forte que W(4b$''$)!}).

Si $\pi_0(f)$ n'était pas surjective, prenant $i$ qui n'est pas dans l'image, on aurait $X_i = \emptyset \longrightarrow Y_i \neq \emptyset$, qui serait dans $W$. Soit donc
\[
f : X = \emptyset \longrightarrow Y
\]
dans $W$. Alors $\forall$ $Z$ dans Cat, $f \times \mathrm{id}_Z : \underbrace{X \times Z}_{= \emptyset} \longrightarrow Y \times Z$ est dans $W$ (car $Z$ est
\printedpage{18}
% head: A. Grothendieck : Dérivateurs, Ch. XII, Édition des Universités de Montpellier II et Paris VII
$W$-fibrée sur $e$), donc $\emptyset \longrightarrow Y \times Z$ est dans $W$. Donc aussi la projection
$\begin{tikzcd}[column sep=small, row sep=small, baseline=(\tikzcdmatrixname-1-1.base)]
Y \times Z \arrow[rr] & & Y. \\
& \emptyset \arrow[ul] \arrow[ur] &
\end{tikzcd}$
Comme cette projection est une $W$-fibration, il s'ensuit par W(4b$''$) que ses fibres sont $W$-asphériques. Comme $Y \neq \emptyset$, il y a une telle fibre, elle est $\simeq Z$. Donc $Z$ est $W$-asphérique. Donc on trouve que \textit{toute} catégorie $Z$ dans Cat est $W$-asphérique. Donc toute flèche $Z \longrightarrow Z'$ dans Cat est dans $W$, q.e.d.

Supposons maintenant $\pi_0(f)$ surjective, mais non injective. Il y aurait une fibre de $\pi_0(X) \longrightarrow \pi_0(Y) = I$ qui serait de cardinal $> 1$. En prenant le $f_i$ correspondant, on trouve $X_i \longrightarrow Y_i$ avec $Y_i$ 0-connexe et $\card \pi_0(X_i) \geq 2$. Donc on est ramené au cas
\[
f : X \longrightarrow Y \qquad\qquad \left\{
\begin{array}{l}
f \in W \\[1ex]
Y \text{ 0-connexe} \\[1ex]
\card \pi_0(X) \geq 2.
\end{array}
\right.
\]
Considérons la factorisation canonique
\[
X \xrightarrow{i_f} \bar{X} \xrightarrow{\bar{f}} Y,
\]
avec $i_f \in W$, donc $\bar{f} \in W$. Donc par W(4b$''$) les fibres de $\bar{f}$ sont $W$-asphériques. Mais on voit tout de suite, en décomposant $\bar{X}$ suivant les composantes connexes $X_i$ de $X$, que ces fibres ont un $\pi_0$ de cardinalité $\geq \card \pi_0(X) \geq 2$. Donc on est ramené au cas où $Y = e$, i.e.
\[
X \ W\text{-asphérique}, \qquad \card \pi_0(X) \geq 2.
\]
Considérons $X \times X$, qui est aussi $W$-asphérique, et l'application diagonale\msp{page 22}
\[
d : X \longrightarrow X \times X = Y,
\]
qui est donc dans $W$. Mais $\pi_0(d)$ s'identifie à l'application diagonale $\pi_0(X) \longrightarrow \pi_0(X) \times \pi_0(X)$, qui n'est pas surjective puisque $\card \pi_0(X) \geq 2$. Donc pour la première partie, $W = \Fl(\Cat)$.

Donc sous forme plus positive :

\textbf{Corollaire.} \textit{Soit} $W \subseteq \Fl(\Cat)$ \textit{qui satisfait à} W(1,2,3,4b$''$)\textit{, et telle que} $W \neq \Fl(\Cat)$. \textit{Alors $W$ satisfait l'axiome de connexité} W(5)\textit{, et par suite (th. 1, page 12)} $W \subseteq W_\infty$.

% ---- from Der-ChXIIa-012_p021-025.tex
\printedpage{19}
% head: A. Grothendieck : Dérivateurs, Ch. XII, Édition des Universités de Montpellier II et Paris VII
% foot: Version 7 décembre 2002

\section*{III\quad Classes remarquables de foncteurs liés à $W$.}

Finalement, j'appelerai \textit{‘localiseur fondamental’} une partie $W$ de Fl(Cat) satisfaisant les trois conditions suivantes : % [sic] appelerai -> appellerai, spelling

\begin{enumerate}
\item[(i)] Saturation W(1) (cf. rappel p. 1).
\item[(ii)] Axiome de l'élément final W(3) (page 3).
\item[(iii)] La variante affaiblie W(2 bis) de l'axiome de localisation W(2) donnée dans le corollaire page 1 (anciennement c'était L4$'$ dans Pursuing Stacks).
\end{enumerate}

Quand j'aurai besoin de plus sur $W$, je le dirai explicitement. Ces trois axiomes ont suffi pour pratiquement tout ce que j'ai fait avec $W$ dans Pursuing Stacks, la seule exception étant un usage (si je me rappelle) unique, quelque part, d'une variante affaiblie de W(2) (pour les $S$-foncteurs \textit{cartésiens} entre $S$-catégories \textit{fibrées} - si sur toute fibre il est dans $W$, il est lui-même dans $W$).

\dgrblock{dgrb477}{The localizer retained for the geometric argument}
\begin{dcomment}{exposition}{p22:Quand j'aurai besoin de plus sur}{The localizer retained for the geometric argument}
Grothendieck now separates the part of his theory that depends on a notion of weak equivalence from the much stronger requirements imposed on fibrations. This is the point of retaining only W(1), W(3), and W(2 bis). The distinction is consequential: a map can have a well-behaved theory of lifting along arrows without preserving every weak equivalence under pullback. Throughout this section, a category is \emph{$W$-aspherical} when its map to the terminal category belongs to $W$. For $W=W_\infty$, this means that its nerve has weakly contractible realization. The word does not mean that only its higher homotopy groups vanish; connectedness and trivial fundamental group are included.

The notation for comma categories fixes the direction of every assertion that follows. For $f:X\to Y$, an object of $X/y$ is a pair $(x,\alpha:f(x)\to y)$, whereas an object of $y\backslash X$ is $(x,\alpha:y\to f(x))$. The strict fibre $X_y$ contains the objects mapping \emph{equal} to $y$ and only those arrows mapping to $1_y$. Even when several objects of $Y$ are isomorphic to $y$, they do not enter this strict fibre. These three categories solve different problems. The comma categories allow a comparison with $y$; the strict fibre requires equality over $y$.

An aspherical functor has all $X/y$ aspherical. W(2 bis) then says that it belongs to $W$. This is the absolute form of Quillen's Theorem A. W(2), by contrast, tests a map $X\to Y$ over an arbitrary common base $S$ through all $X/s\to Y/s$. The formally stronger axiom will be named again at the corollary on manuscript page 30; it is not imported into the preceding elementary proofs. In later terminology, the distinction is between the absolute and relative Theorem A axioms: \citet[\S1.1]{Maltsiniotis2012Exact} calls the version with the absolute axiom a \emph{weak} fundamental localizer and reserves ``fundamental localizer'' for the relative version. \citet[Theorem 6.1.18]{Cisinski2006} proves that the minimal weak fundamental localizer is $W_\infty$. That minimality theorem does not assert that every weak fundamental localizer satisfies the relative axiom. The manuscript's terminology is still settling here.

The next propositions convert these axioms into geometric tests. Their economy comes from arranging a comparison so that a category with a terminal or initial object appears, and then using two-out-of-three. No model structure on $\Cat$, and no derived direct-image functor, is needed for that argument.
\end{dcomment}


On suppose donc par la suite donné un localiseur fondamental
\[
W \subseteq \mathrm{Fl}(\mathrm{Cat}).
\]

\textbf{Proposition 1.} \textit{Soit $f : X \longrightarrow Y$ dans} \textup{Cat}. \textit{Conditions équivalentes :}

\begin{enumerate}
\item[\textit{a)}] \textit{Pour tout isomorphisme local $Y' \longrightarrow Y$, $f' : X' \longrightarrow Y'$ est $\in W$.}
\item[\textit{b)}] \textit{Même condition, dans le cas particulier $Y' = Y/y \longrightarrow Y$, donc $\forall\ y \in Y$, $X/y \simeq X \times_Y (Y/y)$ est $W$-asphérique.}
\end{enumerate}

\msp{page 23}\textbf{Définition 1.} Sous ces conditions équivalentes a), b) de la proposition 1, on dit que $f$ est \textit{$W$-asphérique.} On dit \texttt{[que]} $f$ est \textit{$W$-coasphérique} si $f^{\circ}$ est $W$-asphérique, i.e. si les $y\backslash X$ sont $W$-asphériques (pour $y \in \Ob Y$).

\textsc{Démonstration.} Il suffit de prouver b) $\Longrightarrow$ a). Par définition d'un isomorphisme local, on voit que la condition b) est stable par changement de base $Y' \longrightarrow Y$, avec $Y'$ un isomorphisme local (puisque les $Y'/y' \xrightarrow{\sim} Y/y$, donc $X'/y' \xrightarrow{\sim} X/y$). Il suffit donc de voir que si $f$ satisfait à b), alors $f \in W$. Mais c'est l'axiome W(2 bis).

\begin{center}
\begin{tikzcd}[row sep=large]
X \arrow[d] & & & \\
Y & & Y' \arrow[ll] & \\
& Y/y \arrow[ul] & & Y'/y' \arrow[ll] \arrow[ul]
\end{tikzcd}
\end{center}

\printedpage{20}
% head: A. Grothendieck : Dérivateurs, Ch. XII, Édition des Universités de Montpellier II et Paris VII
% foot: Version 7 décembre 2002

\textbf{Remarque.} Si $f$ a un adjoint à droite $g$, alors pour tout $y \in Y$, $X/y \simeq X/g(y)$, donc $X/y$ a un objet final et est donc $W$-asphérique. Donc si $f$ a un adjoint à droite (à gauche), alors $f$ est $W$-asphérique ($W$-coasphérique).

Mais même si $f$ n'a pas d'adjoint à droite, il y a quand même un adjoint généralisé
\[
g : Y \longrightarrow X^{\wedge}
\]
et on a
\[
X/y \simeq X/g(y) \qquad \text{(pour l'inclusion pleinement fidèle $X \hookrightarrow X^{\wedge}$),}
\]
donc $\mathrm{Top}(X/y)$ s'identifie au topos induit $\mathrm{Top}(X)/g(y)$. La condition $f$ $W$-asphérique signifie donc que les faisceaux $g(y)$ dans $\mathrm{Top}(X)$ sont
\msp{page 24}
$W$-asphériques. Pour de nombreuses autres formulations équivalentes, je renvoie à Pursuing Stacks.

\dgrblock{dgrb478}{The generalized right adjoint is a presheaf of arrows}
\begin{dcomment}{exposition}{p23:Mais même si $f$ n'a pas d'adjoint à droite}{The generalized right adjoint is a presheaf of arrows}
The passage through presheaves explains why the comma categories are the correct substitutes for an adjoint. Define
\[
 g(y)(x)=\operatorname{Hom}_Y(f(x),y),\qquad
 g:Y\longrightarrow [X^{\mathrm{op}},\mathrm{Set}].
\]
Precomposition makes $g(y)$ a presheaf on $X$, and postcomposition makes $g$ covariant in $y$. Its category of elements has objects $(x,\alpha)$ and arrows $h:x\to x'$ satisfying $\alpha=\alpha'f(h)$. It is therefore exactly $X/y$. If $f\dashv r$, the adjunction identifies $g(y)$ with the representable presheaf $X(-,r(y))$; the category of elements then becomes $X/r(y)$ and has terminal object $1_{r(y)}$. The asphericity of a functor admitting a right adjoint is thus the representable case of the construction in the text.

The asserted identification of topoi can also be made explicit. For any presheaf $P$ on $X$, a presheaf $A\to P$ determines on its category of elements the presheaf
\[
 (x,p)\longmapsto\{a\in A(x):a\mapsto p\}.
\]
Conversely, from a presheaf $B$ on $\int_X P$, recover $A(x)$ as the disjoint union of $B(x,p)$ over all $p\in P(x)$. Restriction along $h:x\to x'$ uses the arrow $(x,P(h)p')\to(x',p')$. These operations are inverse up to their canonical isomorphisms and yield
\[
 [ (\textstyle\int_X P)^{\mathrm{op}},\mathrm{Set}]
 \simeq [X^{\mathrm{op}},\mathrm{Set}]/P.
\]
Taking $P=g(y)$ proves the statement about $\operatorname{Top}(X/y)$. Its use here is interpretive: the failure of an actual right adjoint is measured by the homotopy type of a presheaf of possible arrows, rather than by the absence of a distinguished representing object.

The commentary on Chapter~I, manuscript pages~64--66, discusses local constancy and the Yoneda argument. Here the additional point is that an adjoint is sufficient but far from necessary: it supplies a terminal element of every comma category, while $W$-asphericity requires only their homotopical triviality.
\end{dcomment}


\textbf{Proposition 2.} \textit{Soit $X \xrightarrow{f} Y \xrightarrow{g} Z$ dans} \textup{Cat}. \textit{Si $f$ et $g$ sont $W$-asphériques ($W$-coasphériques), de même $gf$. Si $gf$ et $f$ sont $W$-asphériques ($W$-coasphériques), $g$ aussi.}

\textbf{Rappel.} Un morphisme $p : S' \longrightarrow S$ dans Cat est un \textit{isomorphisme local} si $\forall\ s' \in \Ob S'$, le morphisme induit
\[
p_{s'} : S'/s' \xrightarrow{\sim} S/s \qquad (s = p(s'))
\]
est un isomorphisme : $p$ induit un isomorphisme sur les ‘localisées’ de $S'$ et $S$ en $s'$ et $s = p(s')$ respectivement. Il est immédiat que cette notion est stable par composition, que si $gf$ et $g$ sont des isomorphismes locaux, alors $f$ l'est aussi, que tout isomorphisme est un isomorphisme local. De plus, les isomorphismes locaux sont \textit{stables par changement de base,} ce qui résulte du fait (trivial et important) que si on a un carré cartésien dans Cat

\begin{center}
\begin{tikzcd}[row sep=large]
X & X' \arrow[l, "q"'] \arrow[d, "f'"] \\
S \arrow[from=u, "f"'] & S', \arrow[l, "p"']
\end{tikzcd}
\end{center}

et si $x' \in \Ob X'$ a pour image $x$, $s'$, $s$ dans $X$, $S'$, $S$, alors le diagramme

\begin{center}
\begin{tikzcd}[row sep=large]
X/x \arrow[d, "f/x"'] & X'/x' \arrow[l, "q/x'"'] \arrow[d, "f'/x'"] \\
S/s & S'/s' \arrow[l, "p/s'"']
\end{tikzcd}
\end{center}

est également cartésien. (Donc si $p/s'$ est \texttt{[un]} isomorphisme, $q/x'$ aussi.)

\dgrblock{dgrb479}{Local isomorphisms and the scope of cancellation}
\begin{dcomment}{exposition}{p24:est également cartésien.}{Local isomorphisms and the scope of cancellation}
The ``local isomorphisms'' in this argument are discrete fibrations. Indeed, an isomorphism $S'/s'\to S/p(s')$ says that every arrow ending at $p(s')$ has exactly one lift ending at $s'$, and that its factorization triangles lift uniquely. Conversely, that unique-lifting property gives the displayed isomorphism of slice categories. The categories of elements of presheaves supply the basic examples. In particular, $Y/y\to Y$ is such a map.

This description makes Proposition 1 transparent. If every $X/y$ is $W$-aspherical and $q:Y'\to Y$ is a local isomorphism, then, for $y'\in Y'$ with image $y$, strict pullback gives
\[
 (X\times_Y Y')/y'\cong X/y.
\]
Theorem A applied over $Y'$ puts $X\times_Y Y'\to Y'$ in $W$. Conversely, the particular base changes $Y/y\to Y$ recover the same condition because $Y/y$ has a terminal object.

The cancellation in Proposition 2 uses precisely this class of allowed base changes. Given $X\xrightarrow fY\xrightarrow gZ$ and a local isomorphism $Z'\to Z$, form the three pullbacks $X',Y',Z'$. Its pullback $Y'\to Y$ is again a local isomorphism. Thus asphericity of $f$ gives $f'\in W$. If $g$ is also aspherical, $g'\in W$ and hence $g'f'\in W$. If instead $gf$ is aspherical, then $g'f'\in W$ and two-out-of-three gives $g'\in W$. Proposition 1 now gives the two assertions. The opposite-category argument proves the coaspherical forms.

There is no cancellation statement with $g$ and $gf$ implying asphericity of $f$. For example, let $f$ select the terminal object of $[1]$, and let $g:[1]\to e$. Both $g$ and $gf$ are aspherical, whereas the comma category of $f$ over $0$ is empty. For any $W$ with $\varnothing\to e\notin W$, $f$ is not aspherical. This elementary directional distinction is exactly what the local-isomorphism proof preserves.
\end{dcomment}


\printedpage{21}
% head: A. Grothendieck : Dérivateurs, Ch. XII, Édition des Universités de Montpellier II et Paris VII
% foot: Version 7 décembre 2002

\msp{page 25}Revenons à la proposition 2. La stabilité de la notion de morphisme asphérique par composition résulte aussitôt de la forme a) de la condition. De même pour la condition que si $gf$ et $f$ sont $W$-asphériques, $g$ aussi - il faut montrer que si $Z' \longrightarrow Z$ est \texttt{[un]} isomorphisme local, alors $g'$ est dans $W$, or l'hypothèse implique que $g'f'$ et $f'$ le sont.

\begin{center}
\begin{tikzcd}[row sep=large]
X \arrow[d, "f"'] & X' \arrow[l, "r"'] \arrow[d, "f'"] \\
Y \arrow[d, "g"'] & Y' \arrow[l, "q"'] \arrow[d, "g'"] \\
Z & Z' \arrow[l, "p"']
\end{tikzcd}
\end{center}

\textbf{Corollaire 1.} \textit{L'image inverse d'un $f : X \longrightarrow Y$ qui est $W$-asphérique (resp. $W$-coasphérique) par un isomorphisme local (resp. par un coïsomorphisme local) est $W$-asphérique ($W$-coasphérique).}

\begin{center}
\begin{tikzcd}[row sep=huge]
X \arrow[d] & X' \arrow[l] \arrow[d] & X'' \arrow[l] \arrow[d] \\
Y & Y' \arrow[l, "\text{iso local}"] & Y'' \arrow[l, "\text{iso local}"]
\end{tikzcd}
\end{center}

Résulte formellement de la définition par a) et du fait qu'un composé de deux isomorphismes locaux (resp. coïsomorphismes locaux) est itou.

\textbf{NB} Bien sûr, la notion de foncteur $W$-asphérique n'est pas stable par changement de base quelconque. Par exemple, l'inclusion $X = \{y_0\} \longrightarrow Y$ est $W$-asphérique si et seulement si les catégories discrètes $\mathrm{Hom}(y_0, y)$ sont $W$-asphériques, i.e. (sous réserve que l'on ait W(5), axiome de connexité) si et seulement si les $\mathrm{Hom}(y_0, y)$ sont de cardinal 1, i.e. $y_0$ \textit{objet initial} de $Y$. Mais si $Y \neq \{y_0\}$, prenant $y$ dans $Y$, $y \neq y_0$, et posant $Y' = \{y\} \longrightarrow Y$, on a $X' = X \times_Y Y' = \emptyset$ et $\emptyset \longrightarrow e$ n'est pas asphérique!

\msp{page 26}Par contre, on s'attend que les morphismes $W$-asphériques soient stables par changements de base \textit{$W$-lisses} (qui généralisent les isomorphismes locaux), et les morphismes $W$-coasphériques par changements de base $W$-propres. Cf. plus bas.

\dgrblock{dgrb480}{Properness asks for contractible categories of lifts}
\begin{dcomment}{exposition}{p26:Par contre, on s'attend que les morphismes}{Properness asks for contractible categories of lifts}
The anticipated base-change theorem needs a condition on the map along which one pulls back. Grothendieck finds it by inspecting every possible lift of an arrow. For $f:X\to Y$, an object $x\in X$, and $u:f(x)\to y$, write $L_f(x,u)$ for the category whose objects are arrows $a:x\to z$ with $f(a)=u$ and whose morphisms are arrows $h:z\to z'$ over $1_y$ satisfying $ha=a'$. Its morphisms matter as much as its objects: merely asking for at least one lift would lose the homotopy information.

Let $j_y:X_y\to X/y$ send $z$ to $(z,1_y)$. The object $(x,u)$ of $X/y$ gives an equality of categories
\[
 (x,u)\backslash j_y\cong L_f(x,u).
\]
The same category is the strict fibre over $u$ of $x\backslash X\to f(x)\backslash Y$. These identifications prove the equivalence of conditions (a), (b), and (d) in Proposition 3. The edition's intermediate list ``a) b) c)'' on manuscript page 27 should read ``a) b) d)''; the remaining argument establishes (c).

To see that argument without suppressing an endpoint, pull $f$ back along the arrow $u:[1]\to Y$. Write $E=X\times_Y[1]$. For an object in $E_1$, the comma category into $E_1\hookrightarrow E$ has an initial object. For an object $(x,0)$ in $E_0$, that comma category is $L_f(x,u)$. Consequently $E_1\hookrightarrow E$ is coaspherical exactly when all these lift categories are aspherical. This is why a test over the walking arrow is sufficient for \emph{properness}: the test asks for coasphericity of the inclusion, not merely that its total homotopy type be an equivalence.

The dual incoming-lift category consists of arrows $a:z\to x$ over a fixed $u:y\to f(x)$. Its morphisms are vertical factorizations. Requiring these categories to be $W$-aspherical defines smoothness. The terminology records two directions of the same lifting geometry. The commentary on Chapter~I, manuscript pages~52 and~68, explains the coefficient-theoretic origin; here Grothendieck obtains the criterion directly from $W$.
\end{dcomment}


\textbf{Corollaire 2.} \textit{Soit $f : X \longrightarrow Y$ dans} \textup{Cat} \textit{une} \textup{équivalence} \textit{de catégories. Alors $f$ est $W$-asphérique et $W$-coasphérique.}

En effet, $f$ a un adjoint à gauche et un adjoint à droite, et on utilise la remarque page 23.

\textbf{Proposition 3.} \textit{Soit $f : X \longrightarrow Y$ dans} \textup{Cat}. \textit{Conditions équivalentes :}

\printedpage{22}
% head: A. Grothendieck : Dérivateurs, Ch. XII, Édition des Universités de Montpellier II et Paris VII
% foot: Version 7 décembre 2002

\begin{enumerate}
\item[\textit{a)}] \textit{Pour tout $y \in \Ob Y$, le foncteur d'inclusion}
\[
X_y \longrightarrow X/y \tag*{$(*)$}
\]
\textit{est $W$-}co\textit{asphérique.}
\item[\textit{b)}] \textit{Pour tout $x_0 \in X$, le morphisme induit pour les colocalisées}
\[
x_0\backslash X \longrightarrow y_0\backslash Y \qquad \textit{(où $y_0 = f(x_0)$)}
\]
\textit{est à fibres $W$-asphériques.}
\item[\textit{c)}] \textit{Pour tout morphisme $\mathbf{\Delta}^1 = Y' \longrightarrow Y$, considérons $X' = X \times_Y Y'$ sur $\mathbf{\Delta}^1$, l'inclusion $X'_1 \longrightarrow X'$ de la fibre fermée de $X'$ (fibre sur le point fermé de $Y'$) est $W$-coasphérique.}
\item[\textit{d)}] \textit{Pour toute flèche $u_0 : y_0 \longrightarrow y_1$ dans $Y$, et tout $x_0$ dans $X_{y_0}$, la catégorie $X(x_0;u_0)$ formée des flèches $u' : x_0 \longrightarrow x$ dans $X$ de source $x_0$, et qui relèvent $u$} \texttt{[plutôt} $u_0$\texttt{]}\textit{, est $W$-asphérique.}
\end{enumerate}

\msp{page 27}\textsc{Démonstration.} Pour expliciter la condition a), il faut expliciter la catégorie $\xi_0\backslash(X_y)$, pour un objet quelconque $\xi_0 = (x_0, u_0)$ de $X/y$, $x_0 \in \Ob X$,
\[
u_0 : y_0 \longrightarrow y_1 = y \quad \text{dans $Y$, avec $y_0 = f(x_0)$.}
\]

Cette catégorie $\xi_0\backslash X_{y_1}$ s'explicite alors comme la catégorie des $x_1$ dans $X_{y_1}$, \textit{munis} d'une flèche de $\xi_0$ dans l'image $\xi_1 = (x_1, \mathrm{id}_{f(x_1)} = \mathrm{id}_{y_1})$ de $x_1$ dans $X/y_1$ dans $\xi_0$ \texttt{[‘dans} $\xi_0$\texttt{’ à effacer]}, i.e. d'une flèche $u$ de $x_0$ dans $x_1$, telle que $f(u)$ donne lieu au triangle commutatif dans $Y$, i.e. $f(u) = u_0$. % AUDIT resolved: regular CMTT12 editorial bracket; see corrections.csv, PDF leaf 32

\begin{center}
\begin{tikzcd}
x_1 \arrow[dd, dash] & & \\
& & x_0 \arrow[ull, "u"'] \arrow[dd, dash] \\
y_1 & y_1 \arrow[l, "\mathrm{id}"'] & \\
& & y_0 \arrow[ul, "u_0"'] \arrow[ull, "f(u)"]
\end{tikzcd}
\end{center}

Donc c'est la catégorie notée $X(x_0;u_0)$ dans d). On voit tout de suite qu'elle est aussi isomorphe à la fibre de $x_0\backslash X \longrightarrow y_0\backslash Y$ en l'élément $u_0$ de $y_0\backslash Y$, dont il est question dans b). D'où l'équivalence de a) b) c). D'autre part, la donnée de $u_0$ dans $Y$ équivaut à celle d'un foncteur $\mathbf{\Delta}^1 \longrightarrow Y$, comme dans c), et pour exprimer alors que $\underbrace{X'_1}_{\simeq X_{y_1}} \longrightarrow X'$ est coasphérique, il suffit d'exprimer que les $x'\backslash X'_1$ ($x' \in \Ob X'$) sont $W$-asphériques. Or c'est trivial si $x' \in \Ob X'_1$ (objet initial), donc il suffit de le vérifier pour $x'$ dans $X'\backslash X'_1 = X'_0 \simeq X_{y_0}$. Alors $x'$ s'identifie à un élément $x_0$ de $X_{y_0}$, et on vérifie tout de suite que $x_0\backslash X'_1$ est isomorphe à la catégorie $X(x_0;u_0)$ de d). Cela achève la démonstration.

\printedpage{23}
% head: A. Grothendieck : Dérivateurs, Ch. XII, Édition des Universités de Montpellier II et Paris VII
% foot: Version 7 décembre 2002

\msp{page 28}\textbf{Définition 2.} Sous les conditions de la proposition 3, on dit que $f$ est \textit{$W$-propre.} On dit que $f$ est \textit{$W$-lisse} si $f^{\circ}$ est $W$-propre, i.e. si les morphismes d'inclusion $(*)$ \texttt{[plutôt} $X_y \longrightarrow y\backslash X$\texttt{]} sont $W$-asphériques etc. (conditions duales de celles de la proposition 3).

\textbf{Remarque.} La condition a) de la proposition 3 est satisfaite en particulier si l'inclusion $(*)$ a un adjoint à gauche (cf. remarque page 23). Cela signifie que les catégories $X(x_0;u)$ ont un objet initial. Cela signifie aussi que la catégorie $X$ sur $Y$ est Cat-cofibrée sur $Y$ au sens faible, i.e. qu'il existe, pour $u_0 : y_0 \longrightarrow y_1$ flèche dans $Y$, et $x_0$ dans $X_{y_0}$, un objet $u_{0*}(x_0)$ dans $X_{y_1}$ qui représente le foncteur covariant $X_{y_1} \longrightarrow (\mathrm{Ens})$, $x_1 \longmapsto \mathrm{Hom}_u(x_0, x_1)$. Ainsi, si $X$ est Cat-cofibrée au sens faible sur $Y$, et a fortiori si elle est Cat-cofibrée, alors $f : X \longrightarrow Y$ est $W$-propre. Dualement, si $f$ est faiblement Cat-fibrée, $f$ est \texttt{[}$W$\texttt{-]}lisse.

\dgrblock{dgrb481}{Representable lifts and their homotopical relaxation}
\begin{dcomment}{examples}{p28:La condition a) de la proposition 3}{Representable lifts and their homotopical relaxation}
Grothendieck first places ordinary fibred categories inside the new theory. Fix $x\in X$ and $u:f(x)\to y$. An initial object $a:x\to u_*x$ of $L_f(x,u)$ gives, for every $z\in X_y$, a bijection
\[
 \operatorname{Hom}_{X_y}(u_*x,z)
 \xrightarrow{\ \sim\ }
 \{b:x\to z:f(b)=u\},\qquad h\longmapsto ha.
\]
Conversely, such a representing arrow is initial in $L_f(x,u)$. Choosing it for all $(x,u)$ is the weak cofibred condition used here. It only asserts this factorization property for arrows over the fixed arrow $u$; the stronger condition on composition is part of the usual opfibration formalism. Either condition gives properness, since initial categories are $W$-aspherical. The dual argument gives smoothness for prefibrations.

A useful calculation shows how much weaker asphericity can be than representability. Let $C$ be the poset with three objects $b,d,e$ and relations $b<d$, $b<e$. Form its left cone $X=C^{\lhd}$ by adjoining an initial object $a$, and let $p:X\to[1]$ send $a$ to $0$ and $C$ to $1$. The outgoing-lift category from $a$ over $0\to1$ is $C$ itself. It is contractible because $b$ is initial. This example is already a weak opfibration. To remove representability while retaining contractibility, replace $C$ by the four-object fence
\[
 b<d,\qquad e<d,\qquad e<c.
\]
Its nerve is the path $b-d-e-c$, hence contractible, but it has no initial object. In the corresponding cone, $L_p(a,0\to1)\cong C$ is therefore aspherical without being representable by a universal lift. For every identity in the base, the lift category is a coslice in the appropriate fibre and has an initial object. There are no other base arrows to test. Thus this $p$ is $W_\infty$-proper but is not weakly opfibred.

The dual cone gives a smooth map without the prefibration property. \citet[Proposition 4.3 and Example 4.17]{Maltsiniotis2005Asphericite} places these criteria in the broader language of asphericity structures; for a structure coming from a fundamental localizer, the resulting weak smoothness criterion recovers $W$-smoothness. For an arbitrary asphericity structure those notions need not agree.
\end{dcomment}


\textbf{Proposition 4.} (\footnote[43]{Dire aussi : Si $X \xrightarrow{f} Y \xrightarrow{g} Z$, si $gf$ et $f$ $W$-lisses, alors $g$ est lisse en les points de $\mathrm{Im}\, f$, idem pour propre.}) \textit{La notion de morphisme $W$-propre est stable par composition, par changement de base. Les coisomorphismes locaux sont propres (trivial sur b)), a fortiori les isomorphismes sont propres, les immersions fermées sont propres.}

\textit{Énoncé dual pour les morphismes lisses : stabilité par changement de base, par composition, les isomorphismes locaux sont lisses, notamment les immersions ouvertes.}

\msp{page 29}La stabilité par changement de base est immédiate sur la forme c). La stabilité par composition \texttt{[est]} moins triviale :

\begin{center}
\begin{tikzcd}[row sep=large]
{\scriptstyle x} &[-1.5em] X \arrow[d, "f"] & x\backslash X \arrow[l] \arrow[d] & \bar{X} \arrow[l] \arrow[d] \\
{\scriptstyle y} & Y \arrow[d, "g"] & y\backslash Y \arrow[l] \arrow[d] & \bar{Y} \arrow[l] \arrow[d] \\
{\scriptstyle z} & Z & z\backslash Z \arrow[l] & \{u\} \arrow[l]
\end{tikzcd}
\end{center}

Dans le diagramme à carrés cartésiens ci-contre \texttt{[plutôt dont les deux carrés de droite sont cartésiens]} (déduit de $x \in \Ob X$, d'où $y = f(x)$, $z = g(y) = gf(x)$, et de $u \in \Ob z\backslash Z$), il faut prouver que $\bar{X}$ est $W$-asphérique. On sait que $\bar{Y}$ l'est ($g$ propre) et que $\bar{X} \longrightarrow \bar{Y}$ est à fibres $W$-asphériques. Il suffit donc d'utiliser la proposition 5 ci-dessous, et la remarque immédiate % AUDIT resolved: regular CMTT12 editorial bracket; see corrections.csv, PDF leaf 33

\textbf{Lemme.} \textit{Si $f$ est propre, il en est de même de $x\backslash f : x\backslash X \longrightarrow y\backslash Y$, pour tout $x$ dans $X$, $y = f(x)$.}

% ---- from Der-ChXIIa-012_p026-030.tex
\printedpage{24}
% head: A. Grothendieck : Dérivateurs, Ch. XII, Édition des Universités de Montpellier II et Paris VII
% foot: Version 7 décembre 2002
En effet, si $u'$ est dans $x\backslash X$ et $v'$ dans $y\backslash Y$, et si $x'$, $y'$ sont leurs images dans $X$, $Y$ respectivement, alors
\[
u'\backslash(x\backslash X)\simeq x'\backslash X,\qquad v'\backslash(y\backslash Y)\simeq y'\backslash Y,
\]
et la flèche entre les deux s'identifie à la flèche de localisation
\[
x'\backslash X\longrightarrow y'\backslash Y.
\]

\dgrblock{dgrb482}{The composition proof has no circular dependence}
\begin{dcomment}{exposition}{p29:et la flèche entre les deux s'identifie à la flèche de localisation}{The composition proof has no circular dependence}
The proof of composition proceeds by factoring a category of lifts through another category of lifts. It uses Proposition 5, but Proposition 5 itself needs only the definition of properness and the already established base-change stability. Thus its appearance after Proposition 4 does not create a circular proof.

First prove the part of Proposition 5 that will be used. If $q:A\to B$ is proper and every strict fibre $A_b$ is aspherical, then $A_b\to A/b$ is coaspherical and hence belongs to $W$. Two-out-of-three makes $A/b$ aspherical, and W(2 bis) gives $q\in W$. The same reasoning applies after every base change, because properness and the asphericity of the fibres persist. The smooth version follows by taking opposites.

Now let $X\xrightarrow fY\xrightarrow gZ$ be proper. Fix $x\in X$ and $u:gf(x)\to z$. There is a functor
\[
 L_{gf}(x,u)\longrightarrow L_g(f(x),u),\qquad a\longmapsto f(a).
\]
It is the base change of $x\backslash f:x\backslash X\to f(x)\backslash Y$ along the fibre of $f(x)\backslash g$ over $u$. The coslice functor $x\backslash f$ is proper: taking the coslice at an arrow $x\to x'$ reduces its incoming data to $x'\backslash X\to f(x')\backslash Y$, precisely the original lifting test. Its displayed base change is therefore proper. The fibre over a lift $v:f(x)\to y'$ is $L_f(x,v)$, which is aspherical. Proposition 5 makes the displayed functor a $W$-equivalence. Since its target is aspherical by properness of $g$, its source is aspherical too. This proves properness of $gf$.

The same diagram proves the observation in footnote 43. If $f$ and $gf$ are proper, its source is aspherical, and the displayed proper functor still has aspherical fibres. Its target is therefore aspherical. This proves the lifting condition for $g$ at every object $f(x)$ in the strict image of $f$; it does not assert anything about objects outside that image. Duality proves the smooth statement. These arguments require W(2 bis), not the full relative axiom W(2).
\end{dcomment}


\textbf{Proposition 5.} \textit{Soit $f : X\longrightarrow Y$ dans $\Cat$, supposons $f$ $W$-propre ou $W$-lisse, et à fibres $W$-asphériques, alors $f\in W$ (et même $f$ est} \textup{universellement} \textit{dans $W$, i.e. il est dans $W$ et le reste par tout changement de base).}

\leavevmode\msp{page 30}Comme les hypothèses sont stables par changement de base, il suffit de prouver $f\in W$. Supposons p.~ex. $f$ propre, prouvons que $f$ est $W$-\textit{asphérique}, i.e. que les $X/y$ ($y\in\Ob\,Y$) sont $W$-asphériques. Mais $X_y\longrightarrow X/y$, étant $W$-coasphérique par hypothèse de propreté, est dans $W$ (par W(2 bis)), et comme $X_y$ est $W$-asphérique par hypothèse, $X/y$ aussi. Mais par W(2 bis), $f$ $W$-asphérique $\Longrightarrow f\in W$. Le cas $f$ lisse est dual.

\textbf{Remarque.} Ainsi, si $f$ est propre, s'il est à fibres $W$-asphériques il est `universellement $W$-asphérique', i.e. est et reste $W$-asphérique par tout changement de base. Inversement, si $f$ est universellement dans $W$ (pas forcément propre), alors en particulier ses fibres sont $W$-asphériques. Mais ça n'implique bien sûr pas que $f$ soit propre, ou lisse (\footnote[44]{Mais `universellement dans $W$' implique déjà `universellement $W$-asphérique', puisque la notion de $W$-asphéricité est définie par \texttt{[la]} proposition 1 a).}).

\dgrblock{dgrb483}{A universal equivalence with neither lifting property}
\begin{dcomment}{examples}{p30:Mais ça n'implique bien sûr pas}{A universal equivalence with neither lifting property}
The qualification after Proposition 5 has a small model which makes the role of the lifting condition visible. Let $P$ be the poset with objects $u,v,w,z$ and exactly the strict relations
\[
 u<v,\qquad u<z,\qquad w<z.
\]
Define $p:P\to[1]$ by $p(u)=p(v)=0$ and $p(w)=p(z)=1$. Both fibres are intervals, and the nerve of $P$ is the path $v-u-z-w$. More is true than the contractibility of these three nerves: $p$ has a contraction over its base.

The section $s:[1]\to P$ selects $u$ and $z$. Let $R:P\to P$ send $v$ to $u$ and fix $u,w,z$. Each displayed relation is preserved, so $R$ is a functor over $[1]$. The pointwise inequalities give natural transformations over $[1]$
\[
 \operatorname{id}_P\longleftarrow R\longrightarrow sp.
\]
For the first transformation use $u\to v$ at $v$; for the second use $w\to z$ at $w$; all remaining nonidentity components are absent. Since $ps=\operatorname{id}_{[1]}$, these transformations exhibit $p$ as a homotopy equivalence of categories over $[1]$. After any base change $T\to[1]$, the section and both transformations pull back, so $P\times_{[1]}T\to T$ is still a homotopy equivalence. It therefore belongs to every $W$ considered in this section.

Nevertheless $p$ is not $W_\infty$-proper: the arrow $0\to1$ has no outgoing lift from $v$, so its lift category is empty. Nor is it $W_\infty$-smooth: the same arrow has no incoming lift ending at $w$. Thus universal membership in $W$ does not imply either geometric lifting property. The same failure holds for any $W$ for which the empty category is not aspherical.

This also verifies the terminology in footnote 44. A universal $W$-equivalence is universally aspherical because, after any base change and any further local-isomorphism test, one is still taking a base change of the original map. The four-object example is universally aspherical and universally coaspherical without being proper or smooth. The latter properties inspect the lift category at each particular object, which a relative contraction may move.
\end{dcomment}


\textbf{Corollaire.} \textit{Si $W$ satisfait} \textup{W(2) (pas seulement W(2 bis))}\textit{, alors pour un $S$-morphisme $X\longrightarrow Y$ de catégories qui sont toutes deux $W$-propres sur $S$, ou toutes deux $W$-lisses, si $f_s : X_s\longrightarrow Y_s$ est dans $W$ pour tout $s\in S$, alors $f\in W$.}

\leavevmode\msp{page 31}Supposons p.~ex. $f$ propre, considérons, pour $s\in\Ob\,S$
\[
\begin{tikzcd}[row sep=large, column sep=large]
X_s \arrow[r, "f_s"] \arrow[d, "i_s^X"'] & Y_s \arrow[d, "i_s^Y"] \\
X/s \arrow[r, "f/s"] & Y/s,
\end{tikzcd}
\]
comme $X$, $Y$ \texttt{[sont]} propres, les flèches verticales sont dans $W$ (il suffirait de beaucoup moins que la propreté \dots), donc $f_s\in W\Longrightarrow f/s\in W$. Ainsi, `$f$ est dans $W$ localement sur $S$', ce qui implique $f\in W$ par W(2).

\dgrblock{dgrb484}{Fibrewise comparison uses the structural maps}
\begin{dcomment}{exposition}{p31:les flèches verticales sont dans $W$}{Fibrewise comparison uses the structural maps}
The corollary compares two families over $S$; it does not assume that the comparison $f:X\to Y$ is itself proper. The phrase ``Supposons p.~ex.\ $f$ propre'' at the start of the proof is an edition slip. The case being considered is that both structural maps $X\to S$ and $Y\to S$ are proper. This is confirmed by the next sentence and by the two vertical arrows in the displayed square.

In fact, the proof needs only
\[
 X_s\longrightarrow X/s\in W,
 \qquad Y_s\longrightarrow Y/s\in W
 \quad\text{for every }s\in S.
\]
If the fibre map $f_s$ belongs to $W$, the commutative square and two-out-of-three imply $f/s\in W$. The remaining step from all these slice comparisons to $f\in W$ is exactly the relative axiom W(2). Replacing all slices by coslices gives the smooth case, using the opposite form of the same axiom.

This distinction will persist when properness is weakened to preproperness. The geometric hypothesis supplies comparisons between fibres and comma categories; the relative localization axiom assembles the comma comparisons over the whole base. There is no appeal to the converse assertion that a weak equivalence between total categories must be a weak equivalence on fibres. That substantially stronger assertion belongs to the fibration discussion later in this chapter.
\end{dcomment}


\printedpage{25}
% head: A. Grothendieck : Dérivateurs, Ch. XII, Édition des Universités de Montpellier II et Paris VII
% foot: Version 7 décembre 2002
\textbf{Proposition 6.} \textit{Soit}
\[
\begin{tikzcd}[row sep=large, column sep=large]
X & X' \arrow[l, "q"'] \arrow[d, "f'"] \\
S \arrow[from=u, "f"'] & S', \arrow[l, "p"]
\end{tikzcd}
\]
\textit{cartésien dans} $\Cat$. \textit{Si $f$} \texttt{[est]} \textit{$W$-asphérique, $p$ $W$-lisse, alors $f'$ est $W$-asphérique.}

Il faut prouver \texttt{[que]} $X'\times_{S'}S'/s'$ est $W$-asphérique ($s'\in\Ob\,S'$). Soit $s=p(s')$, on a
\[
\begin{tikzcd}[column sep=small]
S & & S' \arrow[ll, "p"'] & \\
& S/s \arrow[ul] & & S'/s', \arrow[ul] \arrow[ll, "p_0"]
\end{tikzcd}
\]
avec $p_0$ $W$-lisse (lemme p.~28 \texttt{[plutôt 29]}) et à fibres $W$-asphériques, donc une $W$-équivalence universelle (prop.~5). Comme l'image inverse de $X$ sur $S/s$ est $W$-asphérique (hypothèse pour $f$), il \texttt{[s'ensuit]}

\leavevmode\msp{page 32}que dans le carré cartésien ci-contre \texttt{[ci-dessous]}
\[
(*)\qquad\qquad
\begin{tikzcd}[row sep=large, column sep=large]
X/s \arrow[d] & X'/s' \arrow[l, "W"'] \arrow[d] \\
S/s & S'/s' \arrow[l, "W"']
\end{tikzcd}
\]
les flèches horizontales sont dans $W$, et la première flèche verticale l'est par hypothèse sur $f$, donc la seconde aussi, donc $X'/s'$ \texttt{[est]} $W$-asphérique.

\dgrblock{dgrb485}{Why smooth base change preserves asphericity}
\begin{dcomment}{exposition}{p32:les flèches horizontales sont dans $W$, et la première flèche verticale}{Why smooth base change preserves asphericity}
Grothendieck's anticipated stability theorem now follows from the same categories of lifts, with the two variables interchanged. For a smooth $p:S'\to S$ and $s'\in S'$, put $s=p(s')$ and consider
\[
 p_{s'}:S'/s'\longrightarrow S/s.
\]
The dual of the coslice lemma makes $p_{s'}$ smooth. Its strict fibre over an arrow $v:t\to s$ is the category of incoming lifts of $v$ ending at $s'$. Smoothness of $p$ says that this category is aspherical. Proposition 5 therefore makes $p_{s'}$ a universal $W$-equivalence.

For an arbitrary $X\to S$, the square
\[
\begin{tikzcd}
 (X\times_S S')/s'\ar[r]\ar[d]&X/s\ar[d]\\
 S'/s'\ar[r,"p_{s'}"']&S/s
\end{tikzcd}
\]
is a strict pullback. An object of either description of its upper-left corner is a pair $(x,t')$ together with an arrow $t'\to s'$ and the equality $f(x)=p(t')$; the arrow conditions agree as well. Its top map belongs to $W$, independently of any assumption on $X$. If $X\to S$ is aspherical, $X/s$ is aspherical, so the top map makes the upper-left category aspherical. This proves Proposition 6.

For the corollary, apply this same universal equivalence once to $X$ and once to $Y$. If $f:X\to Y$ over $S$ satisfies $f/s\in W$, the commutative comparison square with horizontal arrows in $W$ gives $f'/s'\in W$. Thus smooth pullback preserves $W_S^{\mathrm g}$, the class detected on all over-categories. Opposite categories give the preservation of $W_S^{\mathrm d}$ by proper pullback. Neither argument needs the extra inference $W_S^{\mathrm g}\subseteq W$; accordingly the proof works with W(2 bis) alone.
\end{dcomment}


\textbf{Corollaire} (généralise \texttt{[la]} prop.~6, obtenue en faisant $Y=S$). \textit{Soit $p : S'\longrightarrow S$} \texttt{[un]} \textit{morphisme $W$-lisse dans} $\Cat$. \textit{Alors le foncteur changement de base}
\[
\Cat/S\longrightarrow\Cat/S'
\]
\textit{est compatible avec les localiseurs $W^{\mathrm{g}}_S$, $W^{\mathrm{g}}_{S'}$, i.e. transforme $f : X\longrightarrow Y$ qui est dans $W$} localement \textit{sur $S$, dans $f' : X'\longrightarrow Y'$ dans $W$ localement sur $S'$.}

Cette fois, on regarde le diagramme similaire à $(*)$
\[
\begin{tikzcd}[column sep=small]
X/s \arrow[dd] \arrow[dr, "f/s"] & & X'/s' \arrow[ll, "W"'] \arrow[dd] \arrow[dr, "f'/s'"] & \\
& Y/s \arrow[dl] & & Y'/s' \arrow[ll, "W"', crossing over] \arrow[dl] \\
S/s & & S'/s', \arrow[ll, "W"'] &
\end{tikzcd}
\]

\printedpage{26}
% head: A. Grothendieck : Dérivateurs, Ch. XII, Édition des Universités de Montpellier II et Paris VII
% foot: Version 7 décembre 2002
où les flèches horizontales sont dans $W$, donc $f/s\in W$ (hypothèse) implique $f'/s'\in W$.

\leavevmode\msp{page 33}Dualement, si $p$ est $W$-propre, le foncteur image inverse est compatible avec les localiseurs $W^{\mathrm{d}}_S$, $W^{\mathrm{d}}_{S'}$.

\dgrblock{dgrb486}{The passage from lifting categories to exact squares}
\begin{dcomment}{later}{p33:Dualement, si $p$ est $W$-propre}{The passage from lifting categories to exact squares}
The geometry just proved is the ingredient that later formulations package as exactness of squares. The term ``proper'' is justified by a universal base-change property; ``smooth'' describes the corresponding property when it is the map of bases that varies. The connection with the geometric proper and smooth base-change theorems, already discussed in the commentary on I.72, is made explicit in \citet[Introduction and Theorem 1.2]{Maltsiniotis2005Asphericite}.

The precise categorical statement appears in \citet[Propositions 2.12--2.13]{Maltsiniotis2012Exact}: a cartesian square with proper right side or smooth bottom side is $W$-exact, and the universal square condition characterizes properness. Its derivator interpretation is given by Theorem 4.32 of that paper, where cohomological base change is characterized through the appropriate weak exactness condition. These assertions explain why the lifting categories in the present text are sufficient; they are the categories appearing in pointwise base-change comparisons.

There is a convention to keep visible when comparing modern derivator texts. Grothendieck's coefficients here are contravariant diagrams. If a derivator is written with covariant diagrams instead, opposite categories exchange fibrations and opfibrations and interchange the two directions of the statement. In that convention \citet[Propositions 1.24 and 1.26]{Groth2013} proves base change for cartesian squares with a Grothendieck fibration on the right or an opfibration at the bottom, and identifies this property with the pointwise Kan-extension axiom. The apparent reversal has this variance explanation.

These developments do not identify every strict pullback with a homotopy pullback. An exact square controls a specified comparison of Kan extensions. Here the point is that properness or smoothness makes the relevant comma-category comparison cofinal in the required direction. The construction can therefore be transported to other coefficient theories without choosing path objects in each one separately.
\end{dcomment}


Si $f : X\longrightarrow S$ est $W$-propre, on va définir un foncteur canonique
\[
\begin{aligned}
\hot_*^{X/S,W} : S &\quad\longrightarrow\quad \Hot_W\overset{\text{déf}}{=}(\Cat)W^{-1}\\
\text{donné sur les objets par}\qquad s &\quad\longmapsto\quad \hot_W(X_s).
\end{aligned}
\]
Il faut définir, pour une flèche $u$ dans $S$, $u : s\longrightarrow s'$, une flèche
\[
\hot_W(X_s)\longrightarrow\hot_W(X_{s'})
\]
dans $\Hot_W$ - flèche que nous pouvons noter $\hot_*^{X/S,W}(u)$. On la définit par le diagramme
\[
\begin{tikzcd}[row sep=large, column sep=large]
X_s \arrow[r, dotted] \arrow[d, "i_s^X"] & X_{s'} \arrow[d, "i_{s'}^X"] \\
X/s \arrow[r, "X/u"] & X/s',
\end{tikzcd}
\]
où les flèches verticales sont $W$-coasphériques, donc dans $W$. Ce sont donc des isomorphismes dans $\Hot_W$, d'où l'unique flèche en pointillés dans $\Hot_W$, rendant le carré commutatif. Sur cette définition il est immédiat que
\[
u\longmapsto\hot_W(u)
\]
est compatible avec composition, et unités.

Ce foncteur est donc défini en supposant

\leavevmode\msp{page 34}seulement que les
\[
X_y\longrightarrow X/y
\]
sont dans $W$, ce qui est nettement moins fort que propre. Si on suppose que cette hypothèse est satisfaite `universellement' (i.e. reste satisfaite après tout changement de base $S'\longrightarrow S$) (\footnote[45]{On dira que $f$ est $W$-\textit{prépropre.}}), alors on a donc, pour tout $S'$ dans $\Cat/S$, un foncteur
\[
\hot_*^{X'/S',W} : S'\longrightarrow\Hot_W,
\]
et on vérifie tout de suite que ce n'est autre que le composé dans

\printedpage{27}
% head: A. Grothendieck : Dérivateurs, Ch. XII, Édition des Universités de Montpellier II et Paris VII
% foot: Version 7 décembre 2002
\[
(*)\qquad\qquad
\begin{tikzcd}[column sep=small]
S' \arrow[rr, "p"] \arrow[dr, "\hot_W^{X'/S'}"'] & & S \arrow[dl, "\hot_W^{X/S}"] \\
& \Hot_W &
\end{tikzcd}
\]
\texttt{[plutôt} $\hot_*^{X'/S',W}$ \texttt{resp.} $\hot_*^{X/S,W}$\texttt{]}. On a des observations duales si $X$ est lisse sur $S$, plus généralement si les inclusions
\[
X_s\longrightarrow s\backslash X
\]
sont dans $W$. On trouve alors un contrafoncteur
\[
\left\{
\begin{aligned}
S^\circ &\quad\longrightarrow\quad \Hot_W\\
s &\quad\longmapsto\quad \hot_W(X_s)\\
u &\quad\longmapsto\quad \hot^*_{X/S,W}(u)\ \text{ou}\ \hot^*_{X/S}(u)
\end{aligned}
\right.
\]

\leavevmode\msp{page 35}et une compatibilité similaire à $(*)$, si l'hypothèse est satisfaite \textit{universellement} sur $S$, p.~ex. $f$ $W$-lisse.

\dgrblock{dgrb487}{Transport obtained by conjugating the comma functor}
\begin{dcomment}{exposition}{p35:et une compatibilité similaire à}{Transport obtained by conjugating the comma functor}
To construct transport, Grothendieck uses a functor that already exists before localization. Let $Q:\Cat\to\Hot_W$ be the localization. The assignment
\[
 s\longmapsto X/s,\qquad
 (u:s\to t)\longmapsto(X/u):(x,\alpha)\longmapsto(x,u\alpha)
\]
is a strict functor $S\to\Cat$. If $i_s:X_s\to X/s$ belongs to $W$ for every $s$, define
\[
 \lambda(u)=Q(i_t)^{-1}Q(X/u)Q(i_s).
\]
The existence of the inverse is the only use of the hypothesis. For $s\xrightarrow u t\xrightarrow v r$, multiplication cancels $Q(i_t)$ against its inverse and gives
\[
 \lambda(v)\lambda(u)
 =Q(i_r)^{-1}Q(X/v)Q(X/u)Q(i_s)
 =\lambda(vu).
\]
The identity case follows in the same way. Thus the word ``functor'' is justified by the comma construction, not by a choice of lifts in the strict fibres.

For base change $q:S'\to S$, the canonical maps $X'/s'\to X/q(s')$ form a natural transformation of these comma diagrams, and the square with $X'_{s'}\cong X_{q(s')}$ and the two inclusions commutes. If both inclusions belong to $W$, two-out-of-three puts the comma comparison in $W$ and the displayed formula gives precisely the compatibility asserted in the text. Requiring the inclusion hypothesis after every base change is the definition of $W$-preproperness in footnote 45. Opposite categories give $W$-presmoothness, called $W$-prélissité in the source, and a contravariant transport functor.

Properness is a practical sufficient condition because it says more: each $i_s$ is coaspherical. Preproperness keeps the universal weak-equivalence conclusion and discards its local lifting explanation. Neither should be confused with the condition on the two endpoint inclusions over individual arrows considered on manuscript page 36. As the development below shows, those individual comparisons need not obey the displayed composition law.
\end{dcomment}


\textbf{Proposition 7.} \textit{Supposons $f : X\longrightarrow S$ à la fois $W$-propre et $W$-lisse, plus généralement que les morphismes}
\[
X_s\longrightarrow X/s,\qquad X_s\longrightarrow s\backslash X\qquad(s\in S)
\]
\textit{soient dans $W$, et que cette hypothèse reste verifiée après tout changement de base. Alors les foncteurs correspondants} % [sic] verifiée -> vérifiée, accent missing
\[
\begin{aligned}
\lambda &\;=\;\hot_*^{X/S,W} : &&\quad S \quad\longrightarrow\quad \Hot_W\\
\lambda^\circ &\;=\;\hot^*_{X/S,W} : &&\quad S^\circ \quad\longrightarrow\quad \Hot_W
\end{aligned}
\]
\textit{transforment toute flèche $u$ de $S$ en} \texttt{[un]} \textit{isomorphisme de} $\Hot_W$\textit{, et on a}
\[
\lambda^\circ(u)=\lambda(u)^{-1}\qquad\textit{dans}\quad\Hot_W,
\]
\textit{i.e. $\lambda(u)$ et $\lambda^c irc(u)$ sont inverses l'un de l'autre.} % [sic] λ^c irc(u) -> λ°(u), command \circ misprinted as superscript c followed by "irc"

\textsc{Démonstration}. La donnée de $u : s_0\longrightarrow s_1$ équivaut à celle de
\[
S'=\mathbf{\Delta}^1\longrightarrow S,
\]
d'où
\[
X'=X\times_S S'\qquad\text{sur } S'=\mathbf{\Delta}^1,
\]
avec les deux fibres
\[
X'_0\simeq X_{s_0},\qquad X'_1\simeq X_{s_1}.
\]
L'hypothèse assurant l'existence de $\lambda$ s'exprime par la condition que
\[
X'_1\longrightarrow X'/1\simeq X'
\]
\printedpage{28}
% head: A. Grothendieck : Dérivateurs, Ch. XII, Édition des Universités de Montpellier II et Paris VII
% foot: Version 7 décembre 2002
\msp{page 36}est dans $W$, de sorte que le diagramme
\[
X'_0\xrightarrow{i_0}X'\xleftarrow[W]{i_1}X'_1
\]
définit une flèche $\lambda(u) : X'_0\longrightarrow X'_1$ dans $\Hot_W$, savoir $\lambda(u)=i_1^{-1}i_0$. L'hypothèse qui assure l'existence de $\lambda^\circ$ s'exprime par la condition symétrique $i_0\in W$, qui permet de définir la flèche $\lambda^\circ(u) : X'_1\longrightarrow X'_0$ par $\lambda^\circ(u)=i_0^{-1}i_1$. Quand les deux conditions sont satisfaites simultanément, on voit donc que $\lambda(u)$, $\lambda^\circ(u)$ sont inverses l'un de l'autre.

\dgrblock{dgrb488}{Two directions of transport give monodromy}
\begin{dcomment}{examples}{p36:Quand les deux conditions sont satisfaites simultanément}{Two directions of transport give monodromy}
Proposition 7 says that the two canonical transports become inverse when both comma comparison hypotheses hold universally. Pulling back over $u:[1]\to S$ gives $E=X\times_S[1]$ and inclusions $i_0:X_{s_0}\to E$, $i_1:X_{s_1}\to E$. Preproperness makes $i_1$ a $W$-equivalence, and presmoothness makes $i_0$ one. The already constructed functors therefore satisfy
\[
 \lambda(u)=Q(i_1)^{-1}Q(i_0),\qquad
 \lambda^{\mathrm{op}}(u)=Q(i_0)^{-1}Q(i_1).
\]
These formulas prove the claimed inverse relation without an additional saturation assumption. The covariant functor factors through the ordinary groupoid completion $S[\mathrm{Mor}(S)^{-1}]$. It records transport around loops, or monodromy, at the level of the homotopy category.

For a concrete example, let $C_2=\{1,\sigma\}$ act on the discrete set $F=\{0,1\}$ by interchange. Regard $BC_2$ as the one-object groupoid. The action groupoid $E$ has objects $0,1$ and an arrow labelled $g$ from $x$ to $gx$. The projection $E\to BC_2$ is both a fibration and an opfibration: every base arrow has a unique incoming and a unique outgoing lift once the appropriate endpoint is specified. Its strict fibre is $F$. The transport along $\sigma$ interchanges the two points, in both directions, since $\sigma^{-1}=\sigma$.

The action is free and transitive, so there is exactly one arrow between each ordered pair of objects of $E$. Either object is consequently both initial and terminal, and $BE$ is contractible. The base has fundamental group $C_2$, while the fibre consists of two points. This is the usual universal two-sheeted covering expressed entirely by categories. It illustrates why the fibre functor and its monodromy carry information that the homotopy type of the total category alone does not retain. It also illustrates the scope of Proposition 7: it gives a functor into $\Hot_W$, not a claim that ordinary homotopy-category monodromy classifies all homotopy-coherent families.
\end{dcomment}


\textbf{Remarques.} Pour définir $\lambda$ `universellement', il n'était pas indispensable de considérer les $X/s$ et les inclusions $X_s\longrightarrow X/s$, au lieu de cela il suffisait de regarder les changements de base $S'=\mathbf{\Delta}^1\xrightarrow{u}S$, et de faire l'hypothèse que pour tout tel changement de base, l'inclusion $X'_1\xrightarrow{i_1}X'=X_{(\mathbf{\Delta}^1,u)}$ est dans $W$. Cette hypothèse est automatiquement stable par changement de base arbitraire. (La $W$-propreté de $f : X\longrightarrow S$ signifie que $X'_1\longrightarrow X'$ est non seulement dans $W$, mais même $W$-coasphérique.) Dualement, pour la définition `universelle' de $\lambda^\circ$, il suffirait que pour tout changement de base $\mathbf{\Delta}^1\longrightarrow S$ comme dessus, l'inclusion $X'_0\xrightarrow{i_0}X'$ est dans $W$.

\dgrblock{dgrb489}{An eight-object obstruction to edge transport}
\begin{dcomment}{development}{p36:Cette hypothèse est automatiquement stable par changement de base arbitraire.}{An eight-object obstruction to edge transport}
Grothendieck proposes to replace the universal comma comparisons by endpoint comparisons over the walking arrow. Those tests do define a morphism for each arrow of the base, but they do not by themselves make these morphisms functorial. The distinction is already visible on connected components, and also answers the question about pseudofibrations on manuscript page 40.

\textbf{Proposition (edge tests do not force coherence).}\ours
There exists a functor $p:P\to[2]$ such that, for every $u:[1]\to[2]$, both endpoint inclusions into $P\times_{[2]}[1]$ belong to $W_\infty$, but the resulting edge transports fail the composition law. Moreover, $p$ does not satisfy the fibration test over contractible bases.

\emph{Proof.} Indices will be taken modulo $2$. Let $P$ have eight objects
\[
 a_0,a_1,b_0,b_1,d_0,d_1,c_0,c_1,
\]
and exactly the following nonidentity order relations:
\[
 a_i<d_i,\qquad b_i<d_i,\qquad
 b_i<c_i,\qquad a_i<c_{1-i}\quad(i=0,1).
\]
Every relation starts at an $a_i$ or a $b_i$ and ends at a $d_i$ or a $c_i$. Thus no two strict relations are composable, and the list already defines a transitive partial order. Define $p(a_i)=0$, $p(b_i)=p(d_i)=1$, and $p(c_i)=2$. All listed relations preserve this order, so $p$ is a functor.

The fibre $P_0$ is the discrete pair $\{a_0,a_1\}$, the fibre $P_1$ is the disjoint union of the two intervals $b_i<d_i$, and $P_2$ is the discrete pair $\{c_0,c_1\}$. For the nonidentity base arrow $0\to1$, the pullback is the disjoint union of the two posets $a_i<d_i>b_i$. Each component has terminal object $d_i$, and each endpoint fibre meets it in a nonempty contractible subcategory. Both endpoint inclusions are therefore weak homotopy equivalences. For $1\to2$, the two components are $d_i>b_i<c_i$; each has initial object $b_i$, and the same conclusion follows. For $0\to2$, the pullback consists of the two intervals $a_i<c_{1-i}$, so again both endpoint inclusions are weak equivalences. A constant functor $[1]\to[2]$ has pullback $P_j\times[1]$; its endpoint inclusions are categorical homotopy equivalences. These six functors exhaust all $[1]\to[2]$, proving every required edge test.

Label the two connected components of each fibre by $i$. The span through the $0,1$ pullback identifies component $i$ with component $i$; so does the span through the $1,2$ pullback. The $0,2$ pullback instead identifies $i$ with $1-i$. Thus the maps induced on $\pi_0$ satisfy
\[
 \pi_0\lambda(0\to1)=1,\qquad
 \pi_0\lambda(1\to2)=1,\qquad
 \pi_0\lambda(0\to2)=(0\ 1).
\]
Since $\pi_0$ sends $W_\infty$ to bijections, it factors through $\Hot$, and these unequal functions prove that the proposed transports cannot form a functor into $\Hot$.

For completeness, the nerve of $P$ has eight vertices, eight nondegenerate edges, and no higher nondegenerate simplices. Those edges form exactly the cycle
\[
 a_0-d_0-b_0-c_0-a_1-d_1-b_1-c_1-a_0.
\]
Its realization is therefore a circle. Each fibre has two connected components, whereas $P$ is connected. Consequently none of the inclusions $P_j\to P$ is in $W_\infty$, although the base $[2]$ has both an initial and a terminal object. This disproves the stronger fibration condition over contractible bases and completes the proof.

The counterexample also satisfies the edge tests for $W_0$, the localizer detected by connected components, and has the same failure of composition there. The issue is therefore not a higher-dimensional obstruction inside the fibres. The missing information is how transports for composable arrows agree across a triangle in the base. The universal comma hypotheses used in Proposition 7 supply precisely the strict diagram from which that agreement was obtained. We checked the lift-category and smoothness treatments in \citet[\S4]{Maltsiniotis2005Asphericite} and \citet[\S\S1--2]{Maltsiniotis2012Exact}; we did not find this finite example there. This is a development of the manuscript's question, with no claim of priority.
\end{dcomment}


\leavevmode\msp{page 37}La notion la plus faible raisonnable de $W$-fibré, me semble être celle où on suppose les deux choses à la fois. Les hypothèses qu'on a faites jusqu'à présent, à l'exclusion de la forme forte de $W$-propreté ou de $W$-lissité, reviennent toujours à supposer que le \textit{changement de base par $f : X\longrightarrow S$} transforme certaines flèches dans $\Cat/S$ en \texttt{[élément de]} $W$ : Pour définir $\lambda$, ces flèches sont
\[
\text{soit } \{s\}\longrightarrow S/s,\qquad \text{soit } \{1\}\longrightarrow\mathbf{\Delta}^1(\longrightarrow S),
\]
pour définir $\lambda^\circ$, ce sont
\[
\text{soit } \{s\}\longrightarrow s\backslash S,\qquad \text{soit } \{0\}\longrightarrow\mathbf{\Delta}^1(\longrightarrow S),
\]
le `soit' suivant la définition qu'on a adoptée. Dans ces exemples, il s'agit toujours de l'inclusion de l'objet final d'une catégorie $S'$, dans ladite, ou inversement de l'objet initial. Cela montre d'ailleurs (si on exige l'universalité), que la première hypothèse (assurant $\lambda$) est \textit{exactement} cela (cas de l'inclusion d'un objet final $s'$ - puisque pour un tel objet, on a $S'/s'\simeq S'$), tandis que la deuxième hypothèse exige moins, seulement pour le cas de $\mathbf{\Delta}^1$. Et dualement pour $\lambda'$ \texttt{[plutôt $\lambda^\circ$]}. Mais il convient de noter que pratiquement, la façon la plus commode, et de loin,

\leavevmode\msp{page 38}pour vérifier de telles propriétés, c'est soit par la $W$-propreté de $f$, soit par la $W$-lissité - ce qui ramène à vérifier que certaines catégories bien précisées sont $W$-asphériques, sans avoir à se soucier de changements de base et de vérifier que les flèches construites en termes de changements de base arbitraires sont dans $W$. Et l'emploi de la $W$-propreté et de la $W$-lissité est d'autant plus sympathique, que ces notions sont stables par composition,

% ---- from Der-ChXIIa-012_p031-035.tex
\printedpage{29}
% head: A. Grothendieck : Dérivateurs, Ch. XII, Édition des Universités de Montpellier II et Paris VII
alors que toutes les notions dérivées, ainsi que les notions de fibrations à la clef, ont une fâcheuse tendance à ne \textit{pas} être stable par composition.

On pourrait appeler pseudo-$W$-propre resp. quasi-$W$-propre (un peu mieux que pseudo) ce qu'on vient de dégager - j'entends des ricanements! Et si on veut les deux à la fois - $\lambda$ et $\lambda^{\circ}$ associés comme on a dit? Je vois le même principe de description avec des $S''\fl S'$ sur $S$, d'exiger que $X''\fl X'$ soit transformé en des \texttt{[éléments de]} $W$. Avec les choix suivants qui ont l'air sympathique

\begin{itemize}
\item[a$_1$)] \msp{page 39}$Y'$ a un objet final, $Y''=e$. La définition revient donc à la condition que les
\[ X_s\fl X/s, \qquad X/s'\fl X/s \]
(pour $s\in\Ob S$ où $u : s'\fl s$ dans $\mathrm{S}$) soient dans $W$, et que ça reste comme ça après tout changement de base $S'\fl S$.

\item[a$_2$)] $Y'$ a un objet initial, $Y''=e$, c'est la condition duale à coups de
\[ X_s\fl s\backslash X, \qquad s\backslash X\fl s'\backslash X. \]

\item[a$_0$)] $Y'=\boldsymbol{\Delta}^1$, $Y''=e$ (\footnote[46]{NB a$_1$), a$_2$) sont notés pareil p. 4, mais a$_0$) a été oublié.}).

\item[b)] $Y'$ est contractile, $Y''$ aussi (ou au choix, $Y'$ contractile, $Y''=e$)
\end{itemize}

donc la condition revient à ceci : pour toute $Y'$ sur $Y$ \textit{contractile} (notion indépendante de $W$, c'est ça son principal avantage) les inclusions des fibres de $X'=X\times_Y Y'$ sur $Y'$ dans $X'$ sont \texttt{[dans]} $W$.

\dgrblock{dgrb490}{The test objects explain the hierarchy}
\begin{dcomment}{exposition}{p39:donc la condition revient à ceci}{The test objects explain the hierarchy}
The proposed definitions differ by the category that is allowed to serve as a test base. For condition $a_1$, a base $T$ with terminal object $t$ is tested at \emph{every} object $r\in T$, not only at $t$. Testing the terminal object gives preproperness. Once it is available, the unique arrow $r\to t$ identifies the comparison $X_r\to X_T$ with transport to the terminal fibre. Requiring that comparison to belong to $W$ adds the demand that all such transports be weak equivalences. The dual reading gives $a_2$.

Condition $a_0$ restricts the base to $[1]$ and allows either endpoint. The preceding eight-object example shows that it cannot be substituted for $a_1$ or $a_2$, even for ordinary weak homotopy equivalences. In particular, existence of individually invertible edge roofs is not a substitute for a functor of fibres. This is the reason for keeping the preproper or presmooth hypothesis in the equivalences proved later in the section.

For condition $b$, the equivalence of the two tests stated in the text has a short proof. Suppose all fibre inclusions are in $W$ over every categorically contractible test base. If $v:T''\to T'$ is a map between two such bases, choose an object $t''$; contractibility ensures that an object exists. In the pullback square, the fibre at $t''$ identifies with the fibre at $v(t'')$. The two inclusions from that common fibre into the total categories both belong to $W$, so their commutative triangle puts the total pullback map in $W$ by two-out-of-three. The converse is the special case $T''=e$.

Condition $b'$ allows $W$-aspherical bases instead. Condition $c$ allows every map in $W$ as a change of base. Its stronger quantifier is substantial: it asks for preservation of a whole class of weak equivalences, whereas the earlier conditions test selected contractible shapes. The terminology at this point is provisional, but these quantifiers are exact and should be retained when comparing later definitions.
\end{dcomment}


Je sens que même pour des $W$ générales, b) reste du domaine du vérifiable, tandis qu'on mollit si on passe à l'extrême en généralité dans ce sens :

\begin{itemize}
\item[b$'$)] Comme b), mais en supposant seulement $Y'$ $W$-asphérique (\footnote[47]{C'est essentiellement le a$'$) de p. 4}).
\end{itemize}

Enfin (la condition a) de la page 4)

\begin{itemize}
\item[c)] Si $Y''\fl Y'$ est \texttt{[dans]} $W$ sans plus, $X''\fl X'$ aussi.
\end{itemize}
% foot: Version 7 décembre 2002

\printedpage{30}
% head: A. Grothendieck : Dérivateurs, Ch. XII, Édition des Universités de Montpellier II et Paris VII
\msp{page 40}La seule de ces notions qui ait la moindre envie d'être stable par composition, c'est la plus forte c). Par contre, la notion de flèche $W$-parfaite (cf. déf. 3) est stable spontanément par composition :

\textbf{Définition 3.} On dit que $f : X\fl S$ est $W$-\textit{parfaite} si elle est à la fois $W$-propre et $W$-lisse.

\dgrblock{dgrb491}{The interval test and the full fibration condition}
\begin{dcomment}{exposition}{p40:à la fois $W$-propre et $W$-lisse.}{The interval test and the full fibration condition}
Grothendieck is trying to retain the mobility of fibres without asking that every weak equivalence survive every pullback. The latter requirement is exceptionally strong when $W$ forgets some homotopy groups. The useful intermediate condition asks instead that a categorically contractible base make the family constant up to $W$. Here ``contractible'' means that its identity functor is joined to a constant functor by a finite zigzag of natural transformations. This is the notion of homotopy already used in section~I; it must not silently be replaced by weak contractibility of the nerve.

The weakest test in the diagram examines only $[1]\to S$. It gives, for each arrow $u:s\to t$, a span
\[
 X_s\longrightarrow X\times_S[1]\longleftarrow X_t
\]
whose arrows belong to $W$. Inverting them produces an isomorphism between the two fibres in $\Hot_W$. What it does not supply is the compatibility of this isomorphism with composition of arrows of $S$. The development ``An eight-object obstruction to edge transport'' above gives a finite poset over $[2]$ for which all these endpoint tests are weak equivalences, whereas the transports along $0\to1\to2$ and $0\to2$ disagree already on connected components. Its total category has the homotopy type of a circle and each fibre has two components. It therefore answers negatively the question about $a_0$ raised in the next paragraph.

Proposition~8 adds precisely the missing global control. Universal preproperness identifies a fibre with the appropriate comma category; the comma categories carry genuinely functorial transport, because postcomposition in $S$ is an actual functor. This explains why preproperness together with the interval test succeeds where the interval test alone fails. Perfectness is a convenient sufficient condition, since it supplies both universal comma comparisons. It is stronger than the fibration condition being isolated here.
\end{dcomment}


Dans le cas $W=W_\infty$, toutes les variantes de la page précédente, à l'exclusion seulement de a$_0$), sont équivalentes. Ce serait intéressant de dégager des conditions minimales sur $W$ pour assurer qu'il en soit ainsi. (Et il se pose la question si dans le cas de $W_\infty$, a$_0$) suffit pour impliquer le reste - je ne m'y attends guère \dots)

Dans le cas général, on a le diagramme d'implications (et encore je me suis limité)
\[
\begin{tikzcd}[column sep=small]
 & & \mathrm{c} \arrow[d, Rightarrow] & & \text{$W$-fibration ultraforte} \\
 & & \mathrm{b}' \arrow[dd, Rightarrow] & & \text{$W$-fibration $W$-forte} \\
 & \text{$W$-parfait} \arrow[dr, Rightarrow] & & & \\
 & & \mathrm{b} \arrow[dl, Rightarrow] \arrow[dr, Rightarrow] & & \text{$W$-fibration hot-forte (\footnotemark[48])} \\
\text{$W$-fibration gauche} & \mathrm{a}_1 \arrow[dr, Rightarrow] & & \mathrm{a}_2 \arrow[dl, Rightarrow] & \text{$W$-fibration droite} \\
 & & \mathrm{a}_0 & & \text{$W$-pseudofibration}
\end{tikzcd}
\]
\footnotetext[48]{Il semblerait que ce soit ça, finalement, la ‘bonne notion’ de $W$-fibration, \textit{équivalente} à a$_1$ et a$_2$, lesquelles seraient équivalentes entre elles. Cf. prop. 9, p. 48.}

\texttt{[‘$W$-parfait $\Longrightarrow$ b’ encerclé]}.

\msp{page 41}\textbf{Proposition 8.} \textit{Soit $f : X\fl S$ dans} Cat. \textit{Les conditions suivantes sont équivalentes.}

\begin{itemize}
\item[\textit{a)}] \textit{$f$ est une $W$-fibration gauche, i.e. la condition} a$_1$ \textit{de la page 39 est vérifiée.}

\item[\textit{b)}] \textit{$f$ est $W$-prépropre, et $\forall\; s_0\fl s_1$ dans $S$, $X/s_0\fl X/s_1$ est dans $W$.}

\item[\textit{c)}] \textit{$f$ est $W$-prépropre, et une $W$-pseudofibration, i.e. la condition a$_0$ de la page 39 est vérifiée.} % AUDIT resolved: CMTI12 text-italic a with CMR8 subscript; see corrections.csv, PDF leaf 40, ms41
\end{itemize}

(\footnote[49]{\textbf{Corollaire.} \textit{Quand $f$ est} Cat-\textit{cofibrant, alors ces conditions équivalent aussi à ceci : $\forall\; u : s_0\fl s_1$ dans $S$, $u_* : X_{s_0}\fl X_{s_1}$ est dans $W$.}}).
% foot: Version 7 décembre 2002

\printedpage{31}
% head: A. Grothendieck : Dérivateurs, Ch. XII, Édition des Universités de Montpellier II et Paris VII
\textbf{NB} Ces conditions sont satisfaites si $f$ est en même temps $W$-prélisse, cf. prop. 7.

\dgrblock{dgrb492}{Completing the terminal-object test}
\begin{dcomment}{exposition}{p41:Ces conditions sont satisfaites si}{Completing the terminal-object test}
Grothendieck wants a condition that can be checked on the original base, while remaining valid after every change of base. We spell out the argument because the compressed reference to $Y/s$ and $Y/s'$ hides the universal quantifier.

Write $E_u=X\times_S[1]$ for $u:s_0\to s_1$. If the terminal-object test holds, apply it first to $T/t\to T\to S$, for an arbitrary $T\to S$ and $t\in T$. Its terminal fibre is $X'_t$ and its total category is $X'/t$, where $X'=X\times_ST$. Thus $X'_t\to X'/t$ belongs to $W$: this is universal preproperness. Next apply the same test to $S/s_1$. Both the object $(s_1,1)$ and the object $(s_0,u)$ give fibre inclusions into $X/s_1$ in $W$. Since the latter factors as
\[
 X_{s_0}\longrightarrow X/s_0\longrightarrow X/s_1,
\]
two-out-of-three gives $X/s_0\to X/s_1\in W$. This supplies the missing detail in $a)\Rightarrow b)$.

Assume now universal preproperness. The canonical functor $E_u\to X/s_1$ restricts on the terminal fibre to $X_{s_1}\to X/s_1$. Both this inclusion and $X_{s_1}\to E_u$ belong to $W$, so $E_u\to X/s_1$ belongs to $W$. Its restriction to the initial fibre is the composite displayed above. Consequently
\[
 X_{s_0}\to E_u\in W
 \quad\Longleftrightarrow\quad
 X/s_0\to X/s_1\in W.
\]
The other endpoint inclusion is already in $W$. This proves $b)\Longleftrightarrow c)$, including the asserted stability under arbitrary base change, since preproperness and the interval condition both have that stability by definition.

Finally, after any base change to a category $T$ with terminal object $t$, the inclusion of the fibre over $s\in T$ factors through $X'/s\to X'/t=X'$. Both factors belong to $W$, proving $a)$. The dual argument uses initial objects and the categories $s\backslash X$. For an ordinary opfibration the inclusions $X_s\to X/s$ have adjoints and are homotopy equivalences; its postcomposition maps identify, up to these comparisons, with the pushforward functors $u_*$. This proves the corollary in footnote~49 with its stated variance.
\end{dcomment}


\textsc{Démonstration}. On a a) $\Longrightarrow$ b) trivialement (faire $Y''=Y/s$, $Y''=Y/s'$ \texttt{[lien avec l'assertion? Cf. p. 4.]}).

b) $\Longleftrightarrow$ c) par le raisonnement de la proposition 7, car en termes de $X'=X_u=X\times_S(\boldsymbol{\Delta}^1,u)$, la condition b) pour $u : s_0\fl s_1$ donné s'interprète dans le diagramme commutatif
\[
\begin{tikzcd}[column sep=small, row sep=small]
X'_0 \arrow[dr, "i'_0\in W"] \arrow[dd, "\wr"] & & X'_1 \arrow[dr, "i'_1\in W"] \arrow[dd, "\wr"] & \\
 & \underset{(\simeq X'_0)}{X'/0} \arrow[rr, "\varphi'", crossing over] \arrow[dd, "p_0"] & & X'/1\ (\simeq X') \arrow[dd, "p_1"] \\
X_{s_0} \arrow[dr, "i_0\in W"'] & & X_{s_1} \arrow[dr, "i_1\in W"] & \\
 & X/s_0 \arrow[rr, "\varphi"] & & X/s_1
\end{tikzcd}
\]

\texttt{[il y a quelques $W$ redondants dans le diagramme original]} où $i_0$, $i_1$, $i'_0$, $i'_1$ sont dans $W$ à cause de l'hypothèse de $W$-prépropreté (stable par changement de base).

Donc $p_0$, $p_1$ sont \texttt{[dans]} $W$, donc $\varphi\in W$ si et seulement si $\varphi'\in W$, or $\varphi'\in W$ est la condition a$_0$.

\msp{page 42}Prouvons enfin que b), c) impliquent a). Comme c) (donc aussi b)) est stable par changement de base, il suffit de montrer ceci : Si $S$ a un objet final $e$, alors pour tout $s\in\Ob S$, l'inclusion $X_s\xrightarrow{\ i\ }X$ est dans $W$. Or on a le diagramme commutatif
\[
\begin{tikzcd}
X_s \arrow[r, "i_s\in W"] \arrow[dr, "i"'] & X/s \arrow[d, "u\,\in\,W\text{ par c) }\texttt{[plutôt par b)]}"] \\
 & X/e=X,
\end{tikzcd}
\]
qui montre bien $i=ui_s\in W$.

Dualement, si $f$ est $W$-prélisse, alors la condition a$_0$ ($f$ une $W$-pseudofibration) est équivalente à a$_2$ ($f$ une $W$-fibration droite), et signifie aussi que les
\[ s_1\backslash X\fl s_0\backslash X \]
sont dans $W$. Mais la notion de ‘$W$-fibration’ sur laquelle on tombe (si on ne se contente du minable a$_0$!) est différente dans les deux cas duaux l'un de l'autre. On aimerait dans les deux cas trouver des $W$-fibrations ‘hot-fortes’ (condition b), tout au moins en mettant le paquet (s'il le faut) - i.e. pour $f$ $W$-parfait. Ça revient donc à ceci.

\textbf{Question :} \textit{Soient $X\xrightarrow{\ f\ }S$, avec $f$ $W$-parfait, $S$ contractile, $s\in\Ob S$. L'inclusion $X_s\fl X$, est-elle dans $W$?}
% foot: Version 7 décembre 2002

\printedpage{32}
% head: A. Grothendieck : Dérivateurs, Ch. XII, Édition des Universités de Montpellier II et Paris VII
\msp{page 43}On va supposer seulement que $f$ est $W$-lisse et une $W$-fibration gauche \texttt{[plutôt que $f$ est simplement une $W$-fibration à droite, i.e que $f$ est $W$-prélisse et une $W$-pseudofibration]} (cf. prop. 8). On va se ramener au cas où $X$ est même fibrant sur $S$, en introduisant $\Phi_S(X)$ \texttt{[catégorie fibrée définie par le foncteur $S^{\circ}\fl\Cat$, $s\longmapsto s\backslash X$]} et % [sic] i.e que -> i.e. que, missing period of the abbreviation
\[ X\xrightarrow{\ i_X\ }\Phi_S(X). \]

\dgrblock{dgrb493}{The split replacement and its actual comparison maps}
\begin{dcomment}{exposition}{p43:On va se ramener au cas}{The split replacement and its actual comparison maps}
Grothendieck now replaces a family by a split fibration so that transport can be written before localization. The replacement is legitimate for a more concrete reason than an appeal to a previously constructed derivator.

An object of $\Phi_S(X)$ is a triple $(s,x,a)$ with $a:s\to p(x)$. A morphism to $(t,y,b)$ is a pair $(u:s\to t,v:x\to y)$ satisfying $p(v)a=bu$. The structural projection sends this object to $s$. Its fibre at $s$ is $s\backslash X$, and an arrow $u:s\to t$ acts by precomposition. This is a strict functor $S^{\mathrm{op}}\to\Cat$, hence a split Grothendieck fibration.

The comparison $i_X$ sends $x$ to $(p(x),x,1)$. The forgetful functor $r:\Phi_S(X)\to X$ satisfies $ri_X=1_X$, and the arrows $(a,1_x)$ define a natural transformation $1_{\Phi_S(X)}\Rightarrow i_Xr$. In fact $r$ is left adjoint to $i_X$. Thus $i_X$ is a categorical homotopy equivalence and belongs to every $W$ under consideration. On the fibre at $s$ it is exactly
\[
 X_s\longrightarrow s\backslash X.
\]
For a universally prelisse $X/S$, these fibre comparisons belong to $W$ after every base change. The comparison of total categories and the comparison of individual fibres are therefore two separate, verified assertions.

If $X/S$ is also a pseudofibration, the dual of Proposition~8 says that every precomposition $t\backslash X\to s\backslash X$ is in $W$. Consequently $\Phi_S(X)$ is a split fibration with transport in $W$. In a square comparing a fibre inclusion for $X$ with the corresponding inclusion for $\Phi_S(X)$, both horizontal arrows are in $W$; two-out-of-three transfers the desired conclusion. No saturation argument is needed for this step. The familiar free-fibration construction and its variance were discussed in the commentary on Chapter~VI, manuscript pp.~53--60; the additional point here is the universal fibre comparison that permits its use over an arbitrary new base.
\end{dcomment}


L'ennui, c'est que j'ai fait la théorie de $\HOT_W$ jusqu'à maintenant avec l'hypothèse (qui s'avère idiote) que $W\supseteq W_\infty$ - je serai donc obligé de revoir rapidement sous quelles hypothèses plus raisonnables sur $W$ les résultats familiers sont valables. Par exemple j'aimerais que $i_X\in W$ (\footnote[50]{en fait, $i_X\in(W_\omega)^{\mathrm{d}}_S\subseteq W^{\mathrm{d}}_S\subseteq W_S$.}), de sorte qu'on ait
\[
\begin{tikzcd}[column sep=large, row sep=large]
X_s \arrow[r, "i_s\in W", "\text{car $X$ $W$-lisse}"'] \arrow[d, "j^X_s"'] & \Phi_S(X)_s=s\backslash X \arrow[d, "j_s^{\Phi_S(X)}"] \\
X \arrow[r, "i_X\in W"] & \Phi_S(X),
\end{tikzcd}
\]
on voit donc que pour prouver que l'inclusion $j^X_s$ est \texttt{[dans]} $W$, il suffit de prouver la même chose pour $\Phi_S(X)$, qui est fibrant (pas seulement $W$-lisse). De plus, l'hypothèse pseudo-fibration sur $X$ passe sûrement à $\Phi_S(X)$ (je réserve les deux vérifications à plus tard (\footnote[51]{O.K.})). On est donc ramené au cas de
\[ X\in\Ob\uFib_0S, \qquad S\text{ contractile,} \]
à prouver que l'inclusion
\[ X_s\fl X \]
\msp{page 44}est dans $W$. Je vais prouver ceci :

\textbf{Lemme.} \textit{Soient $S$ dans} Cat, \textit{$X\in\Ob\uFib_0S$, et considérons, pour un $S'$ dans} Cat,
\[
\begin{aligned}
f &\quad\longmapsto\quad f^*(X/S)=X\times_S(S',f)\\
\uHom(S',S)^{\circ} &\quad\longrightarrow\quad \uFib_0S'.
\end{aligned}
\]
\textit{(Ce foncteur est défini grâce au fait que $X$ est fibrant sur $S$, et peut s'interpréter donc, à $S$-équivalence près, en termes d'un foncteur $S^{\circ}\fl$} Cat \textit{\dots\ (\footnote[52]{\textbf{NB} $\Phi_S(X)$ était même déjà scindée, définie par le foncteur $s\longmapsto s\backslash X$.})). Les flèches de $\uHom(S',S)$ sont transformées en flèches $\in W^{\mathrm{d}}_{S'}$.}
% foot: Version 7 décembre 2002

\printedpage{33}
% head: A. Grothendieck : Dérivateurs, Ch. XII, Édition des Universités de Montpellier II et Paris VII
En effet, en termes de foncteurs
\[ S'^{\circ}\fl\Cat, \]
si $X$ est définie par
\[ S^{\circ}\xrightarrow{\ \mathscr{X}\ }\Cat, \]
alors $f^*(X)$ est définie par $\mathscr{X}\circ f$, et la flèche
\[ \mathscr{X}\circ f\fl\mathscr{X}\circ g \]
définie par $u : f\fl g$ n'est autre que $\mathscr{X}*u$. Sa valeur en $s'\in\Ob S'$ est obtenue en prenant
\[ u(s') : f(s')\fl g(s'), \]
d'où par $\mathscr{X}$
\[ \mathscr{X}(u(s')) : \underbrace{\mathscr{X}(g(s'))}_{(\mathscr{X}\circ g)(s')}\fl\underbrace{\mathscr{X}(f(s'))}_{(\mathscr{X}\circ f)(s')} \]
(\footnote[53]{à dégager un énoncé à part.}). Mais l'hypothèse $X\in\uFib_0S$ signifie que les $\mathscr{X}(v)$, $v\in\mathrm{Fl}(S)$, sont tous dans $W$. Donc
\msp{page 45}$\mathscr{X}$ est dans $W$ fibre par fibre, donc dans $W^{\mathrm{d}}_{S'}$.

\dgrblock{dgrb494}{Transport along a natural transformation}
\begin{dcomment}{exposition}{p45:est dans $W$ fibre par fibre}{Transport along a natural transformation}
The purpose of the lemma is to turn a homotopy of indexing functors into an actual zigzag of families. For this purpose it is useful to work with the split fibration just constructed. Let $F:S^{\mathrm{op}}\to\Cat$ be its classifying functor, with every $F(u)$ in $W$, and let $\alpha:f\Rightarrow g$ be a natural transformation of functors $T\to S$.

Contravariance gives the strict natural transformation
\[
 F(g(-))\longrightarrow F(f(-)),\qquad
 t\longmapsto F(\alpha_t).
\]
Its Grothendieck construction is a Cartesian functor $g^*X\to f^*X$ over $T$. Each component belongs to $W$. The relative localization axiom, applied in the dual comma convention, implies that this total functor belongs to $W$ as well. Explicitly, the inclusion of the fibre of a Grothendieck fibration into its under-comma category has an adjoint, so the componentwise assertion gives an equivalence on those comma categories; the dual of W(2) then gives the total equivalence.

The first displayed arrow in the source proof points in the opposite direction from the correctly written component immediately below it. It is an edition slip: the transformation is $F\circ g^{\mathrm{op}}\Rightarrow F\circ f^{\mathrm{op}}$. Composition is nevertheless completely controlled: for $\beta:g\Rightarrow h$, the identity
\[
 F(\beta_t\alpha_t)=F(\alpha_t)F(\beta_t)
\]
identifies the transport for $\beta\alpha$ with the composite of the two transports, in the contravariant order. Units are identities. After localization every such transport is invertible, and one may use its inverse when writing the covariant version of the functor in the next paragraph.

The factorization through the groupoid completion of $\underline{\mathrm{Hom}}(T,S)$ follows from its universal property. Connected functors therefore induce isomorphic pullbacks, but a connected component does not choose a particular path, and different paths may induce different isomorphisms. This is the precise reason for the parenthetical qualification ``non canoniquement''.
\end{dcomment}


\textbf{Corollaire 1.} \textit{Le foncteur du lemme, composé avec}
\[ \uFib_0S'\fl\HOT^{\mathrm{lc}}_W(S')\overset{\text{déf}}{=}(\uFib_0S')(W^{\mathrm{d}}_{S'})^{-1} \]
\textit{se factorise par le groupoïde fondamental de} $\uHom$ \textit{:}
\[
\begin{tikzcd}[row sep=huge]
\uHom(S',S)^{\circ} \arrow[r] \arrow[d] & \uFib_0S' \arrow[d] \\
\textstyle\prod(\uHom(S',S))^{\circ} \arrow[r] & \HOT^{\mathrm{lc}}_W(S').
\end{tikzcd}
\]
\textit{Si $f$, $g$ sont dans une même composante connexe de $\uHom(S',S)$, alors $f^*X$ et $g^*X$ sont isomorphes (non canoniquement, en général) dans} $\HOT^{\mathrm{lc}}_W(S')$.

\dgrblock{dgrb495}{A loop can act on the fibre even when the functor is fixed}
\begin{dcomment}{examples}{p45:isomorphes (non canoniquement, en général) dans}{A loop can act on the fibre even when the functor is fixed}
Grothendieck's qualification about noncanonical isomorphisms concerns genuine monodromy. It can be seen without higher homotopy groups. Let $S=BC_2$, the one-object category with automorphism group $C_2=\{1,\tau\}$, and let $F:S^{\mathrm{op}}\to\Cat$ take its value in the discrete category $D=\{0,1\}$, with $F(\tau)$ exchanging the two objects. All transports are isomorphisms.

There is only one functor $e\to BC_2$, but it has two natural automorphisms. Pullback along this functor is the fibre $D$. The identity natural transformation induces $1_D$, whereas the transformation $\tau$ induces the transposition of $D$. These are different even after localization at $W_\infty$, because they induce different maps on $\pi_0(D)$. Thus ``the two indexing functors are homotopic'' cannot single out an isomorphism between their pullbacks.

The total category $\int F$ is equally explicit. Its objects are $0,1$, and a morphism $d\to d'$ is an element $g\in C_2$ with $d=F(g)(d')$. There is exactly one such element for each ordered pair. The total category is therefore equivalent to $e$, whereas the fibre has two components. Its map to $BC_2$ is the categorical model of the universal two-sheeted covering. The inclusion of the fibre is not a weak equivalence; the base is not contractible. This example also shows why the word ``contractile'' in the next argument is a mathematical hypothesis, not a convenient choice of coordinates.
\end{dcomment}


Il est plus astucieux de ne pas fixer $X$ et de considérer
\[ \uHom(S',S)\fl\uHom(\HOT^{\mathrm{lc}}_W(S),\HOT^{\mathrm{lc}}_W(S')) \tag*{$(*)$} \]
correspondant à
\[
\begin{aligned}
\uHom(S',S) &\quad\times\quad \HOT^{\mathrm{lc}}_W(S) &&\quad\longrightarrow\quad &&\HOT^{\mathrm{lc}}_W(S')\\
(f &\quad,\quad\ \xi) &&\quad\longmapsto\quad && f^*(\xi).
\end{aligned}
\]
Le lemme dira que les flèches de $S'\fl S$ sont transformées par $(*)$ en isomorphismes, d'où
\[ \prod\uHom(S',S)\fl\uHom(\HOT^{\mathrm{lc}}_W(S),\HOT^{\mathrm{lc}}_W(S')), \]
% foot: Version 7 décembre 2002

% ---- from Der-ChXIIa-012_p036-040.tex
\printedpage{34}
% head: A. Grothendieck : Dérivateurs, Ch. XII, Édition des Universités de Montpellier II et Paris VII
% foot: Version 7 décembre 2002

\msp{page 46}
donc si $f,g : S' \rightrightarrows S$ sont deux foncteurs homotopes, alors
\[ f^* \simeq g^* : \HOTlc(S) \rightrightarrows \HOTlc(S'). \]

Il s'ensuit aussitôt ceci

\textbf{Corollaire 2.} \textit{Si $f$ est un homotopisme, $f^* : \HOTlc(S) \longrightarrow \HOTlc(S')$ est une équivalence de catégories.}

En effet, prenant $g : S \longrightarrow S'$ telle que
\[ fg \sim \id_S, \qquad gf \sim \id_{S'}, \]
on trouve que
\[ \begin{aligned}
g^*f^* \simeq (\id_S)^* &= \id_{\HOTlc(S)} \\
f^*g^* \simeq (\id_{S'})^* &= \id_{\HOTlc(S')},
\end{aligned} \]
donc on trouve que (quitte à ajuster ces isomorphismes fonctoriels) $f^*$ et $g^*$ sont des équivalences quasi-inverses l'une de l'autre.

J'ai l'impression (\footnote[54]{O.K.}) que tout ça ne doit pas utiliser plus que les axiomes standard sur $W$ comme localiseur fondamental, tout au plus W(2) au lieu de la forme faible W(2 bis).

Revenons à $X$ lisse sur $S$ contractile. Le corollaire 2 implique que
\[ \begin{aligned}
\Hot_W &\longrightarrow \HOTlc(S) \\
\xi &\longmapsto p^*_S(\xi),
\end{aligned} \]
où $p_S : S \longrightarrow e$, est une équivalence. Donc dans la catégorie $\HOTlc(S)$, $X$ est isomorphe \msp{page 47} à un Cat-fibré de la forme $Z \times S \xrightarrow{\mathrm{pr}_2} S$, avec $Z$ dans Cat. Or pour $X_1 = Z \times S \longrightarrow S$, il est immédiat que $X_{1\,s} \longrightarrow X_1$ est dans $W$, car un produit cartésien de flèches \texttt{[dans]} $W$ (savoir $\id_Z$ et $s \longrightarrow S$) est dans $W$. (En fait, l'inclusion est même un homotopisme!) Notons maintenant que si $X_0$, $X_1$ sont deux objets dans $\underline{\mathrm{Liss}}_0S$, tels qu'il existe une $S$-flèche $X_0 \longrightarrow X_1$ qui soit dans $\Wd$, alors $X_{0\,s} \longrightarrow X_0$ est dans $W$ si et seulement si $X_{1\,s} \longrightarrow X_1$ l'est. En effet, on a
\[ \begin{tikzcd}
X_{0\,s} \arrow[r, "\in W"] \arrow[d, "i_s^{X_0}"'] & X_{1\,s} \arrow[d, "i_s^{X_1}"] & \\
X_0 \arrow[r, "\in W", ""{name=A, below}] & X_1 & \\
 & & \text{\footnotesize\begin{tabular}{@{}l@{}}ici on utilise sans doute\\ W(2), non seulement W(2 bis),\end{tabular}}
\arrow[from=A, to=3-3, bend right]
\end{tikzcd} \]
\printedpage{35}
% head: A. Grothendieck : Dérivateurs, Ch. XII, Édition des Universités de Montpellier II et Paris VII
% foot: Version 7 décembre 2002
qui prouve le résultat.

\dgrblock{dgrb496}{The finite zigzag that makes contraction effective}
\begin{dcomment}{exposition}{p47:qui prouve le résultat.}{The finite zigzag that makes contraction effective}
Grothendieck correctly notices, immediately below, that an isomorphism in a localization does not by itself provide a zigzag consisting entirely of the specified weak equivalences. The contraction supplies such a zigzag directly.

Suppose $F:S^{\mathrm{op}}\to\Cat$ sends arrows into $W$, and choose a finite zigzag
\[
 1_S=\varphi_0\;\Longleftrightarrow\;\varphi_1
 \;\Longleftrightarrow\;\cdots
 \;\Longleftrightarrow\;\varphi_m=c_t
\]
of natural transformations, with the displayed double arrows indicating a choice of one direction at each step. Set $X_j=\int(F\circ\varphi_j^{\mathrm{op}})$ over $S$. Transport along the chosen transformation gives, in the appropriate direction, a functor $X_j\to X_{j+1}$ over $S$ which is in $W$ on every fibre. By the argument of the preceding block it is also in $W$ on total categories.

For every $s\in S$ the resulting square
\[
\begin{tikzcd}
 (X_j)_s\arrow[r]\arrow[d] & (X_{j+1})_s\arrow[d]\\
 X_j\arrow[r] & X_{j+1}
\end{tikzcd}
\]
has horizontal arrows in $W$; if the chosen natural transformation points backwards, reverse the horizontal arrows. Two-out-of-three says that one vertical arrow belongs to $W$ exactly when the other does. At the last stage the family is $F(t)\times S\to S$. The map $e\to S$ selecting $s$ is a categorical homotopy equivalence: a contraction to $t$, evaluated at $s$, supplies a zigzag between $s$ and $t$, which joins the constant functors $c_s$ and $c_t$. Multiplying this homotopy by $F(t)$ shows that
\[
 F(t)\cong F(t)\times\{s\}\longrightarrow F(t)\times S
\]
is a categorical homotopy equivalence, hence is in $W$. Moving back through the finite zigzag proves $X_s\to X\in W$.

The replacement comparison at p.~43 transfers this conclusion to every prelisse pseudofibration; the dual argument treats preproper pseudofibrations. This proves the question of p.~42 and the contractible-base part of Proposition~9 under W(1,2,3). It uses the explicit contraction, componentwise transport, the relative localization axiom, and two-out-of-three. Strong saturation is not needed for this particular conclusion.
\end{dcomment}


Il est vrai que le fait que $X = X_0$ et $X_1 = Z \times S$ soient isomorphes dans $\HOTlc(S)$ n'implique pas, par lui-même, l'existence d'une \msp{page 48} chaîne d'isomorphismes provenant de flèches dans $\Wd$. Mais en fait, si dans l'énoncé du corollaire 2 ci-dessus on compare $X$ dans $\underline{\mathrm{Fib}}_0S$ à $g^*f^*(X)$, grâce à une homotopie $\id_S \sim fg$, on trouve bel et bien une suite de foncteurs $\varphi_i : S \longrightarrow S$, avec $\varphi_0 = \id_S$, $\varphi_N = fg$, et $\varphi_i$ et $\varphi_{i+1}$ reliés par une flèche dans $\underline{\mathrm{Hom}}(S,S)$, donc on a une suite d'objets $X_i = \varphi_i^*(X)$ de $\underline{\mathrm{Fib}}_0S$ (avec $X_0 = X$, $X_N = g^*f^*X$), et $X_i$ et $X_{i+1}$ reliés par une flèche \textit{dans} $\underline{\mathrm{Fib}}_0S$ (pas seulement dans $\HOTlc(S)$), laquelle flèche est dans $\Wd = \Wf$. Donc on gagne, sous réserve de vérifications de routine diverses (\footnote[55]{O.K.}). Donc modulo celles-ci, on aura

\textbf{Proposition 9} (\footnote[56]{C'est ok en supposant que $W$ est \textit{fortement} saturé, i.e. $\hot_W(f)$ isomorphisme $\Longrightarrow f \in W$.},\footnote[57]{Donc $W$-fibration à gauche $\Longleftrightarrow$ $W$-fibration à droite $\Longleftrightarrow$ $W$-fibration.}). \textit{Supposons que $X$ dans $\mathrm{Cat}/S$ soit une $W$-fibration à gauche (cf. prop. 8). Alors c'est même une $W$-fibration (au sens de b), page 40, où c'est appelé `$W$-fibration \textup{hot}-forte'), en d'autres termes, si $S'' \longrightarrow S'$ est sur $S$ avec $S''$, $S'$ contractiles, alors $X'' \longrightarrow X'$ est dans $W$ (\footnote[58]{on trouve seulement que $\hot_W(X'') \longrightarrow \hot_W(X')$ \texttt{[est un]} isomorphisme.}). On a sans doute mieux : si $S'' \longrightarrow S'$ est un homotopisme, alors $X'' \longrightarrow X'$ est dans $W$.}

\dgrblock{dgrb497}{What the localization argument proves}
\begin{dcomment}{exposition}{p48:On a sans doute mieux}{What the localization argument proves}
Grothendieck next asks for invariance under an arbitrary categorical homotopy equivalence of bases. Here the argument changes: he identifies the map of total categories with a map of homotopy colimits and applies an adjunction. Let $v:T\to S$ be a homotopy equivalence and let $\xi$ be a locally constant coefficient represented by the family. The contraction argument for the two composites of a homotopy inverse shows that restriction
\[
 v^*:\HOT_W^{\mathrm{lc}}(S)\longrightarrow\HOT_W^{\mathrm{lc}}(T)
\]
is an equivalence. Once the comparison with strict pullback and the adjunction $p_!\dashv p^*$ have been established, the displayed Hom calculation on p.~50 proves that the actual map of total categories becomes invertible in $\Hot_W$.

This is an assertion about the specified map, rather than merely an abstract isomorphism between its source and target. The last inference to membership in $W$ requires strong saturation. The footnote makes that distinction correctly. Later theory proves that every fundamental localizer in the modern, stronger sense is strongly saturated; see the status discussion at p.~16 and \citep[Proposition~4.2.4]{Cisinski2006}. Here the stronger sense includes the relative axiom W(2), not merely W(2 bis) in the terminology fixed at the beginning of section~III. Theorem~3 explicitly assumes W(1,2,3), so that later theorem removes its additional saturation hypothesis. We are not asserting this consequence for an arbitrary weak fundamental localizer satisfying only W(2 bis).

The finite-zigzag argument in the preceding block already proves the contractible-base tests directly in $W$. Thus it is consistent for Theorem~3 to state those equivalences under W(1,2,3), while presenting its more general assertion $d)$ first in the localized category.
\end{dcomment}


\msp{page 49}
Prouvons (si nous pouvons) ce complément sympathique. Je précise que j'ai eu à faire l'hypothèse de $W$-prélissité (où de $W$-prépropreté) pour être sûr que l'image inverse $X \times_S S'$, par $f' : S' \longrightarrow S$, coïncide bien avec $f^*$ sur les $\HOTlc$ - autrement je n'ai aucune chance. % [sic] où -> ou, the conjunction "or" is meant
`On sait' (tout est à revoir!) que c'est O.K. dans les trois cas $X$ $W$-lisse, $X$ $W$-propre, $X$ une $W$-fibration (mais en quel sens???). La vérification à faire (quand on travaille avec les $\Wd$) est celle-ci :
\[ \begin{aligned}
&X \xrightarrow{i_X} \Phi_S(X) = Y \text{ est dans } \Wd \text{ (sans mérite), et} \\
&\text{le reste par tout changement de base.}
\end{aligned} \tag*{$(*)$} \]

Mais pour des $X, Y \in \mathrm{Ob}\,\mathrm{Cat}/S$ qui sont $W$-prélisse, la $\Wd$-équivalence coïncide avec $\Wf$, laquelle est préservée par changement de base, ainsi que l'hypothèse sur $X$, $Y$. Donc on est tranquille.

On est ramené à prouver que si $X$ est $W$-prélisse sur $S$ et une $W$-pseudofibration, et si $S' \longrightarrow S$ est un homotopisme, alors $X' \longrightarrow X$ est dans $W$. Par l'argument de tantôt utilisant $X \longrightarrow \Phi_S(X)$, on se ramène au cas où $X$ est fibrée scindée.

\printedpage{36}
% head: A. Grothendieck : Dérivateurs, Ch. XII, Édition des Universités de Montpellier II et Paris VII
% foot: Version 7 décembre 2002

\msp{page 50}
\[ \begin{tikzcd}
X' \arrow[rr] \arrow[d] & & X \arrow[d] \\
S' \arrow[rr, "f"] \arrow[dr, "p'"'] & & S \arrow[dl, "p"] \\
 & e &
\end{tikzcd} \]

La flèche $X' = f^*(X) \longrightarrow X$ s'interprète dans $\Hot_W$ comme
\[ p'_!(f^*(\xi)) \longrightarrow p_!(\xi) \ , \qquad \text{où } \xi = \hot_{S,W}(X). \tag*{$(*)$} \]

Je dis que cette flèche est \texttt{[un]} isomorphisme, i.e. que pour tout $\eta \in \Hot_W$, la flèche transposée
\[ \Hom_W(p_!(\xi),\eta) \longrightarrow \Hom_W(p'_!(f^*(\xi)),\eta) \tag*{$(**)$} \]
est \texttt{[une]} bijection. Admettant la formule d'adjonction, la flèche $(**)$ s'interprète comme
\[ \Hom_W(\xi,\underbrace{p^*(\eta)}_{\eta_S}) \longrightarrow \Hom_W(f^*(\xi),\underbrace{p'^*(\eta)}_{f^*(\eta_S)}), \]
i.e.
\[ \Hom_W(\xi,\eta_S) \longrightarrow \Hom_W(f^*(\xi),f^*(\eta_S)), \]
que ce soit \texttt{[un]} isomorphisme résulte donc de $f^*$ pleinement fidèle sur les $\HOTlc$. Or on a vu que $f^*$ est une équivalence (cor. 2, page 46).

\msp{page 51}
Je vais reprendre les propositions 9 et 8, un peu mal foutues.

\textbf{Théorème 3.} \textit{On suppose que $W$ satisfait à \textup{W(1,2,3)}. Soit $f : X \longrightarrow S$. Considérons les conditions suivantes.}

\begin{itemize}[leftmargin=4em]
\item[\textit{a)}] \textit{$\forall\ S'' \longrightarrow S'$ sur $S$, avec $S''$, $S'$} contractiles, \textit{$X'' \longrightarrow X'$ est dans $W$.}

\item[\textit{a}$'$\textit{)}] \textit{$\forall\ S'$ sur $S$, avec $S'$ contractile, et $s' \in S'$, l'inclusion} % AUDIT resolved: the label is a prime, not d prime; see corrections.csv, PDF leaf 46
\[ X_{s'} \hookrightarrow X' = X \times_S S' \]
\textit{est dans $W$.}

\item[\textit{a}$_1$\textit{)}] \textit{Comme a$'$), mais en supposant que $S'$ a un objet final (au lieu de $S'$ contractile).} % AUDIT resolved: the reference is a prime; see corrections.csv, PDF leaf 46

\item[\textit{a}$_2$\textit{)}] \textit{Comme a$'$), mais en supposant que $S'$ a un objet initial.} % AUDIT resolved: the reference is a prime; see corrections.csv, PDF leaf 46

\printedpage{37}
% head: A. Grothendieck : Dérivateurs, Ch. XII, Édition des Universités de Montpellier II et Paris VII
% foot: Version 7 décembre 2002

\item[\textit{b}$_1$\textit{)}] \textit{$f$ est $W$-}prépropre \textit{(i.e. a$'$) est satisfait quand $s'$ est un objet final de $S'$, sans autre condition sur $S'$ dans $\mathrm{Cat}/S$), et de plus pour toute flèche $s_0 \longrightarrow s_1$ dans $S$, la flèche induite $X/s_0 \longrightarrow X/s_1$ est dans $W$ (i.e. a$'$ \textup{\texttt{[plutôt a]}}) est satisfait pour $Y/s_0 \longrightarrow Y/s_1$ \textup{\texttt{[plutôt} $S/s_0 \longrightarrow S/s_1$\texttt{]}}).}

\item[\textit{b}$_2$\textit{)}] \textit{$f$ est $W$-}prélisse \textit{(i.e. a$'$) est satisfait quand $s'$ est un objet initial de $S'$, sans autre condition sur $S'$), et de plus les flèches $s_0\backslash X \longrightarrow s_1\backslash X$ \textup{\texttt{[plutôt} $s_1\backslash X \longrightarrow s_0\backslash X$\texttt{]}} (pour $s_0 \longrightarrow s_1$ dans $S$) sont dans $W$, (i.e. a$'$ \textup{\texttt{[plutôt a]}}) satisfait pour $s_1\backslash Y \longrightarrow s_0\backslash Y$ \textup{\texttt{[plutôt} $s_1\backslash S \longrightarrow s_0\backslash S$\texttt{]}}).}

\item[\textit{c}$_1$\textit{)}] \textit{$f$ est $W$-prépropre et une $W$-pseudofibration, i.e. la condition a$'$) est satisfaite dans les deux cas suivants: 1$^\circ$) $s'$ est un objet final de $S'$, 2$^\circ$) $S' = \mathbf{\Delta}^1$.}

\item[\textit{c}$_2$\textit{)}] \textit{$f$ est $W$-prélisse et une $W$-pseudofibration, i.e. a$'$) dans les deux cas: 1$^\circ$) $s'$ est objet initial de $S'$, 2$^\circ$) $S' = \mathbf{\Delta}^1$.}
\end{itemize}

\textit{Toutes ces 8 conditions sont équivalentes. Elles impliquent la suivante :}

\begin{itemize}[leftmargin=4em]
\item[\textit{d)}] \textit{Soit $S'' \xrightarrow{u} S'$ un homotopisme dans $\mathrm{Cat}/S$, et considérons $u_X : X'' \longrightarrow X'$. Alors $\hot_W(u_X)$ est \textup{\texttt{[un]}} isomorphisme. (Donc si $W$ est fortement saturé, $u_X \in W$. Donc sous cette condition sur $W$, d) équivaut aux conditions précédentes.)}
\end{itemize}

\msp{page 52}
\textbf{Définition 4.} On dit que $f : X \longrightarrow Y$ dans Cat est une \textit{$W$-fibration}, si elle satisfait les conditions équivalentes a) à c$_2$) de la proposition 9.

\dgrblock{dgrb498}{Reading the eight equivalent conditions}
\begin{dcomment}{exposition}{p52:si elle satisfait les conditions équivalentes}{Reading the eight equivalent conditions}
Grothendieck has now isolated a usable fibration notion. Its eight formulations organize two independent requirements: strict fibres must compute comma fibres universally, and the latter must vary by weak equivalences. The reference to ``proposition~9'' in Definition~4 means the immediately preceding Theorem~3; Proposition~9 did not contain eight labelled conditions.

Here is the logical structure, including the point at which the contractible-base argument is used. Condition $a)$ implies $a')$ by taking $S''=e$. Conversely, for a map $u:T\to U$ between contractible categories over $S$, choose $t\in T$. The triangle of inclusions of the same strict fibre into $X_T$ and $X_U$ commutes. Both inclusions are in $W$ by $a')$, so $X_T\to X_U\in W$. Contractible categories are nonempty; there is no missing empty-base case.

The implications $a')\Rightarrow a_1),a_2)$ follow because categories with terminal or initial objects admit categorical contractions. Proposition~8 gives $a_1)\Longleftrightarrow b_1)\Longleftrightarrow c_1)$ and its dual gives $a_2)\Longleftrightarrow b_2)\Longleftrightarrow c_2)$. Starting with either $c_1)$ or $c_2)$, use the corresponding split replacement and the finite transport zigzag to obtain the fibre inclusion over every contractible base. This proves the reverse implication to $a')$ and therefore all eight equivalences.

For an ordinary Grothendieck fibration, the comma comparisons are automatic and the remaining test is simply that every pullback functor $u^*$ lies in $W$; the opfibred statement uses $u_*$. For a family that is both preproper and prelisse, the two interval endpoint inclusions already lie in $W$, so the pseudofibration test is automatic. These observations prove the three corollaries. They do not identify a $W$-fibration with a strict Grothendieck fibration, nor with a perfect functor: the first supplies strict Cartesian lifts, and the second supplies stronger asphericity of lift categories.
\end{dcomment}


\textbf{Corollaire 1.} \textit{Si $f$ est $W$-propre ou $W$-lisse, pour que $f$ soit une $W$-fibration, il faut et il suffit que $f$ soit une $W$-pseudofibration, i.e. que pour $S' = \mathbf{\Delta}^1 \longrightarrow S$, les inclusions des deux fibres}
\[ X'_0 \xrightarrow{i_0} X' \xleftarrow{i_1} X'_1 \]
\textit{de $X' = X \times_S S'$ soient dans $W$.}

\textbf{Corollaire 2.} \textit{Supposons $f$ \textup{Cat}-fibrant. Alors $f$ est une $W$-fibration si et seulement si pout toute flèche $u : s \longrightarrow s'$ dans $S$, $u^* : X_{s'} \longrightarrow X_s$ est dans $W$. Énoncé dual si $f$ est \textup{Cat}-cofibrant.} % [sic] pout -> pour, misprint

\textbf{Corollaire 3.} \textit{Supposons $f$ $W$-prépropre et $W$-prélisse (par exemple $f$ $W$-parfait). Alors $f$ est une $W$-fibration.}

\textbf{Proposition 10.} \textit{Soit $f : X \longrightarrow Y$ une flèche dans $\underline{\mathrm{Fib}}_WS$ (\footnote[59]{$X$ et $Y$ tous deux prépropres ou tous deux prélisses suffit.}) (sous-catégorie pleine de $\mathrm{Cat}/S$ formée des $W$-fibrés sur $S$). Si $\forall\ s \in S$, $f_s : X_s \longrightarrow Y_s$ est dans $W$, alors $f \in W$.}

\dgrblock{dgrb499}{Fibrewise equivalence gives equivalence of total categories}
\begin{dcomment}{exposition}{p52:est dans $W$, alors $f \in W$.}{Fibrewise equivalence gives equivalence of total categories}
Proposition~10 is a comparison theorem in one direction. It says that equivalent fibres give equivalent totals when the structural maps make those fibres compute the relevant comma categories. It does not yet assert that an equivalence of totals can be detected fibre by fibre; that converse is the separate axiom W(4b).

This argument uses the full relative axiom W(2), as assumed in Theorem~3; the weaker standing definition of section~III alone is not the stated justification. Suppose first that $X$ and $Y$ are prelisse over $S$, and that $h:X\to Y$ is in $W$ on every fibre. For each $s$ the commutative square
\[
\begin{tikzcd}
 X_s\arrow[r,"h_s"]\arrow[d] &Y_s\arrow[d]\\
 s\backslash X\arrow[r] &s\backslash Y
\end{tikzcd}
\]
has its upper and vertical arrows in $W$. Hence its lower arrow is in $W$. The dual form of W(2) gives $h\in W$. If the structural maps are preproper, use $X/s$ and $Y/s$ instead. This proves exactly the strengthening in footnote~59, without any pseudofibration hypothesis.

The source's replacement proof packages the same argument differently. The maps $i_X,i_Y$ are homotopy equivalences on total categories and are fibrewise equivalences under prelissity. Their split replacements are smooth over $S$, so the corollary to Proposition~5 applies to the induced map. Referring to that corollary, rather than requiring the map between the replacements itself to be smooth, is the correct reading of the final sentence.
\end{dcomment}


Considérons
\printedpage{38}
% head: A. Grothendieck : Dérivateurs, Ch. XII, Édition des Universités de Montpellier II et Paris VII
% foot: Version 7 décembre 2002
\[ \begin{tikzcd}
X \arrow[rr, "f"] \arrow[dr, "i_X"'] & & Y \arrow[dr, "i_Y"] & \\
 & \Phi_S(X) \arrow[rr, "\Phi_S(f)"'] & & \Phi_S(Y),
\end{tikzcd} \]
où $i_X, i_Y \in W$, donc $\Phi_S(f) \in W \Rightarrow f \in W$. Or le fait que $X$ \texttt{[est]} $W$-fibrée sur $S$ implique que $i_X \in \Wf$ (il suffit même que $X$ soit prépropre \texttt{[plutôt prélisse]} sur $S$), itou pour $Y$, $i_Y$. Donc l'hypothèse $f \in \Wf$ implique $\Phi_S(f) \in \Wf$, donc $\Phi_S(f) \in W$ par \texttt{[la]} proposition 5 (page 29).

% ---- from Der-ChXIIa-012_p041-044.tex
\printedpage{39}
% head: A. Grothendieck : Dérivateurs, Ch. XII, Édition des Universités de Montpellier II et Paris VII

\msp{page 53}

\section*{IV\quad Sur l'axiome W(4) et ses variantes.}

(\footnote[60]{Cf. aussi partie VI (HOT$_W$), p. 186-196}). J'ai passé des jours à essayer de prouver que toutes ces variantes sont équivalentes. Je vais donner une liste de celles, parmi ces variantes, qui avec le recul me paraissent encore avoir une utilité. (Avec des notations différentes de celles des pages 13-14.)

\dgrblock{dgrb500}{The two directions of a fibration axiom}
\begin{dcomment}{exposition}{p53:J'ai passé des jours à essayer}{The two directions of a fibration axiom}
Grothendieck separates the two properties that ordinary homotopy theory makes look inseparable. The first concerns a change of base: a weak equivalence $T\to S$, pulled back along a suitable family $X\to S$, should remain a weak equivalence. The second concerns a map between families over one fixed base: if its total map is a weak equivalence, each fibre map should be one. These are W(4a) and W(4b), respectively. W(4c) is the special case of the second in which the second family is $S\to S$ itself.

The primes restrict the objects on which these properties are tested. In W(4a') and W(4c') the map is perfect for the minimal localizer; in W(4b') the structural families are split and the comparison preserves the splittings; in W(4b'') the structural families are $W$-perfect. These restrictions are not interchangeable merely by renaming a fibration. In particular, the argument of section~I used a stronger base-change condition than the eventual definition on p.~52. The editors' objection to the arrow $a'\Rightarrow b$ records a genuine dependence on that change of definition.

There is also a change of lettering. The weak perfect-fibre axiom called W(4b'') in section~I is now W(4c'). Keeping its content fixed is safer than carrying its old label into a later argument. The subsequent proof establishes one implication through locally constant coefficients; it does not prove every arrow that the surrounding discussion anticipates.
\end{dcomment}


\begin{itemize}[leftmargin=5em,labelsep=0.5em]
\item[\fbox{W(4a)}] Toute $W$-fibration (cf. th. 3, p. 51, déf. 3, p. 52) est ultraforte, ce qui revient à ceci: pour
\[
\begin{tikzcd}
X & X' \arrow[l,"q"'] \arrow[d,"f'"] \\
Y \arrow[from=u,"f"'] & Y' \arrow[l,"p"']
\end{tikzcd}
\]
cartésien dans Cat, si $f \in \underline{\Fib}_W$, $p \in W$, alors $q \in W$.

\item[W(4a$'$)] Comme a), en supposant $f \in \Parf_\omega$, i.e. $f$ $W_\omega$-parfait (donc une $W_\omega$-fibration), où $W_\omega \subseteq W$ est la plus petits partie $\Sigma$ de Fl(Cat) satisfaisant W(1, 2 bis, 3). % [sic] petits -> petite, accord avec « partie »

\item[\fbox{W(4b)}] Soit un diagramme commutatif
$\begin{tikzcd}[ampersand replacement=\&, column sep=small, row sep=small, baseline=(X.base)]
|[alias=X]| X \arrow[rr,"f"] \arrow[dr,"p"'] \& \& Y \arrow[dl,"q"] \\
\& S \&
\end{tikzcd}$
dans Cat, avec $p, q \in \Fib_W$, $f \in W$.
Alors $\forall s \in S$, $f_s : X_s \longrightarrow Y_s$ \texttt{[est dans} $W$\texttt{]}.

\item[W(4b$'$)] Comme b), en supposant \texttt{[de plus]} $X$ et $Y$ fibrés scindés sur $S$, $f$ compatible avec \texttt{[les]} scindages (donc foncteur cartésien).

\item[W(4b$''$)] Comme b), en supposant $p$, $q$ $W$\textit{-parfaits.}

\item[\fbox{W(4c)}] Soit $f : X \longrightarrow Y$, avec $f \in \Fib_W \cap W$, alors les fibres de $f$ sont $W$-asphériques.

\item[\dashedbox{W(4c$'$)}] Comme c), en supposant $f \in \Parf_\omega \cap W$. (C'est l'ancien W(4b$''$), dépouillé du $W_\infty$-parasite.)
\end{itemize}

\msp{page 54}

J'ai cerclé les notions \texttt{[encadrées ici]} qui ont été vraiment utiles. Pratiquement, on peut parier que c'est la notion la plus forte, W(4a), qui est la seule importante. Voici les implications tautologiques (à l'exception de W(4a$'$) $\Longrightarrow$ W(4b), qui est démontré dans la démonstration de la proposition 4 (pages 6-10) (\footnote[61]{Cette démonstration utilise une notion plus forte de $W$-fibration que celle du théorème 3, p. 51.})).

\printedpage{40}
% head: A. Grothendieck : Dérivateurs, Ch. XII, Édition des Universités de Montpellier II et Paris VII

\[
\begin{tikzcd}[column sep=small]
& \mathrm{a} \arrow[d,Rightarrow] & \\
& \mathrm{a}' \arrow[dd,Rightarrow,"{\text{\rule[0.6ex]{2cm}{0.4pt}\ \begin{tabular}[t]{@{}l@{}}implication douteuse,\\ il faut affaiblir\\ un peu b), en\\ renforçant la\\ notion de\\ $W$-fibration\\ utilisée\\ pour énoncer b)\end{tabular}}}"] & \\
& & \\
& \mathrm{b} \arrow[dl,Rightarrow] \arrow[dr,Rightarrow] \arrow[dd,Rightarrow] & \\
\mathrm{b}' & & \mathrm{b}'' \\
& \mathrm{c} \arrow[d,Rightarrow] & \\
& \mathrm{c}' &
\end{tikzcd}
\]

J'ai surtout sué pour montrer que c$'$) $\Longrightarrow$ c), ou c) $\Longrightarrow$ b), sans y arriver. L'effort m'a permis de me débarrasser de diverses idées fausses, et de les remplacer par des justes, cf. plus bas. Il y a quand même un résultat positif qui surnage :

\textbf{Proposition 1.} \textit{La condition b) (et sans doute même b$'$) et b$''$)} (\footnote[62]{et en supposant dans le cas de b$''$, que $W$ satisfait la condition des limites inductives W(6 bis), cf. prop. 5 du `corollaire' de la, page 14.})\textit{, mais je ne vais pas le prouver ici \dots) implique sans doute a) pour $\bar{W}$, le saturé fort de $W$. En tous cas, si $W$ est fortement saturé, les conditions a), a$'$), b) (et sûrement aussi b$'$) et b$''$)) sont équivalentes.}

\dgrblock{dgrb501}{The later resolution of the weakest fibre axiom}
\begin{dcomment}{status}{p54:En tous cas, si $W$ est fortement saturé}{The later resolution of the weakest fibre axiom}
Grothendieck's reservations here belong to the search for a proof from a deliberately small list of axioms. The minimal weak-localizer theorem gives $W_\omega=W_\infty$ even with W(2 bis), by \citet[Theorem~6.1.18]{Cisinski2006}. The proof in section~I, completed by the nondegeneracy argument at p.~20, then shows that W(4c') forces
\[
 W=W_\infty\quad\hbox{or}\quad W=\operatorname{Fl}(\Cat).
\]
The second possibility must be retained: all the fibration axioms hold vacuously when every functor is a weak equivalence. This argument uses the perfect replacement and iterated homotopy-fibre construction, together with the minimality theorem; it does not require filtered-colimit stability or strong saturation of an arbitrary weak localizer. Those properties follow after the displayed classification. For the nontrivial alternative, ordinary homotopy theory gives the full pullback and fibrewise comparison statements. A related published characterization under the full relative-localization hypothesis is \citet[Theorem~6.4.22 and Remark~6.4.23]{Cisinski2006}, formulated using locally constant functors.

The split comparison test $b')$ also detects the weakest perfect-fibre test. If $E\to B$ is $W_\infty$-perfect and belongs to $W$, compare the split replacements $\Phi_B(E)\to\Phi_B(B)$. Its total map is in $W$, since both replacement units are. Their transports are weak equivalences, and the comparison is strictly Cartesian. Thus $b')$ makes each $b\backslash E\to b\backslash B$ a $W$-equivalence. The perfect fibre comparison $E_b\to b\backslash E$ and the initial object of $b\backslash B$ then show that $E_b$ is $W$-aspherical. The implication from the perfect pullback test $a')$ uses the mapping-path comparison already audited in section~I, rather than the disputed reuse of an older fibration definition.

Accordingly the weak perfect test, the split test, and the full tests have the same eventual alternatives. This retrospective conclusion does not turn the unfinished implications in the manuscript into proofs written there. In the claim register their original qualifications remain conjectural or sketched, with the later resolution recorded separately.
\end{dcomment}


\textsc{Démonstration} de b) $\Longrightarrow$ a), quand $W$ est fortement saturé. En fait, moyennant b) (et sans même supposer $W$ fortement saturé), nous allons prouver que si on a
\[
\begin{tikzcd}
X & X' \arrow[l,"q"'] \arrow[d,"f'"] \\
Y \arrow[from=u,"f"'] & Y' \arrow[l,"p"']
\end{tikzcd}
\]
cartésien dans Cat, avec
\msp{page 55}
$f$ une $W$-fibration, et $p \in W$, alors $\underline{q \in \bar{W}}$, i.e. $\hot_W(q)$ \texttt{[est un]} isomorphisme. Pour ceci, nous interprétons la flèche $q$ dans $\Hot_W$ de la façon suivante. Soit
\[
\xi = \hot_{Y,W}(X/Y), \qquad \xi' = p^*(\xi),
\]
\printedpage{41}
% head: A. Grothendieck : Dérivateurs, Ch. XII, Édition des Universités de Montpellier II et Paris VII
qui dit que $q$ s'identifie à la flèche canonique
\[
H_\bullet^{\HOT_W}(Y',\xi') \longrightarrow H_\bullet^{\HOT_W}(Y,\xi) \qquad \text{dans } \Hot_W = \HOT_W(e), \tag*{$(*)$}
\]
où on note de plus que $\xi \in \HOT_W(Y)$ est un coefficient \textit{localement constant} sur $Y$.

Établissons d'abord cette interprétation de $q$. Notons que comme $X \in \underline{\Fib}_W Y$, l'image inverse $\xi' = p^*(\xi)$ n'est autre que $\hot_{Y',W}(X' = X \times_Y Y')$. D'autre part, si $g : Y \longrightarrow e$, $g' : Y' \longrightarrow e$ sont les flèches structurales, les deux membres de $(*)$ s'interprètent resp. comme $g'_!(\xi')$ et $g_!(\xi)$ (les $g'_!$ et $g_!$ au sens du prédérivateur $\HOT_W$). Or le calcul de $g'_!$, $g_!$ nous donne que les deux membres s'identifient resp. à
\[
\hot_W(X'), \qquad \hot_W(X),
\]
et nous admettons que la flèche
\[
\hot_W(X') \longrightarrow \hot_W(X)
\]
n'est autre que $\hot_W(q)$. D'ailleurs, comme $X$ est $W$-fibrée sur $S$, par définition pratiquement, $\xi$ est dans $\HOT^{\mathrm{lc}}_W(S)$. Donc il nous suffira de prouver le

\msp{page 56}

\textbf{Corollaire.} \textit{Supposons que $W$ satisfait à W(4b), et soit $p : S' \longrightarrow S$ dans $W$, $\xi \in \HOT^{\mathrm{lc}}_W(S)$, $\xi' = p^*(\xi)$. Alors}
\[
H_\bullet(S',\xi') \xrightarrow{\ \sim\ } H_\bullet(S,\xi)
\]
\textit{est un isomorphisme dans} $\Hot_W$.

\dgrblock{dgrb502}{Why fully faithful restriction controls integration}
\begin{dcomment}{exposition}{p56:est un isomorphisme dans}{Why fully faithful restriction controls integration}
Grothendieck's proof is driven by the fact that a homotopy colimit is characterized by its maps into constant coefficients. We can state the argument in a form which makes every required comparison visible. Let $p:T\to S$, and suppose that restriction is fully faithful on locally constant coefficients. Let $q:S\to e$ and $r:T\to e$ be the structural maps, and suppose the adjunctions $q_!\dashv q^*$ and $r_!\dashv r^*$ are available. For a locally constant $\xi$, there is a canonical comparison
\[
 c:r_!p^*\xi\longrightarrow q_!\xi.
\]
For every coefficient $A$ at the point, adjunction identifies precomposition by $c$ with
\[
 \operatorname{Hom}(\xi,q^*A)
 \longrightarrow
 \operatorname{Hom}(p^*\xi,p^*q^*A).
\]
Both $\xi$ and $q^*A$ are locally constant, so full faithfulness makes this a bijection. Yoneda implies that $c$ is an isomorphism. Essential surjectivity of restriction is stronger than this argument requires.

The application to a literal pullback $X_T\to X$ needs two further identifications: strict pullback must represent $p^*\xi$, and the chosen integration functors must send these representatives to their total categories, with $c$ represented by the actual projection. Prelissity supplies the first identification; the construction of $q_!$ supplies the second. The source acknowledges these comparisons instead of proving them anew. The same formal Hom argument appears in \citet[Proposition~1.10]{Cisinski2009LocallyConstant}, and its Theorem~1.3 supplies ordinary weak-homotopy invariance for model-category-valued locally constant diagrams. These later results clarify the argument without making W(4b) an unmentioned consequence of W(1,2,3).
\end{dcomment}


\textsc{Démonstration}. Il suffit de prouver que pour tout $\eta \in \mathrm{Ob}\,\Hot_W$ on a une bijection pour les $\Hom(-,\eta)$. Par la formule d'adjonction sur $S$, $S'$, cette application entre les Hom dans $\Hot_W$ s'identifie \texttt{[à]}
\[
\Hom_{\HOT_W(S)}(\xi,\eta_S) \longrightarrow \Hom_{\HOT_W(S')}(\xi',\eta_{S'}),
\]
où $\eta_S = p^*(\eta)$, $\eta_{S'} = p'^*(\eta)$ \texttt{[plutôt} $\eta_S = g^*(\eta)$\texttt{,} $\eta_{S'} = g'^*(\eta)$\texttt{, où} $g : S \longrightarrow e$\texttt{,} $g' : S' \longrightarrow e$ \texttt{sont les flèches structurales]}, donc $\eta_{S'} = p^*(\eta_S)$, et la flèche $(*)$ est la flèche sur les Hom induite par le foncteur
\[
p^*_{\mathrm{lc}} : \HOT^{\mathrm{lc}}_W(S) \longrightarrow \HOT^{\mathrm{lc}}_W(S').
\]
Or moyennant l'axiome W(4b) (que nous n'avons pas utilisé jusqu'à présent), ledit foncteur est une \textit{équivalence de catégories} si $p \in W$ (cf. VI, p. 204).

\hfill q.e.d.

De peur de les oublier, je vais maintenant faire une liste des petites questions que j'ai élucidées chemin faisant en gribouillant.

\dgrblock{dgrb503}{Why preservation of Cartesian arrows matters}
\begin{dcomment}{examples}{p56:De peur de les oublier}{Why preservation of Cartesian arrows matters}
Grothendieck's first observation separates fibrewise smoothness from smoothness of the total functor. A small example makes the lost information explicit. Let $C$ be the one-object category with endomorphism monoid $(\mathbf N,+)$, and put $X=Y=C\times[1]$, both projected to $[1]$. Define an endofunctor $h:X\to Y$ over $[1]$ to be the identity on both fibres and to send a morphism crossing from $0$ to $1$, labelled $n\in\mathbf N$, to the morphism labelled $n+1$.

This is a functor. A composite crosses from $0$ to $1$ at most once, so adding $1$ commutes with precomposition and postcomposition by fibre morphisms. Each fibre functor is the identity and hence is smooth and proper. Both structural projections are products, hence split fibrations. But $h$ does not preserve Cartesian arrows: the Cartesian crossing labelled $0$ is sent to the non-Cartesian crossing labelled $1$.

To test smoothness of $h$, take the target object over $1$ and the arrow in $Y$ from the object over $0$ labelled $0$. A lift with this prescribed image would be a crossing morphism labelled $n$ in $X$ with $n+1=0$. There is none. Its lifting category is empty, so $h$ is not $W_\infty$-smooth. This is a fully explicit instance of the failure claimed in item~1. If the functor preserves the splitting, an arbitrary arrow over $u:s\to t$ factors through the chosen Cartesian lift, and its lift category is the corresponding lift category of the fibre functor at $s$. That identification proves the positive assertion in item~1.
\end{dcomment}


\begin{itemize}[leftmargin=5em,labelsep=0.5em]
\item[1°)] Si $X$, $Y$ sont $W$-lisses sur $S$, et $f : X \longrightarrow Y$ un $S$-morphisme tel que $f_s : X_s \longrightarrow Y_s$ soit $W$-lisse pour tout $s \in \mathrm{Ob}\,S$, il ne s'ensuit \textit{pas} que $f$ soit $W$-lisse. Mais c'est pourtant vrai quand $X$ et $Y$ sont fibrées scindées sur $S$ (ce
\printedpage{42}
% head: A. Grothendieck : Dérivateurs, Ch. XII, Édition des Universités de Montpellier II et Paris VII
\msp{page 57}
qui implique déjà leur $W$-lissité sur $S$), et $f$ compatible avec les scindages et a fortiori cartésien. Énoncé dual pour les catégories cofibrées scindées et les foncteurs strictement cartésiens entre elles qui sont $W$-propres fibre par fibre : ils sont $W$-propres.

\item[2°)] J'étais intéressé, dans un cas de catégories fibrées scindées $X$, $Y$ sur $S$ et d'un foncteur strictement cartésien $f : X \longrightarrow Y$, supposé $W$-parfait fibre par fibre, de conclure que $f$ soit aussi $W$-parfait. (C'était un cas où de plus on sait que $X$, $Y$ sont aussi des $W$-fibrés sur $S$, i.e. les fonteurs changement de base $X_{s'} \longrightarrow X_s$, $Y_{s'} \longrightarrow Y_s$ sont dans $W$. Mais ça n'a pas beaucoup servi \dots) La question qui restait, en 1°), était celle de la $W$\textit{-propreté} de $f$. Mais c'est sans espoir (même pour $Y = S$), en vertu de la brillante % [sic] fonteurs -> foncteurs, lettre omise

\textbf{Proposition 2.} \textit{Supposons que} $S = 0$
$\begin{tikzcd}[ampersand replacement=\&, row sep=tiny, column sep=small, baseline=(Z.base)]
\& a \\
|[alias=Z]| \phantom{0} \arrow[ur] \arrow[dr] \& \\
\& b
\end{tikzcd}$
\textit{, et soit $Z$ une catégorie sur $S$. Alors $\Phi(Z)$ est propre sur $S$ si et seulement si $Z = \emptyset$.}

\dgrblock{dgrb504}{The fork obstructs outgoing lifts}
\begin{dcomment}{exposition}{p57:Alors $\Phi(Z)$ est propre sur $S$}{The fork obstructs outgoing lifts}
The fork is being used for a second purpose here: a free split fibration supplies incoming Cartesian lifts, while properness requires the appropriate outgoing lift categories to be aspherical. The commentary on Chapter~VI, manuscript pp.~190--192, gives the full fork calculation and the support obstruction; we refer to it rather than repeat the calculation.

For the present application, one may see the obstruction at an object $(0,z,a)$ of $\Phi_S(Z)$ whose label in $S$ lies on one branch. The arrow from $0$ to the other branch has no compatible outgoing lift from that object: it would require a morphism in the fork from the first terminal branch to the other. If $Z$ has an object over either branch, this construction applies. If every object of the nonempty category $Z$ lies over $0$, then an outgoing lift towards either terminal branch is again impossible, because no object of $Z$ can label its target. These two cases exhaust the possibilities. In each case the relevant lift category is empty.

Thus the assertion is correct for ordinary properness, and for $W$-properness whenever $W$ is nontrivial: the empty category is then not $W$-aspherical. If $W=\operatorname{Fl}(\Cat)$, every lift category, including the empty one, is $W$-aspherical and every functor is $W$-proper. That degenerate case must be excluded from the relative reading of ``if and only if $Z=\varnothing$''. Fibrewise perfectness of the structural map to $S$ carries no information about this outgoing obstruction, since every map from a category to the point is perfect.
\end{dcomment}


Prenons alors dans 2°) ci-dessus $X = \Phi(Z)$, $Y = S$, $f : \Phi(Y) \longrightarrow S$ \texttt{[plutôt} $f : \Phi(Z) \longrightarrow S$\texttt{]} le morphisme structural (donc les $f_s$ sont $W$-parfaits), alors $f$ n'est propre que si $X = \emptyset$. Pourtant, si on prend $Z$ $W$-fibrée sur $S$, $\Phi(Z)$ l'est aussi, et on a un contre-exemple à 2°).
\end{itemize}

Mais on s'est quand même convaincu du résultat suivant :

\textbf{Proposition 3.} \textit{Soient $X$, $Y$ catégories fibrées scindées sur $S$, $f : X \longrightarrow Y$ $S$-foncteur strictement cartésien, tel que}
\msp{page 58}
\textit{pour tout $s \in S$, $f_s : X_s \longrightarrow Y_s$ soit} $W$-fibrant. \textit{Alors,} si $W$ satisfait W(4a), $f$ \textit{est $W$-fibrant.}

\dgrblock{dgrb505}{A missing hypothesis in the last proposition}
\begin{dcomment}{development}{p58:si $W$ satisfait W(4a), $f$}{A missing hypothesis in the last proposition}
Grothendieck is trying to lift a fibrewise fibration statement through the Grothendieck construction. The preceding paragraph mentions an additional assumption that both structural families are $W$-fibrations over $S$. That assumption is absent from Proposition~3 as printed. The literal statement is false, even for ordinary weak equivalences.

\textbf{A strictly Cartesian counterexample.}\ours\quad
Take $S=[1]$. Let $F,G:S^{\mathrm{op}}\to\Cat$ be diagrams of discrete categories
\[
 F(0)=F(1)=e,\qquad
 G(0)=e,\quad G(1)=\{a,b\},
\]
with the unique transition maps. Let $\eta_0$ be the identity and let $\eta_1:e\to\{a,b\}$ select $a$. These components form a strict natural transformation. Its Grothendieck construction is the inclusion
\[
 h:\{o\to a\}\longrightarrow\{o\to a,\ o\to b\}.
\]
Both source and target are split fibrations over $[1]$, and $h$ preserves the chosen Cartesian lifts: the only nonidentity chosen lift in the source is sent to $o\to a$. Each component of $\eta$ is perfect. Indeed every functor between discrete categories is both a Grothendieck fibration and an opfibration; the only arrows to lift in its target are identities. It is consequently a $W_\infty$-fibration.

Pull back $h$ along the arrow $o\to b$ in the target. The resulting category has one object, lying over $0\in[1]$. Its fibres are $e$ at $0$ and $\varnothing$ at $1$. The second endpoint inclusion is $\varnothing\to e$, which is not a weak equivalence. Thus $h$ is not even a $W_\infty$-pseudofibration. This proves the counterexample. Notice that the target structural family is not a $W_\infty$-fibration: its transport is $\{a,b\}\to e$. This explains exactly how the extra assumption mentioned in the preceding discussion would exclude the example.

Here is a sufficient repair for the fibrewise-perfect situation that motivated the proposition. It states explicitly the two comparison properties used in the proof rather than hiding them in an appeal to ordinary intuition.

\textbf{A sufficient fibrewise-smooth repair.}\ours\quad
Assume W(4a) and W(4b). Let $F,G:S^{\mathrm{op}}\to\Cat$ send every arrow into $W$, and let $\eta:F\Rightarrow G$ be a strict natural transformation. Suppose every $\eta_s$ is both $W$-smooth and a $W$-fibration; fibrewise perfectness is sufficient. Then $h=\int\eta$ is $W$-smooth and a $W$-fibration.

First, an arrow in $\int G$ from $(s,y)$ to the image of $(t,x)$ is given by $u:s\to t$ and $\beta:y\to\eta_sF(u)x$. Its lift category for $h$ is exactly the lift category of $\beta$ for $\eta_s$. This proves smoothness, because the latter category is $W$-aspherical.

Next fix $u:s\to t$. Form
\[
 P=G(t)\times_{G(s)}F(s).
\]
The projection $P\to F(s)$ belongs to $W$ by W(4a), since it is the pullback of $G(u)\in W$ along the $W$-fibration $\eta_s$. The naturality map $F(t)\to P$ composes with this projection to $F(u)\in W$, so it also belongs to $W$. Both $F(t)\to G(t)$ and $P\to G(t)$ are $W$-fibrations. W(4b) therefore shows that for every $z\in G(t)$ the induced relative-fibre transport
\[
 F(t)_z\longrightarrow F(s)_{G(u)z}
\]
belongs to $W$.

To finish, factor an arbitrary arrow of $\int G$ into an arrow within a fibre followed by a chosen Cartesian arrow. For an arrow within $G(s)$, transport on the under-comma categories of $h$ is in $W$ because $\eta_s$ is a $W$-fibration. The comparison with its total under-comma category is a homotopy equivalence: an object above $t$ with base arrow $u:s\to t$ is sent to its Cartesian pullback $F(u)x$ above $s$, giving a right adjoint to the inclusion of the fibre under-comma category. For a chosen Cartesian arrow, the displayed relative-fibre transport and smoothness identify the corresponding under-comma map as an element of $W$. Composition now gives the assertion for every arrow of $\int G$. The dual form of Proposition~8, together with smoothness, proves that $h$ is a $W$-fibration.

This repair covers the preceding discussion when its structural families and fibrewise-perfect hypotheses are retained. It does not claim that adding only the omitted structural-fibration hypothesis proves the entire printed formulation. We checked the later treatments of smooth functors and locally constant families cited in this commentary and did not locate this exact small counterexample or this explicit sufficient repair there. These are deductions given here, without a claim of priority.
\end{dcomment}


(J'espérais avoir que $f$ est $W$-fibrant `à l'\oe{}il', pour prouver que W(4c) implique W(4b). Mais j'étais obligé de supposer W(4a) \dots)


\dgrblock{dgrb506}{The coherence condition in a later gluing theorem}
\begin{dcomment}{later}{p58:J'espérais avoir que $f$}{The coherence condition in a later gluing theorem}
The last failed shortcut has a useful later analogue. \citet[Theorem~7.1]{Rezk1998Sharp} proves that the homotopy colimit of a diagram of sharp maps is sharp when every naturality square is homotopy Cartesian. This square condition is the coherence missing from a bare list of fibrewise fibration assertions. In the counterexample just given, the square of discrete fibres
\[
\begin{tikzcd}
 e\arrow[r]\arrow[d] & e\arrow[d]\\
 \{a,b\}\arrow[r] & e
\end{tikzcd}
\]
is not homotopy Cartesian: the homotopy pullback of the other three corners has two points.

This is a comparison of mathematical mechanisms, not an identification of all the models. Thomason's comparison between the nerve of a Grothendieck construction and a homotopy colimit is a weak equivalence; such a comparison alone does not transport the property that every \emph{strict} pullback is homotopy Cartesian. To deduce a statement about the literal functor of Proposition~3, the comparison must also control the pullbacks that occur in its fibration test. That is why the source's dependence on W(4a), and the omitted control of structural transport, cannot be removed by a slogan about homotopy colimits.
\end{dcomment}
\begin{otherlanguage}{english}
\dgroverview{XI-XIIa}{Chapter XII \textperiodcentered\ \textup{XI-XIIa} \textperiodcentered\ Overview}
\begin{otherlanguage}{english}
\noindent\textbf{Fibrations, localizers, and the coherence of fibres}

\medskip
The pages assembled here turn a useful geometric intuition into a problem of foundations. A family of categories should behave like a fibration when its strict fibres calculate homotopy fibres and vary coherently over the base. Grothendieck first characterizes that behaviour for ordinary weak homotopy equivalences. He then asks which axioms on an abstract class \(W\) of functors permit the same constructions, and discovers that several apparently natural requirements already force the whole ordinary homotopy theory. The resulting revisions are part of the argument: the word “fibration” changes its tests as the desired scope of the theory becomes clearer.

These pages comprise Chapter XI, manuscript pp.~127--136, and Chapter XII, manuscript pp.~1--58, through the end of section~IV. These are separate manuscript sequences. The edition explicitly says that Chapter XI has no pp.~1--126; their absence is not a gap in the text. We write \(BC=|NC|\) for the classifying space of a category and \(W_\infty\) for the functors whose nerves are weak homotopy equivalences.

\medskip
\noindent\textbf{The geometric theorem and its coefficient proof.}
Chapter XI begins with a property of the actual pullback functor:
\[
 f^*:\Cat/Y\longrightarrow\Cat/X
 \quad\hbox{preserves }W_\infty.
\]
This is a universal condition on changes of base, not a condition merely on the homotopy types of the individual fibres. Theorem~1 identifies it with universal comma-category tests and with base-change and local-constancy statements for direct images of locally constant coefficients. In the terminology later used for model categories, these are sharp maps for the Thomason weak equivalences \citep[\S2]{Rezk1998Sharp}. Quillen's Theorem B explains the geometry of the comma tests: when their transition functors are weak equivalences, the strict fibre, compared with the appropriate comma category, represents the homotopy fibre \citep[\S1]{Quillen1973}.

Most of the coefficient argument is complete as written. The degree-zero comparison and the first-quadrant Leray comparison require no new finiteness hypothesis. The nonabelian passage needs more care. A component of a torsor groupoid has the automorphism group of a representative torsor; this group can be an inner form of the coefficient group, and its direct image is not generally the fixed group printed in the proposed obstruction sequence. We therefore replace that sequence by an argument with the groupoids of torsors on \(X/y\). The universal degree-zero hypothesis gives full faithfulness of their transition functors; the degree-one hypothesis gives essential surjectivity. The same two hypotheses identify their values with the torsor groupoids on strict fibres. Descent reconstructs global torsors from this locally constant diagram, and homotopy invariance of such diagrams proves the required comparison \citep[Theorem~1.3]{Cisinski2009LocallyConstant}.

Theorem~2 then detects an equivalence between two such families on one fibre in each connected component of the base. The source's local reduction for torsors becomes precise through the same descent argument. Its converse uses a contractible replacement of the base. We give an independent construction by factoring a point in Thomason's model category \citep[Theorem~4.9]{Thomason1980}, and also check the specific path-category dependency that Chapter XII recalls from Chapter VII. This separates the existence of a suitable replacement from the stronger task of constructing it explicitly with both properness and smoothness.

\medskip
\noindent\textbf{What the axioms force.}
Chapter XII starts with identities and cancellation, a theorem-A principle, and the asphericity of categories with a terminal object. Two forms of theorem A occur. W(2) is relative: equivalences on comma categories over every object of a third category imply equivalence of the total functor. W(2 bis) is absolute: aspherical comma categories imply that a functor is an equivalence. The manuscript deliberately uses the latter in several arguments. Modern “fundamental localizer” normally includes the former. We keep the distinction throughout the audit.

The minimum conjectured in these pages is now known. Cisinski identifies \(W_\infty\) as the smallest fundamental localizer \citep[Theorem~2.2.11]{Cisinski2004Minimal}, and as the smallest weak fundamental localizer with the absolute theorem-A axiom \citep[Theorem~6.1.18]{Cisinski2006}. Strong saturation and stability under filtered colimits are also established for full fundamental localizers; those results are not assertions about every class satisfying only the weaker axiom. An inference from “invertible in the localization” to “belongs to \(W\)” is justified at the point where its applicable hypothesis is available.

The fibration axioms add a rigidity that the basic localizer axioms do not possess. Their weakest perfect-fibre form, combined with nondegeneracy, makes homotopy fibres of \(W\)-maps \(W\)-aspherical. Iterating the loop construction then detects all homotopy groups. We repair the source's nondegeneracy argument without its unjustified restriction to a component: a missed component produces an empty homotopy fibre, while two identified components produce a disconnected aspherical fibre, whose diagonal in its square again has an empty homotopy fibre. Asphericity of the empty category collapses the theory. Thus the alternatives are \(W_\infty\) and the class of every functor. The second alternative cannot be discarded from the opening assertion about derivators; the terminal derivator realizes it.

The truncation examples explain why a more flexible fibration notion is needed. The class \(W_n\) retains homotopy groups only through degree \(n\), and is a genuine localizer. But the based-path family of an \((n+1)\)-sphere has an invisible total space and an invisible base while its fibre retains a nonzero \(n\)-th homotopy group, or a nontrivial component set when \(n=0\). Hence the strongest fibre axioms fail. In degree zero, the exact sequence must be read in pointed sets and with the appropriate group action: a cokernel of groups does not describe the component set in general.

\medskip
\noindent\textbf{Lift categories supply the missing structure.}
Section~III develops proper and smooth functors through categories of lifts with prescribed source or target. Those categories turn universal comma comparisons into workable tests, prove stability under base change and composition, and explain why Grothendieck (op)fibrations furnish examples. They also support the passage from fibrewise equivalence to equivalence of total categories. The later formulations by Maltsiniotis describe the corresponding exact squares, while the derivator account explains the associated base-change transformations \citep{Maltsiniotis2005Asphericite,Maltsiniotis2012Exact,Groth2013}.

The proposed weakening to endpoint equivalences over every arrow loses a necessary compatibility. We give a poset with eight objects over \([2]\). Its fibres have two contractible components, every interval pullback has both endpoint inclusions in \(W_\infty\), and transport along \(0\to1\) and \(1\to2\) preserves the component labels. Transport along \(0\to2\) exchanges them. Its total nerve is an eight-cycle. Consequently the edge transports do not compose, and the pseudofibration question raised on manuscript p.~40 has a negative answer. The complete proof is independent of the accompanying finite computation.

The successful replacement uses universal preproperness or prelissity in addition to the interval test. Comma categories then provide actual functorial transport. A split replacement converts this into a diagram of categories whose arrows lie in \(W\). For a categorically contractible base, a finite zigzag of natural transformations gives an explicit finite zigzag of fibrewise and total \(W\)-equivalences. This proves the contractible-base portion of Theorem~3 directly, without extracting a zigzag from an abstract isomorphism in a localization. Here “contractible” has the manuscript's categorical meaning; a weakly contractible nerve need not admit such a finite contraction.

\medskip
\noindent\textbf{The last proposition and the limits of the repair.}
Section~IV returns to the relations among the fibration axioms, then proposes a fibrewise gluing statement. Its final Proposition~3 is false literally as printed. Over \([1]\), take the strictly Cartesian inclusion of the split fibration \(\{o\to a\}\) into \(\{o\to a,o\to b\}\). Every fibre map is a functor between discrete categories and is perfect. Pulling the total functor back along \(o\to b\) gives endpoint fibres \(e\) and \(\varnothing\); it is not even a pseudofibration. The target structural family is not a \(W_\infty\)-fibration, precisely the extra condition mentioned in the preceding discussion and omitted from the proposition.

We prove a sufficient repair for that preceding, fibrewise-perfect situation: the structural transports lie in \(W\), the fibre maps are smooth \(W\)-fibrations, and W(4a) and W(4b) control the relevant pullback and fibre comparison. We do not claim that adding only the omitted structural condition repairs the full printed statement. Rezk's later gluing theorem has an analogous additional requirement that naturality squares be homotopy Cartesian \citep[Theorem~7.1]{Rezk1998Sharp}; passing between models does not by itself remove the need to control strict pullbacks.

The marked developments give these deductions with complete proofs and bounded accounts of the literature examined, without a claim of priority. The registers distinguish what the source proves, sketches, assumes, or conjectures from the status supplied by the commentary. They preserve edition slips in the French text and apply only corrections of the transcription. Questions not decided under their exact weaker hypotheses are identified as limitations of this audit, rather than declared open in the literature.

\end{otherlanguage}
\end{otherlanguage}
\dgrchunkend
